\documentclass[11pt]{amsart}
\usepackage{amsmath,amssymb,amsthm,mathtools,booktabs,hyperref,url,array}
\usepackage[margin=1.1in]{geometry}
\newtheorem{theorem}{Theorem}[section]
\newtheorem{proposition}[theorem]{Proposition}
\newtheorem{conjecture}[theorem]{Conjecture}
\newtheorem{lemma}[theorem]{Lemma}
\newtheorem{corollary}[theorem]{Corollary}
\theoremstyle{definition}
\newtheorem{definition}[theorem]{Definition}
\newtheorem{remark}[theorem]{Remark}
\newtheorem{example}[theorem]{Example}
\newcommand{\Q}{\mathbb{Q}}
\newcommand{\Z}{\mathbb{Z}}
\newcommand{\F}{\mathbb{F}}
\newcommand{\R}{\mathbb{R}}
\newcommand{\C}{\mathbb{C}}
\newcommand{\T}{\mathbb{T}}
\newcommand{\degphi}{\deg\varphi_E}
\newcommand{\ord}{\operatorname{ord}}
\newcommand{\Phig}{\Phi^{\mathrm{geom}}}
\newcommand{\Gal}{\operatorname{Gal}}
\newcommand{\ad}{\operatorname{ad}^0}
\newcommand{\Sym}{\operatorname{Sym}^2}
\newcommand{\nv}{\mathrm{nv}}
\newcommand{\Frob}{\operatorname{Frob}}
\newcommand{\Fitt}{\operatorname{Fitt}}
\newcommand{\Tam}{\operatorname{Tam}}
\newcommand{\len}{\operatorname{length}}
\newcommand{\Ns}{3\mathrm{Ns}}
\newcommand{\cyc}{\mathrm{cyc}}
\newcommand{\SL}{\mathrm{SL}}
\newcommand{\GL}{\mathrm{GL}}
\setlength{\emergencystretch}{3em}

\title[Local factors of the adjoint and the modular degree]{Local factors of the adjoint $L$-function and the $\ell$-adic valuations of the modular degree}
\author{Dharmaj Soni}
\email{dharmaj.soni@gmail.com}
\urladdr{https://orcid.org/0009-0007-6397-3351}
\date{24 October 2026}

\begin{document}

\begin{abstract}
Let $E/\Q$ be the $X_0(N)$-optimal curve of its isogeny class, $\varphi_E\colon X_0(N)\to E$ its modular parametrisation and $f$ its newform. For an odd prime $\ell$ not dividing $N$ such that the mod $\ell$ representation of $E$ is absolutely irreducible on $\Gal(\overline\Q/\Q(\sqrt{\ell^*}))$ (for $\ell\ge5$ this only asks that $E[\ell]$ be irreducible), we prove
\[
v_\ell(\degphi)\ =\ \len_{\Z_\ell}H^1_f(\Q,\ad T_\ell E\otimes\Q_\ell/\Z_\ell)+\sum_{q\mid N}v_\ell\bigl(\Tam_q(\ad T_\ell E)\bigr)+\sum_{q\mid N}e_\ell(E,q),
\]
where $e_\ell(E,q)$ is the $\ell$-adic valuation of the Euler factor at $s=1$ of the adjoint $L$-function $L(\ad f,s)$ that the naive adjoint $L$-function omits (it is nonzero only at the additive primes: for a tame $q\ge5$ it is $v_\ell(q-\chi_{-3}(q))$ for the Kodaira types $\mathrm{II},\mathrm{II}^*,\mathrm{IV},\mathrm{IV}^*$, $v_\ell(q-\chi_{-4}(q))$ for $\mathrm{III},\mathrm{III}^*$, $v_\ell(q^2-1)$ for $\mathrm{I}_n^*$ and $v_\ell(q-1)+v_\ell((q+1)^2-a_q(E')^2)$ for $\mathrm{I}_0^*$, with $E'$ the twist of good reduction), and the Tamagawa factors are trivial at the potentially good primes and explicit at the potentially multiplicative ones, where their valuation is at least $t_\ell(E,q)$, the $\ell$-part of the valuation of the Tate parameter. In particular
\[
v_\ell(\degphi)\ \ge\ \sum_{q\mid N}e_\ell(E,q)+\sum_{q\mid N}t_\ell(E,q).
\]
The proof assembles the theorem of Diamond, Flach and Guo identifying the length of the $\Sigma$-imprimitive adjoint Selmer group with the valuation of a congruence ideal, their formula for the variation of that ideal with $\Sigma$, the local Tamagawa formula of Fontaine and Perrin-Riou, a computation of the Tamagawa factor of the adjoint at every bad prime, and the observation that for the optimal curve the congruence ideal of Diamond, Flach and Guo is generated by the modular degree, by the projection formula; it also uses an identification of their integral structure with the cohomology of $X_0(N)$, which we prove. The Tamagawa factor raises the bound by at least one whenever a Steinberg prime $q\equiv1\pmod\ell$ with $\ell\mid n_q$ has a Tate parameter whose unit part is an $\ell$-th power. Without the irreducibility over $\Q(\sqrt{\ell^*})$ the terms of the types $\mathrm{I}_0^*$ and $\mathrm{I}_n^*$ at the odd primes still hold, by a twisting argument and the level-raising formula of Diamond, Flach and Guo, but the inequality itself fails: for $\ell=3$, among the curves of the tables without complex multiplication, exactly the curves whose mod $3$ image is the normaliser of a split Cartan subgroup violate it, by exactly one, in $31\%$, $23\%$ and $18\%$ of the cases on three ranges, the one being the length of $H^0(\Q,\ad\rho(1)\otimes\Q_3/\Z_3)$. With that correction term the inequality holds with no exception, with Euler factors computed by PARI at every bad prime, for every curve of conductor at most $500000$ that the tables determine to be optimal, with $E[\ell]$ irreducible and without complex multiplication, and for $893$ curves outside the tables, at $\ell=3,5,7,11,13$, also at primes $\ell$ dividing $N$ where no theorem applies, with equality in $57\%$ of the cases with a positive bound at $\ell=3$. We describe the extra terms at $q=\ell=3$ and the Eisenstein case, where we show that the bound $2$ for the deficit at $\ell=3$ follows from a conjectured bound by the rational torsion.
\end{abstract}

\maketitle

\section{Introduction}

Let $E/\Q$ be an elliptic curve of conductor $N$ and let $f=f_E\in S_2(\Gamma_0(N))$ be the newform attached to it. Among the curves in the isogeny class of $E$ there is a distinguished one, the $X_0(N)$-optimal curve, for which the modular parametrisation $\varphi_E\colon X_0(N)\to E$ has minimal degree; throughout, $E$ denotes this curve and $\degphi$ its \emph{modular degree}. When $E[\ell]$ is irreducible, every curve of the class is linked to $E$ by a cyclic isogeny of degree prime to $\ell$, so $v_\ell$ of the degree of its minimal parametrisation is $v_\ell(\degphi)$: at such an $\ell$ the choice of the optimal curve does not matter. The modular degree is computable \cite{Cremona97,Watkins02}, and its size is tied to deep questions (the conjectured bound $\degphi\ll_\varepsilon N^{2+\varepsilon}$ is equivalent to a form of the $abc$ conjecture, see \cite[\S2]{ARS12} and the references there), but its prime factorisation is not understood. The purpose of this note is to describe, and for the most part prove, a local-to-global identity: the $\ell$-adic valuation of $\degphi$ is the length of an adjoint Selmer group plus local invariants of $E$ at the bad primes, and the right invariants are the local factors of the adjoint $L$-function.

The classical route to the degree is the Petersson norm. Up to explicit factors $\degphi$ is $\langle f,f\rangle$ divided by the area of the period lattice of $E$ (Zagier \cite{Zagier85}, Watkins \cite{Watkins02}), and $\langle f,f\rangle$ is the special value of the naive symmetric square $L$-function at $s=2$ (Rankin \cite{Rankin39}, Shimura \cite{Shimura76}; see \cite[Lemma 2.12]{DFG04}); Remark~\ref{rem:degree} gives the exact formula. Flach's finiteness theorem for the symmetric square \cite{Flach92} is the first place where the modular degree meets the Selmer group of $\Sym f$. By a theorem of Agashe, Ribet and Stein \cite[Theorem 2.2]{ARS12}, $\degphi$ divides the \emph{congruence number} $r_E$, the largest integer $r$ such that $f\equiv g\pmod r$ coefficientwise for some integral cusp form $g$ of level $N$ orthogonal to $f$, with $v_p(\degphi)=v_p(r_E)$ for every prime $p$ with $p^2\nmid N$. The congruence number is the congruence ideal of $f$ in the full Hecke algebra of level $N$ (Lemma~\ref{lem:crit}), and the theory of such ideals, from Hida \cite{Hida81} to Diamond, Flach and Guo \cite{DFG04}, expresses them through the adjoint Selmer group and the adjoint $L$-value. The $L$-function that appears there is the \emph{naive} one, $L^{\nv}(\Sym f,s)=\zeta^{(N)}(2s-2)\sum_{n\ge1}a_{n^2}(f)n^{-s}$ (Shimura \cite{Shimura75}), which has the true local factor at the primes of good or multiplicative reduction and the trivial factor at the additive primes. The point of this note is that the primes at which the naive and the true adjoint $L$-functions differ are exactly the primes at which the modular degree acquires extra local divisibility, that the amount is the valuation of the omitted factor, and that this is a consequence of the theorems of \cite{DFG04} once they are read for elliptic curves, with a precise boundary at the one hypothesis they need.

\subsection{The local terms}\label{sec:localterms}
Let $V=V_\ell E$, $T=T_\ell E$, $\rho=\rho_{E,\ell}$, $\bar\rho$ its reduction, and $A_f=\ad V$ the trace-zero adjoint. Throughout, $\Frob_p$ is an arithmetic Frobenius at $p$ and $\Frob_p^{-1}$ a geometric one; $\chi_\ell$ is the $\ell$-adic cyclotomic character and $\bar\chi_\ell$ its reduction; the quadratic character of $\Q(\sqrt d)$, $d$ a discriminant, is $\chi_d$, so that $\chi_{-3}$, $\chi_{-4}$ and, for an odd prime $q$ and $q^*=(-1)^{(q-1)/2}q$, the character $\chi_{q^*}$ of conductor $q$ appear below; at $\ell=3$ the cyclotomic character is written $\chi_3^{\cyc}$, and $\bar\chi_3^{\cyc}=\chi_{-3}$. We follow the conventions of \cite[\S1.7.2]{DFG04}: for a prime $p\ne\ell$ let $\delta_p=\dim V^{I_p}$, so $\delta_p=2$, $1$, $0$ for good, multiplicative, additive reduction; $L_p(A_f,s)=\det(1-\Frob_p^{-1}p^{-s}\mid A_f^{I_p})^{-1}$ is the true local factor, defined with the geometric Frobenius as in \cite[\S1.1.1]{DFG04}, and the naive factor is $L_p^{\nv}(A_f,s)=L_p(A_f,s)$ if $\delta_p>0$ and $1$ if $\delta_p=0$. (At $p=\ell$ the definitions use $D_{\mathrm{cris}}$; we need them only for $p\ne\ell$.) The matrices in the proof of Lemma~\ref{lem:tamfactor} are those of the arithmetic Frobenius. In the normalisation of the symmetric square, $L_p(\Sym f,s)=L_p(A_f,s-1)$, and PARI's \texttt{lfunsympow(E,2)} returns these factors at every prime, including $2$ and $3$ \cite{PARI}.

\begin{definition}\label{def:terms}
For a bad prime $q$ put
\[
t_\ell(E,q)=\begin{cases}v_\ell(-v_q(j_E))&\text{if }v_q(j_E)<0,\\0&\text{otherwise,}\end{cases}
\]
and, if $q\ne\ell$,
\[
e_\ell(E,q)=v_\ell\Bigl(\frac{L_q(A_f,1)^{-1}}{L_q^{\nv}(A_f,1)^{-1}}\Bigr)=v_\ell\Bigl(\frac{P_q(q^{-2})}{1-a_q(E)^2q^{-2}}\Bigr),
\]
where $P_q(X)$ is the polynomial with $L_q(\Sym f,s)=P_q(q^{-s})^{-1}$. So $e_\ell$ is the valuation of the Euler factor that the naive $L$-function omits; it is $0$ at a multiplicative prime, where the naive factor is the true one. The term $t_\ell$ is the $\ell$-part of the valuation of the Tate parameter at a potentially multiplicative prime, which is $v_\ell(n)$ for a Kodaira symbol $\mathrm{I}_n$ at any $q$ and $\mathrm{I}_n^*$ at odd $q$, while a symbol $\mathrm{I}_n^*$ at $q=2$ can have potentially good reduction, and then $t_\ell=0$. The term $t_\ell$ is invariant under quadratic twist, and $e_\ell(E,q)$ under quadratic twists unramified at $q$; a twist ramified at $q$ can change $e_\ell(E,q)$ (the curve 11a1 has type $\mathrm{I}_5$ at $11$ and $e_3=0$, its twist 121d2 by $\chi_{-11}$ has type $\mathrm{I}_5^*$ and $e_3=v_3(120)=1$), because the true factor is twist invariant and the naive one is not.
\end{definition}

For a tame prime $q\ge5$ the omitted factor is read off the Kodaira type (Lemma~\ref{lem:euler}). With $\chi_{-3}$, $\chi_{-4}$ the quadratic characters of $\Q(\sqrt{-3})$ and $\Q(i)$, $E'$ the quadratic twist of $E$ by $\Q(\sqrt{q^*})$, and $\alpha,\beta$ the roots of $X^2-a_q(E')X+q$ (for type $\mathrm{I}_0^*$, where $E'$ has good reduction at $q$):
\begin{center}\footnotesize\setlength{\tabcolsep}{4pt}
\begin{tabular}{lll}
\toprule
type at $q\ge5$ & $L_q(A_f,s)^{-1}$ & $e_\ell(E,q)$\\
\midrule
$\mathrm{I}_n$ & $1-q^{-1-s}$ & $0$\\
$\mathrm{I}_n^*$ & $1-q^{-1-s}$ & $v_\ell(q^2-1)$\\
$\mathrm{II},\mathrm{II}^*,\mathrm{IV},\mathrm{IV}^*$ & $1-\chi_{-3}(q)q^{-s}$ & $v_\ell(q-\chi_{-3}(q))$\\
$\mathrm{III},\mathrm{III}^*$ & $1-\chi_{-4}(q)q^{-s}$ & $v_\ell(q-\chi_{-4}(q))$\\
$\mathrm{I}_0^*$ & $(1-\tfrac{\alpha}{\beta}q^{-s})(1-q^{-s})(1-\tfrac{\beta}{\alpha}q^{-s})$ & $v_\ell(q-1)+v_\ell\bigl((q+1)^2-a_q(E')^2\bigr)$\\
\bottomrule
\end{tabular}
\end{center}
The row $\mathrm{I}_n$ holds at every $q\ne\ell$, $q=2$ included, the rows $\mathrm{I}_n^*$ and $\mathrm{I}_0^*$ at every odd $q\ne\ell$, and the row $\mathrm{IV}$, $\mathrm{IV}^*$ at every $q\ne3,\ell$ (Lemmas~\ref{lem:IV} and~\ref{lem:euler}); at every bad prime, $2$ and $3$ included, $P_q$ has one of four shapes (Proposition~\ref{prop:shapes}).

\subsection{The theorem}
For a prime $p\ne\ell$ and a $G_{\Q_p}$-stable $\Z_\ell$-lattice $M$ the \emph{Tamagawa factor} is $\Tam_p(M)=\Fitt_{\Z_\ell}\bigl((H^1(I_p,M)_{\mathrm{tor}})^{G_{\Q_p}}\bigr)$, an ideal of $\Z_\ell$ \cite[p.~708]{DFG04}, \cite[I.4.2.2]{FPR94}; for $M=\ad T$ at a prime of potentially multiplicative reduction its valuation is at least $t_\ell(E,p)$, with an explicit excess (Lemma~\ref{lem:tamfactor}), at a prime of potentially good additive reduction it is $0$ (Lemma~\ref{lem:tamadd}), and at a prime of good reduction it is $0$. Let $H^1_f(\Q,\ad T\otimes\Q_\ell/\Z_\ell)$ be the Bloch-Kato Selmer group \cite[\S2.1]{DFG04}.

\begin{theorem}\label{thm:main}
Let $E$ be $X_0(N)$-optimal and let $\ell$ be an odd prime not dividing $2N$ such that $\bar\rho_{E,\ell}$ is absolutely irreducible when restricted to $\Gal(\overline\Q/\Q(\sqrt{\ell^*}))$, $\ell^*=(-1)^{(\ell-1)/2}\ell$. Then
\begin{equation}\label{eq:identity}
v_\ell(\degphi)\ =\ \len_{\Z_\ell}H^1_f(\Q,\ad T\otimes\Q_\ell/\Z_\ell)\ +\ \sum_{q\mid N}v_\ell\bigl(\Tam_q(\ad T)\bigr)\ +\ \sum_{q\mid N}e_\ell(E,q).
\end{equation}
In particular
\begin{equation}\label{eq:thm}
v_\ell(\degphi)\ \ge\ \sum_{q\mid N}e_\ell(E,q)+\sum_{q\mid N}t_\ell(E,q).
\end{equation}
\end{theorem}

The Tamagawa sum is explicit (Corollary~\ref{cor:explicit}), so \eqref{eq:identity} says that the excess of $v_\ell(\degphi)$ over an explicit local bound is the length of the adjoint Selmer group. The hypothesis is the one under which Diamond, Flach and Guo prove the Tamagawa number conjecture for the adjoint motive of $f$ \cite[Theorem 0.1]{DFG04}; in their notation it says that $\ell$ is not in the exceptional set $S_f$. For $\ell\nmid 2N$ it holds if and only if $E[\ell]$ is irreducible and, when $\ell=3$, the mod $3$ image is not the normaliser of a split Cartan subgroup (Proposition~\ref{prop:Sf}). Theorem~\ref{thm:main} is an assembly of their results. Their Theorem 3.7 identifies the length of the $\Sigma$-imprimitive adjoint Selmer group, $\Sigma$ the set of bad primes, with the valuation of the naive $\Sigma$-finite congruence ideal; their Proposition 1.4(c) says that this ideal is the congruence ideal at level $N$ times the naive Euler factors at the multiplicative primes; their Lemma 2.1 compares the $\Sigma$-imprimitive Selmer group with the Bloch-Kato one, the local indices being, by the formula of Fontaine and Perrin-Riou, the true Euler factors times the Tamagawa factors; and the congruence ideal at level $N$ is generated by $\degphi$, by the projection formula for $\varphi_E$ (Proposition~\ref{prop:pullback}). The naive factors at the multiplicative primes cancel, the true factors at the additive primes are the omitted ones, and \eqref{eq:identity} appears. One input is not a published statement in the form we need: the identification of the integral structure of \cite{DFG04} with $H^1(X_0(N),\Z_\ell)$ (Proposition~\ref{prop:integral}), which we prove; it holds integrally, with no localisation. What this note adds is the identification of the congruence ideal with the modular degree, the reading of the local terms for elliptic curves (Lemmas~\ref{lem:euler}, \ref{lem:tamfactor} and~\ref{lem:tamadd}, Proposition~\ref{prop:shapes}), the twisting theorems that hold without the hypothesis on $\Q(\sqrt{\ell^*})$ (Theorem~\ref{thm:quant} and Lemma~\ref{lem:step1}), the identification of the curves for which that hypothesis fails and of what happens to them (Propositions~\ref{prop:ns} and~\ref{prop:ns-structure}, Section~\ref{sec:anomaly}), and the computations, which test \eqref{eq:identity} well beyond the inequality.

\subsection{Beyond the hypotheses}
The data suggest that the inequality survives the loss of the hypotheses, with two modifications: a weight at the prime $\ell$ itself when $\ell\mid N$, and a correction term when $\ell=3$ and the representation is induced from $\Q(\sqrt{-3})$.

\begin{conjecture}\label{conj:euler}
Let $E$ be $X_0(N)$-optimal and without complex multiplication, let $\ell$ be an odd prime such that $E[\ell]$ is irreducible, and let $h_\ell(E)$ be the length of $H^0(\Q,\ad T(1)\otimes\Q_\ell/\Z_\ell)$. Then
\begin{equation}\label{eq:rule}
v_\ell(\degphi)\ \ge\ \sum_{q\mid N,\ q\ne\ell}e_\ell(E,q)\ +\ \sum_{q\mid N}t_\ell(E,q)\ +\ [\ell=3]\,w_3(E)\ -\ h_\ell(E),
\end{equation}
where $w_3(E)$ is the weight of the prime $3$ itself when $3\mid N$ and the reduction at $3$ is potentially good: $1$ for the types $\mathrm{II},\mathrm{IV},\mathrm{III}^*$, $2$ for $\mathrm{II}^*,\mathrm{IV}^*$ and $0$ for $\mathrm{III},\mathrm{I}_0^*$.
\end{conjecture}

The term $h_\ell(E)$ is at most $1$; it is positive exactly when $\bar\rho\simeq\bar\rho\otimes\bar\chi_\ell$, which forces $\ell=3$, where $\bar\chi_3^{\cyc}=\chi_{-3}$; and it vanishes under the hypothesis of Theorem~\ref{thm:main} (Lemma~\ref{lem:h0}). For $\ell=3$ it is positive precisely when the image of $\bar\rho_{E,3}$ is the normaliser of a split Cartan subgroup, with or without complex multiplication, and then $h_3(E)=1$ (Proposition~\ref{prop:ns}); in the tables this is the image labelled $\Ns$ (we use the labels of the tables, in the notation of \cite{Sutherland16}, together with their intrinsic description: for $\Ns$, a $3$-division polynomial that is a product of two quadratics, Proposition~\ref{prop:ns}). For the curves of the tables with this image and without complex multiplication the inequality without the correction fails by exactly one in $31\%$, $23\%$ and $18\%$ of the cases on the three ranges of Section~\ref{sec:data}, so the hypothesis of Theorem~\ref{thm:main} is sharp. The correction is the term that Lemma 2.1 of \cite{DFG04} produces when $H^0(\Q,\ad T(1)\otimes\Q_\ell/\Z_\ell)\ne0$, provided $H^1_f(\Q,V^D)=0$, which is assumed and not known for these curves (Remark~\ref{rem:ns}). The weights $w_3$ at $q=\ell=3$ are empirical (Section~\ref{sec:wild}); at $q=\ell\ge5$ the rule holds with no term, although the data suggest a sharper statement there (Remark~\ref{rem:oq1}). The terms at the wild primes $2$ and $3$ are the ones of Definition~\ref{def:terms} with PARI's Euler factor, and they reproduce every minimum observed in the tables, including the irregular behaviour of the prime $2$ (Section~\ref{sec:wild}). On the four ranges of Section~\ref{sec:data}, \eqref{eq:rule} holds with no exception at $\ell=3,5,7,11,13$, including the curves with $\ell\mid N$, which no theorem covers.

\subsection{Local component groups}
The rule contains two simpler statements, in terms of component groups. For a bad prime $q$ let $\Phig_q=\Phi_q(\overline\F_q)$ be the group of components of the special fibre of the N\'eron model of $E$ at $q$ over the algebraic closure. Its order is read off the Kodaira type: $n$ for type $\mathrm{I}_n$; $1$ for $\mathrm{II},\mathrm{II}^*$; $2$ for $\mathrm{III},\mathrm{III}^*$; $3$ for $\mathrm{IV},\mathrm{IV}^*$; $4$ for $\mathrm{I}_0^*$ and $\mathrm{I}_n^*$ \cite[IV.9]{SilvermanATAEC}. The Tamagawa number $c_q=\#\Phi_q(\F_q)$ counts only the components defined over $\F_q$; the statements below are about $\Phig_q$, not $c_q$. For an odd prime $q$ let $E^{(q)}$ be the quadratic twist of $E$ by $\chi_{q^*}$, that is, by $\Q(\sqrt{q^*})$; it has the same reduction type as $E$ at every prime other than $q$, while at $q$ the Kodaira types are exchanged (Lemma~\ref{lem:twist}):
\[
\mathrm{I}_n\leftrightarrow\mathrm{I}_n^*,\qquad \mathrm{II}\leftrightarrow\mathrm{IV}^*,\qquad \mathrm{IV}\leftrightarrow\mathrm{II}^*,\qquad \mathrm{III}\leftrightarrow\mathrm{III}^*,\qquad \mathrm{I}_0\leftrightarrow\mathrm{I}_0^*
\]
(all five at every odd $q$; at $q=3$ the three middle exchanges keep the conductor exponent, which can exceed $2$). Define the \emph{twist-symmetrised} weight
\[
w_\ell(E,q)=\begin{cases}
v_\ell(n) & \text{if the type at $q$ is $\mathrm{I}_n$, or $\mathrm{I}_n^*$ with $q$ odd},\\
1 & \text{if $\ell=3$ and the type is $\mathrm{IV}$ or $\mathrm{IV}^*$, or $\mathrm{II}$ or $\mathrm{II}^*$ with $q$ odd},\\
0 & \text{otherwise},
\end{cases}
\]
so that
\[
w_\ell(E,q)=\max\bigl(v_\ell(\#\Phig_q(E)),\ v_\ell(\#\Phig_q(E^{(q)}))\bigr)\ \text{for odd }q,\qquad w_\ell(E,2)=v_\ell(\#\Phig_2(E))
\]
for odd $\ell$, by the exchanges above.
The two inequalities are
\begin{equation}\label{eq:main}
v_\ell(\degphi)\ \ge\ \sum_{q\mid N} v_\ell\bigl(\#\Phig_q\bigr),
\end{equation}
\begin{equation}\label{eq:twist}
v_\ell(\degphi)\ \ge\ \sum_{q\mid N} w_\ell(E,q),
\end{equation}
of which the second implies the first. Both hold with no exception on the same ranges, under the hypotheses of Conjecture~\ref{conj:euler} (Section~\ref{sec:forms}).

\begin{proposition}\label{prop:forms}
Under the hypothesis of Theorem~\ref{thm:main}, \eqref{eq:main} and \eqref{eq:twist} hold; the same is true under Conjecture~\ref{conj:euler} with $h_\ell(E)=0$.
\end{proposition}
\begin{proof}
Compare the two sides term by term with \eqref{eq:thm}, respectively \eqref{eq:rule}; $w_\ell(E,q)\ge v_\ell(\#\Phig_q)$, so it suffices to treat \eqref{eq:twist}. At a prime of type $\mathrm{I}_n$ both sides have $v_\ell(n)=t_\ell(E,q)$. At a prime of type $\mathrm{I}_n^*$ with $q$ odd, $t_\ell(E,q)=v_\ell(n)=w_\ell(E,q)$; with $q=2$, $w_\ell(E,2)=v_\ell(4)=0$. At $q\ge5$ of type $\mathrm{II},\mathrm{II}^*,\mathrm{IV},\mathrm{IV}^*$ and $\ell=3$, $e_3(E,q)=v_3(q-\chi_{-3}(q))\ge1=w_3(E,q)$ since $q\equiv\chi_{-3}(q)\pmod3$; for $\ell\ge5$ both weights are $0$. At $q=3$ of these types and $\ell=3$, $w_3(E)\ge1=w_3(E,3)$ by definition, and for $\ell\ge5$ the weight is $0$. At $q=2$ of type $\mathrm{II},\mathrm{II}^*$, $w_\ell(E,2)=0$; of type $\mathrm{IV},\mathrm{IV}^*$ and $\ell=3$, $w_3(E,2)=1=e_3(E,2)$ by Lemma~\ref{lem:IV}. The types $\mathrm{III},\mathrm{III}^*,\mathrm{I}_0^*$ have $w_\ell=0$ for odd $\ell$. All other terms of \eqref{eq:thm} and \eqref{eq:rule} are nonnegative.
\end{proof}
On the $92{,}812$ working-set curves with type $\mathrm{IV}$ or $\mathrm{IV}^*$ at $2$, PARI's factor at $2$ is $1+2X$ for every one of them, as Lemma~\ref{lem:IV} requires.

\subsection{History, and the statements proved without the hypothesis on \texorpdfstring{$\Q(\sqrt{\ell^*})$}{Q(sqrt l*)}}
The multiplicative part of \eqref{eq:main} has a long history. Ribet and Takahashi \cite{RibetTakahashi97,Takahashi01} compared, for semistable $E$, $\degphi$ with the degree of a parametrisation by a Shimura curve and found the product of the $n_q$ over the primes of the quaternion discriminant, exactly when $E$ is the only curve in its isogeny class and otherwise up to primes where $E[\ell]$ is reducible; Pasten \cite{Pasten} controls that error term. Pollack and Weston \cite{PollackWeston11} conjectured, and proved for weight $2$ and square-free level when $\bar\rho$ is ramified at two primes at least and under their hypothesis CR ($\bar\rho$ surjective and ramified at every prime of the relevant part of the level that is $\pm1$ modulo $\ell$) \cite[Theorem 6.11]{PollackWeston11}, and Kim and Ota \cite{KimOta} proved under the hypotheses recalled in Remark~\ref{rem:ko}, that moving a prime $q$ into the ``new'' part of the level changes the congruence ideal of $f_E$ in the full space $S_2(\Gamma_0(N))$ by exactly the \emph{Tamagawa exponent} $t_q$, the largest $t$ with $\rho_{E,\ell}\bmod\ell^t$ unramified at $q$; for multiplicative $q$ one has $t_q=v_\ell(n_q)=t_\ell(E,q)$ whether or not the reduction is split (Lemma~\ref{lem:tamagawa-exp}). Theorem~\ref{thm:main} contains the inequality $v_\ell(\degphi)\ge\sum_{q\,\|\,N}t_\ell(E,q)$ for every multiplicative prime, including the primes $q\equiv\pm1\pmod\ell$ and the prime $\ell=3$, which the formula of Kim and Ota excludes; it does not contain their equality.

Two statements hold under weaker hypotheses than Theorem~\ref{thm:main}, by level lowering and twisting rather than by the Taylor-Wiles method, and they matter exactly where that method does not apply.

\begin{theorem}\label{thm:qual}
Let $E/\Q$ have conductor $N$ and let $\ell$ be an odd prime not dividing $N$ such that $\bar\rho_{E,\ell}$ is irreducible and, if $\ell=3$, absolutely irreducible when restricted to $\Gal(\overline\Q/\Q(\sqrt{-3}))$. Let $q\neq\ell$ be a prime at which $E$ has additive reduction of one of the following kinds:
\begin{enumerate}
\item[(a)] type $\mathrm{IV}$ or $\mathrm{IV}^*$, with $q\ge5$ or $q=2$, and $\ell=3$;
\item[(b)] type $\mathrm{II}$ or $\mathrm{II}^*$, with $q\ge5$ and $\ell=3$;
\item[(c)] type $\mathrm{I}_n^*$ with $\ell\mid n$, $q$ odd.
\end{enumerate}
Then $\ell$ divides the modular degree of the $X_0(N)$-optimal curve in the isogeny class of $E$.
\end{theorem}

For $\ell\ge5$ Theorem~\ref{thm:qual} needs only the irreducibility of $\bar\rho$; its proof (Section~\ref{sec:proof}) lowers the level of $\bar\rho$, or of its twist $\bar\rho\otimes\chi_{q^*}$, and twists back. The quantitative version of case (c), and of the type $\mathrm{I}_0^*$, is the following, which needs no hypothesis beyond the irreducibility of $E[\ell]$ for the inequality, and gives an equality under the hypothesis of Theorem~\ref{thm:main}.

\begin{theorem}\label{thm:quant}
Let $\ell$ be an odd prime with $\ell\nmid 2N$ and $E[\ell]$ irreducible, let $S$ be the set of odd primes $q$ at which $E$ has type $\mathrm{I}_0^*$ or $\mathrm{I}_n^*$, and let $E_S$ be the twist of $E$ by $\chi_S=\prod_{q\in S}\chi_{q^*}$. Then
\[
v_\ell(\degphi)\ \ge\ v_\ell(r_{E_S})+\sum_{q\in S}e_\ell(E,q)\ \ge\ \sum_{q\in S}e_\ell(E,q).
\]
If moreover $\ell\notin S_f$ (for $\ell\ge5$ this is automatic, Proposition~\ref{prop:Sf}), the first inequality is an equality, and $v_\ell(r_{E_S})=v_\ell(\deg\varphi_{E_S'})$ for the optimal curve $E_S'$ of the isogeny class of $E_S$:
\[
v_\ell(\degphi)\ =\ v_\ell(\deg\varphi_{E_S'})+\sum_{q\in S}e_\ell(E,q).
\]
\end{theorem}

The proof (Section~\ref{sec:twisting}) uses the formula of \cite[Proposition 1.4(c)]{DFG04} for the depleted form of $f_{E_S}$, which lives at level $N$, and an elementary transport of congruences along the twist, which identifies $v_\ell(r_E)$ with the Hecke congruence ideal of that depleted form at every prime $\ell$ outside $S$, with no other hypothesis on $\ell$ (Lemma~\ref{lem:step1}); under the hypothesis of Theorem~\ref{thm:main} the identity \eqref{eq:identity} for $E$ and for $E_S'$ gives the equality, with no multiplicity-one hypothesis.

\begin{proposition}\label{prop:ns}
Let $E/\Q$ have $E[3]$ irreducible. Then $\bar\rho_{E,3}$ is absolutely irreducible on $\Gal(\overline\Q/\Q(\sqrt{-3}))$ unless its image is conjugate to the normaliser $N_s(\F_3)$ of a split Cartan subgroup of $\GL_2(\F_3)$, and in that case $\bar\rho_{E,3}\otimes\F_9$ is induced from a character of order $4$ of $\Gal(\overline\Q/\Q(\sqrt{-3}))$, $\bar\rho_{E,3}\simeq\bar\rho_{E,3}\otimes\chi_{-3}$, $H^0(\Q,\ad\bar\rho_{E,3}(1))\ne0$ and $h_3(E)=1$. The $3$-division polynomial of $E$ then factors as a product of two irreducible quadratics, and otherwise it is irreducible.
\end{proposition}

The last statement concerns the Eisenstein case, in which \eqref{eq:main} can fail.

\begin{conjecture}\label{conj:eis}
Let $E$ be optimal, not CM, with a rational $\ell$-isogeny, $\ell$ odd, and let $D$ be the deficit $\sum_q v_\ell(\#\Phig_q)-v_\ell(\degphi)$ of \eqref{eq:main}. Then
\[
D\ \le\ v_\ell\bigl(\#E(\Q)_{\mathrm{tors}}\bigr)+\bigl[\,E[3]\simeq\Z/3\oplus\mu_3\,\bigr],
\]
where $[\,\cdot\,]$ is $1$ if the condition holds and $0$ otherwise, and also $D\le$ the number of curves in the isogeny class of $E$ with a rational point of order $\ell$. In particular $D\le0$ when $E$ has no rational $\ell$-torsion point, $D\le1$ for $\ell\ge5$, and $D\le2$ for $\ell=3$, the last because a curve without complex multiplication with $E[3]\simeq\Z/3\oplus\mu_3$ has no rational point of order $9$ (Proposition~\ref{prop:no9}).
\end{conjecture}

Section~\ref{sec:local} collects the local facts, Section~\ref{sec:proof} proves Theorem~\ref{thm:qual}, Section~\ref{sec:quant} proves Theorems~\ref{thm:main} and~\ref{thm:quant}, Section~\ref{sec:boundary} proves Propositions~\ref{prop:ns} and~\ref{prop:ns-structure} and explains the correction term, Section~\ref{sec:data} describes the computations and Section~\ref{sec:two} discusses the prime $2$, where the analogous statements overlap with the literature on Watkins's conjecture.

\section{Local invariants}\label{sec:local}

Let $q$ be a prime of bad reduction. We use Tate's algorithm \cite{Tate75}, \cite[IV.9]{SilvermanATAEC} for the Kodaira type and $c_q$, and the following facts.

\begin{lemma}\label{lem:tamagawa-exp}
Let $q$ be a prime of multiplicative reduction, $n_q=v_q(\Delta_E)$, and $\ell\ne q$ a prime. Then $\rho_{E,\ell}\bmod\ell^t$ is unramified at $q$ if and only if $\ell^t\mid n_q$. This holds for split and for nonsplit reduction.
\end{lemma}

\begin{proof}
Over the unramified quadratic extension of $\Q_q$ (which does not change the inertia group) $E$ is a Tate curve with parameter $q_E$, $v_q(q_E)=n_q$. The $\ell^t$-torsion is generated by $\mu_{\ell^t}$ and an $\ell^t$-th root of $q_E$. By Kummer theory over $\Q_q^{\mathrm{ur}}$, which contains $\mu_{\ell^t}$ as $q\ne\ell$, inertia acts trivially on $E[\ell^t]$ if and only if $q_E$ is an $\ell^t$-th power in $\Q_q^{\mathrm{ur}}$; the units of its valuation ring are $\ell^t$-th powers by Hensel's lemma, so this holds if and only if $\ell^t\mid n_q$.
\end{proof}

At an additive prime $q\ne\ell$ the inertia invariants of the torsion are those of the component group: $E[\ell^k]^{I_q}\simeq\Phig_q[\ell^k]$ for every $k$. Indeed $E[\ell^k]^{I_q}=E(\Q_q^{\mathrm{ur}})[\ell^k]$; the group $E(\Q_q^{\mathrm{ur}})$ is an extension of $\Phig_q$ by the group $E_0(\Q_q^{\mathrm{ur}})$ of points reducing to the identity component \cite[IV.9]{SilvermanATAEC}, and $E_0(\Q_q^{\mathrm{ur}})$ is an extension of the additive group of $\overline\F_q$ by a pro-$q$ group, so it is uniquely $\ell$-divisible; the snake lemma for multiplication by $\ell^k$ gives the isomorphism \cite[\S1]{SerreTate68}.

\begin{lemma}\label{lem:IV}
Let $q\ne3$ be a prime at which $E$ has reduction of type $\mathrm{IV}$ or $\mathrm{IV}^*$. Then the reduction at $q$ is potentially good, the image of the inertia group $I_q$ in $\operatorname{Aut}(T_\ell E)$ is cyclic of order $3$ for every prime $\ell\ne q$, and its image in $\operatorname{Aut}(E[3])$ is a nontrivial unipotent subgroup. Consequently $E[3]^{I_q}$ is a line, $\bar\rho_{E,3}$ is tamely ramified at $q$, the Serre conductor of $\bar\rho_{E,3}$ has exponent $1$ at $q$ and the conductor of $E$ has exponent $2$ at $q$. Moreover, for $\ell\ne q$, $(\ad V_\ell E)^{I_q}$ is a line on which $\Frob_q$ acts by $\chi_{-3}(q)$, so that $P_q(X)=1-\chi_{-3}(q)qX$ and $e_\ell(E,q)=v_\ell(q-\chi_{-3}(q))$; in particular $P_2(X)=1+2X$ and $e_3(E,2)=1$ when $q=2$.
\end{lemma}

\begin{proof}
By the remark above, $E[3]^{I_q}\simeq\Phig_q[3]\simeq\Z/3$ is a line. Inertia acts nontrivially on $E[3]$, since otherwise $\Phig_q[3]$ would have order $9$. If the reduction were potentially multiplicative, $E$ would be over $\Q_q$ the twist of a Tate curve by a quadratic character $\chi$, ramified since the reduction is additive \cite[V.5]{SilvermanATAEC}; then $\rho_{E,3}(\sigma)=\chi(\sigma)\begin{psmallmatrix}1&*\\0&1\end{psmallmatrix}$ on $I_q$ in a suitable basis, the corner being $n$ times the $3$-adic tame character; the subgroup of $I_q$ on which the corner is $0$ modulo $3$ has index dividing $3$, so it is not contained in the kernel of $\chi$, and $\bar\rho_{E,3}(I_q)$ would contain $-1$, which fixes no line. So the reduction is potentially good, and $\rho_{E,3}(I_q)$ is finite \cite[\S2]{SerreTate68}. It lies in $\SL_2(\Z_3)$, since $\chi_3^{\cyc}$ is unramified at $q\ne3$, and reduction modulo $3$ is injective on it, the kernel of $\GL_2(\Z_3)\to\GL_2(\F_3)$ being torsion-free. The elements of $\SL_2(\F_3)$ fixing a nonzero vector form a group of order $3$ of unipotent elements; so $\bar\rho_{E,3}(I_q)$ is that group and $\rho_{E,3}(I_q)$ is cyclic of order $3$. The kernel of $\rho_{E,\ell}|_{I_q}$ does not depend on $\ell\ne q$ \cite[\S2]{SerreTate68}, so the same holds on $T_\ell E$. The order $3$ is prime to $q$, so the image is a quotient of the tame inertia group and the ramification is tame. The Serre conductor exponent of $\bar\rho_{E,3}$ at $q$ is $\dim E[3]-\dim E[3]^{I_q}=1$, and the conductor exponent of $E$ is $\dim V-\dim V^{I_q}=2$, there being no wild part. Let $\sigma$ generate $\rho(I_q)$; its eigenvalues are the two primitive cube roots of unity, so $\sigma^2+\sigma+1=0$, and the trace-zero endomorphisms commuting with $\sigma$ form the line spanned by $\sigma-\sigma^{-1}=2\sigma+1$. The arithmetic Frobenius acts on the tame quotient of $I_q$ by $x\mapsto x^q$, so $\rho(\Frob_q)\sigma\rho(\Frob_q)^{-1}=\sigma^q=\sigma^{\chi_{-3}(q)}$, and $\Frob_q$, and therefore also $\Frob_q^{-1}$, acts on the line by $\chi_{-3}(q)$. Hence $L_q(A_f,s)^{-1}=1-\chi_{-3}(q)q^{-s}$, $P_q(X)=1-\chi_{-3}(q)qX$, and, the naive factor being $1$, $e_\ell(E,q)=v_\ell(1-\chi_{-3}(q)q^{-1})=v_\ell(q-\chi_{-3}(q))$. For $q=2$: $\chi_{-3}(2)=-1$ and $e_3(E,2)=v_3(3/2)=1$.
\end{proof}

At every prime $q\ne3$ at which $\rho_{E,3}(I_q)$ is cyclic of order $3$, $q=2$ included, a generator $\sigma$ satisfies $\sigma^2+\sigma+1=0$ on $T_3E$, so $T_3E$ is a torsion-free module of $\Z_3$-rank $2$ over the discrete valuation ring $\Z_3[X]/(X^2+X+1)\simeq\Z_3[\zeta_3]$, with $X$ acting as $\sigma$, hence free of rank one over it; Lemma~\ref{lem:tamadd} uses this, also at $q=2$. For $q\ge5$ the order $3$ is also the semistability defect $12/\gcd(v_q(\Delta_E),12)$ of Kraus \cite{Kraus90}; at $q=2$ the conductor exponent $2$ can also be read off \cite{DokDok15}. At $q=3$ the types $\mathrm{IV}$, $\mathrm{IV}^*$ are wildly ramified, and the lemma does not apply.

\begin{remark}\label{rem:types}
For the other additive types at $q\ge5$ the image of inertia mod $3$ is not unipotent: it is $\{\pm1\}$ (type $\mathrm{I}_0^*$), cyclic of order $4$ (types $\mathrm{III}$, $\mathrm{III}^*$), minus a unipotent, of order $6$ (types $\mathrm{II}$, $\mathrm{II}^*$), or a ramified quadratic character times a unipotent (type $\mathrm{I}_n^*$). For all these types $E[3]^{I_q}=0$ and the Serre conductor of $\bar\rho_{E,3}$ has exponent $2$ at $q$. The twist by $\chi_{q^*}$ lowers that exponent for $\mathrm{I}_0^*$, $\mathrm{I}_n^*$, $\mathrm{II}$, $\mathrm{II}^*$; it lowers the conductor of $E$ only for $\mathrm{I}_0^*$ and $\mathrm{I}_n^*$.
\end{remark}

\begin{lemma}\label{lem:twist}
Let $q$ be an odd prime, $\chi=\chi_{q^*}$ the quadratic character of conductor $q$ and $E^{(q)}$ the twist of $E$ by $\chi$, that is, by $\Q(\sqrt{q^*})$.
\begin{enumerate}
\item[(i)] At every prime $p\ne q$ the curves $E$ and $E^{(q)}$ have the same Kodaira type, the same geometric component group and the same conductor exponent.
\item[(ii)] If $E$ has type $\mathrm{I}_n$ at $q$ ($n\ge0$) then $E^{(q)}$ has type $\mathrm{I}_n^*$ at $q$, and conversely, with the same $n$. If $E$ has type $\mathrm{II}$, $\mathrm{III}$, $\mathrm{IV}$ at $q$ then $E^{(q)}$ has type $\mathrm{IV}^*$, $\mathrm{III}^*$, $\mathrm{II}^*$ respectively, and conversely. For the exchanges $\mathrm{I}_n\leftrightarrow\mathrm{I}_n^*$ ($n\ge0$) at every odd $q$, and for all the exchanges at $q\ge5$, the conductor exponent of $E^{(q)}$ at $q$ is $2$ when its reduction is additive and $1$ when it is of type $\mathrm{I}_n$, $n\ge1$. At $q=3$ the exchanges $\mathrm{II}\leftrightarrow\mathrm{IV}^*$, $\mathrm{III}\leftrightarrow\mathrm{III}^*$, $\mathrm{IV}\leftrightarrow\mathrm{II}^*$ leave the conductor exponent unchanged; it can exceed $2$.
\item[(iii)] (Atkin-Li \cite{AtkinLi78}) Let $g$ be a newform of weight $2$ on $\Gamma_0(M)$ with $q^2\nmid M$. Then the twist $g\otimes\chi$, with coefficients $a_n(g\otimes\chi)=\chi(n)a_n(g)$, is a newform of weight $2$ on $\Gamma_0(M')$, $M'=\operatorname{lcm}(M,q^2)$; that is, $M'=Mq^2$ if $q\nmid M$ and $M'=Mq$ if $q\,\|\,M$.
\end{enumerate}
\end{lemma}

\begin{proof}
(i) The character $\chi$ is unramified at $p$, so $E$ and $E^{(q)}$ become isomorphic over an unramified extension of $\Q_p$, and the Kodaira type, the geometric component group and the conductor exponent are invariant under unramified base change. (ii) If $(a_1,a_2,a_3,a_4,a_6)$ is a minimal model of $E$ at $q$ with invariants $c_4,c_6,\Delta$, the twist has a model with invariants $(q^*)^{2}c_4$, $(q^*)^{3}c_6$, $(q^*)^{6}\Delta$ and the same $j$-invariant; only their $q$-adic valuations are used below. For potentially multiplicative reduction, $v_q(j)<0$, the type is $\mathrm{I}_n$ or $\mathrm{I}_n^*$ with $n=-v_q(j)$ according as the reduction is multiplicative or additive, and a ramified quadratic twist exchanges the two. For $n=0$: a model $y^2=x^3+a_2x^2+a_4x+a_6$ with good reduction at $q$ has twist $y^2=x^3+q^*a_2x^2+(q^*)^2a_4x+(q^*)^3a_6$, for which the cubic of Tate's algorithm, after removing the powers of $q$, is $T^3+\epsilon a_2T^2+a_4T+\epsilon a_6$ with $\epsilon=q^*/q=\pm1$; its roots are $\epsilon$ times those of $x^3+a_2x^2+a_4x+a_6$, so they are distinct modulo $q$, the criterion for type $\mathrm{I}_0^*$; conversely a curve of type $\mathrm{I}_0^*$ at an odd $q$ has a model of the latter shape, and its twist has good reduction. For $q\ge5$ and potentially good reduction the type is determined by $v_q(\Delta)$ of a minimal model, which lies in $\{0,2,3,4,6,8,9,10\}$ for the types $\mathrm{I}_0,\mathrm{II},\mathrm{III},\mathrm{IV},\mathrm{I}_0^*,\mathrm{IV}^*,\mathrm{III}^*,\mathrm{II}^*$; the twisted model has $v_q(\Delta)$ increased by $6$, and is minimal exactly when the result is less than $12$, otherwise $12$ is subtracted. For the exchanges $\mathrm{I}_n\leftrightarrow\mathrm{I}_n^*$ at every odd $q$ and for all exchanges at $q\ge5$ the ramification is tame: $V_\ell E^{(q)}=V_\ell E\otimes\chi$ with $\chi$ tame and $V_\ell E$ multiplicative, good, or (for $q\ge5$) tame; so the conductor exponents are the tame ones, $\dim V-\dim V^{I_q}$. At $q=3$, for the types $\mathrm{II},\mathrm{III},\mathrm{IV}$ and their duals, the ramification can be wild, and we use the formula of Ogg and Saito \cite{Ogg,Saito88}, \cite[IV.11]{SilvermanATAEC}: $v_3(\Delta_{\min})=f_3+m-1$, where $f_3$ is the conductor exponent and $m$ the number of irreducible components of the special fibre of the minimal regular model, which is $1,2,3,5,7,8,9$ for the types $\mathrm{II},\mathrm{III},\mathrm{IV},\mathrm{I}_0^*,\mathrm{IV}^*,\mathrm{III}^*,\mathrm{II}^*$. The reduction of $E^{(3)}$ is potentially good, as $j$ is unchanged, and not of type $\mathrm{I}_0$ or $\mathrm{I}_0^*$, since by the case $n=0$ above that would make $E$ of type $\mathrm{I}_0^*$ or $\mathrm{I}_0$. The character $\chi_{-3}$ is tame, so it does not change the Swan conductor; and neither $V_\ell E$ nor $V_\ell E^{(3)}$ has nonzero inertia invariants, since the image of $I_3$ on each is not contained in $\{\pm1\}$ (otherwise a quadratic twist would have good reduction and the type would be $\mathrm{I}_0$ or $\mathrm{I}_0^*$), so it contains an element of order $3$, $4$ or $6$, which has no eigenvalue $\pm1$. Hence $f_3(E^{(3)})=f_3(E)$. The twisted model has discriminant $3^6\Delta$, so $v_3(\Delta_{\min})$ changes by $6$ modulo $12$; with $f_3$ fixed, Ogg's formula gives $m(E^{(3)})-1\equiv m(E)-1+6\pmod{12}$, and as $0\le m-1\le8$ on both sides, $m\mapsto m\pm6$: $1\leftrightarrow7$, $2\leftrightarrow8$, $3\leftrightarrow9$, which are the three exchanges (the types with $m=7,8,9$ and potentially good reduction being $\mathrm{IV}^*$, $\mathrm{III}^*$, $\mathrm{II}^*$, as $\mathrm{I}_n^*$ with $n\ge1$ is potentially multiplicative at odd $q$). For the classification at $3$ compare \cite{Kraus89,Papadopoulos93}. (iii) is \cite[Theorem 3.1]{AtkinLi78}; in the case $q\,\|\,M$, the local component of $g$ at $q$ is an unramified twist of the Steinberg representation, whose twist by a ramified quadratic character has conductor exponent $2$.
\end{proof}

\begin{lemma}[the omitted factors at tame primes]\label{lem:euler}
Let $q\ge5$ be a bad prime, $q\ne\ell$. Then $L_q(A_f,s)^{-1}$ and $e_\ell(E,q)$ are as in the table of Section~\ref{sec:localterms}. The rows $\mathrm{I}_n^*$ and $\mathrm{I}_0^*$ hold for every odd prime $q\ne\ell$, the row $\mathrm{I}_n$ for every $q\ne\ell$.
\end{lemma}

\begin{proof}
The inertia group acts on $V$ through a finite cyclic quotient when the reduction is potentially good, and through a quadratic character times a unipotent when it is potentially multiplicative; $L_q(A_f,s)^{-1}=\prod(1-\lambda q^{-s})$ over the eigenvalues $\lambda$ of the geometric Frobenius $\Frob_q^{-1}$ on $(\ad V)^{I_q}$, in the normalisation in which a good prime has the eigenvalues $\alpha/\beta,1,\beta/\alpha$ on $\ad V$.

\emph{Type $\mathrm{I}_n$.} The local component of $f$ is an unramified twist of the Steinberg representation; $(\ad V)^{I_q}$ is a line on which the geometric Frobenius has eigenvalue $q^{-1}$, so $L_q(A_f,s)^{-1}=1-q^{-1-s}$, which is also the naive factor since $a_q^2=1$: $e_\ell=0$. This holds at every $q\ne\ell$. \emph{Type $\mathrm{I}_n^*$, $q$ odd.} $V=V'\otimes\chi_{q^*}$ with $V'$ Steinberg at $q$ (Lemma~\ref{lem:twist}), $\ad V=\ad V'$ and $a_q=0$, so the naive factor is $1$ and $e_\ell=v_\ell(1-q^{-2})=v_\ell(q^2-1)$. \emph{Types $\mathrm{II},\mathrm{II}^*,\mathrm{IV},\mathrm{IV}^*$.} The image of $I_q$ in $\operatorname{Aut}V$ is cyclic of order $3$ or $6$ \cite{Kraus90}, generated by $\sigma$ or $-\sigma$ with $\sigma$ of order $3$ and eigenvalues $\zeta_3^{\pm1}$. On $\ad V$ the sign disappears and the invariants form the line spanned by $2\sigma+1$, whose square is $-3$, inside $\operatorname{End}_{I_q}(T)=\Z_\ell[\sigma]\simeq\Z_\ell\otimes\Z[\zeta_3]$ (not a field when $\ell\equiv1\pmod3$); Frobenius centralises $\sigma$ if $q\equiv1\pmod3$ and inverts it if $q\equiv2$, for either choice of Frobenius, so it acts on the line by $\chi_{-3}(q)$, and $e_\ell=v_\ell(1-\chi_{-3}(q)q^{-1})=v_\ell(q-\chi_{-3}(q))$. \emph{Types $\mathrm{III},\mathrm{III}^*$.} The image has order $4$, generated by $\tau$ with eigenvalues $\pm i$; the invariant line of $\ad V$ is spanned by $\tau$ itself, which has trace $0$ (it is diagonalisable only over $\Q_\ell(i)$), and Frobenius acts on it by $\chi_{-4}(q)$, so $e_\ell=v_\ell(q-\chi_{-4}(q))$. \emph{Type $\mathrm{I}_0^*$, $q$ odd.} $V=V'\otimes\chi_{q^*}$ with $V'$ unramified with arithmetic Frobenius eigenvalues $\alpha,\beta$; $\ad V$ is unramified with eigenvalues $\alpha/\beta,1,\beta/\alpha$ for either Frobenius, and
\[
L_q(A_f,1)^{-1}=\Bigl(1-\frac{\alpha}{\beta q}\Bigr)\Bigl(1-\frac1q\Bigr)\Bigl(1-\frac{\beta}{\alpha q}\Bigr)=\frac{q-1}{q}\cdot\frac{(q+1)^2-a_q(E')^2}{q^2},
\]
using $\alpha\beta=q$ and $\alpha+\beta=a_q(E')$. Only Lemma~\ref{lem:twist} for the exchanges $\mathrm{I}_n\leftrightarrow\mathrm{I}_n^*$ and $\mathrm{I}_0\leftrightarrow\mathrm{I}_0^*$ was used in the last two rows, and it holds at every odd $q$; at $q=3$ the formulas read $L_3(A_f,1)^{-1}=\frac89$ for $\mathrm{I}_n^*$ and $\frac23\cdot\frac{16-a_3(E')^2}9$ for $\mathrm{I}_0^*$. In each case the values were confirmed against PARI's \texttt{lfunsympow} on every curve of the tables.
\end{proof}

The shape of the factor at every bad prime, $2$ and $3$ included, is determined by the image of inertia.

\begin{proposition}\label{prop:shapes}
Let $q$ be a bad prime and $\ell\ne q$ an odd prime. Then $P_q(X)$ has one of the following four shapes.
\begin{enumerate}
\item[(i)] $P_q(X)=1-X$, if $v_q(j_E)<0$; then $e_\ell(E,q)=0$ if the reduction is multiplicative and $e_\ell(E,q)=v_\ell(q^2-1)$ if it is additive.
\item[(ii)] $P_q(X)=(1-qX)(1-(a^2-2q)X+q^2X^2)$, if the reduction is potentially good and $\rho(I_q)=\{\pm1\}$; here $a=a_q(E')$ for a quadratic twist $E'$ of $E$ with good reduction at $q$, and $e_\ell(E,q)=v_\ell(q-1)+v_\ell((q+1)^2-a^2)$.
\item[(iii)] $P_q(X)=1-cqX$ with $c=\pm1$, if the reduction is potentially good and $\rho(I_q)$ is cyclic of order $3$, $4$ or $6$; here $c$ is the eigenvalue of Frobenius on the line $(\ad V)^{I_q}$, and $e_\ell(E,q)=v_\ell(q-c)$.
\item[(iv)] $P_q(X)=1$, if the reduction is potentially good and $\rho(I_q)$ is not abelian, which happens only for $q=2,3$; then $e_\ell(E,q)=0$.
\end{enumerate}
\end{proposition}

\begin{proof}
(i) Over $\Q_q$, $\rho\simeq\rho'\otimes\chi$ with $\rho'$ the representation of a Tate curve, possibly twisted by an unramified character, and $\chi$ quadratic \cite[V.5]{SilvermanATAEC}; so $\ad V\simeq\ad V'$, inertia acts on $V'$ through a nontrivial unipotent, and as in the proof of Lemma~\ref{lem:tamfactor} the invariants of $\ad V$ are the line spanned by $e$, on which the geometric Frobenius acts by $q^{-1}$. So $L_q(A_f,s)^{-1}=1-q^{-1-s}$ and $P_q(X)=L_q(A_f,s-1)^{-1}=1-X$ with $X=q^{-s}$; the naive factor is the true one if $\delta_q=1$ and $1$ if $\delta_q=0$. In the other cases $\rho(I_q)$ is a finite subgroup of $\SL_2(\Z_\ell)$ whose elements have characteristic polynomials in $\Z[X]$ \cite[\S2]{SerreTate68}, so their orders are $1,2,3,4,6$; the reduction is additive, $\delta_q=0$, and the naive factor is $1$. (ii) The quadratic character $\psi=\rho|_{I_q}$ of $I_q$ is invariant under conjugation by $G_{\Q_q}$, and it extends to a character $\chi$ of $G_{\Q_q}=I_q\rtimes\langle\phi\rangle$, $\phi$ a Frobenius lift generating a complement isomorphic to $\widehat\Z$, by $\chi(i\phi^k)=\psi(i)$. Every quadratic character of $G_{\Q_q}$ is the restriction of a quadratic Dirichlet character $\chi'$ (for $q$ odd a product of $\chi_{q^*}$ and a character unramified at $q$; for $q=2$ a product of $\chi_{-4}$, $\chi_8$ and $\chi_5$), and $E'=E\otimes\chi'$ has good reduction at $q$. Then $\ad V\simeq\ad V_\ell E'$ is unramified with eigenvalues $\alpha/\beta,1,\beta/\alpha$, $\alpha+\beta=a$, $\alpha\beta=q$, so $P_q(X)=(1-\alpha^2X)(1-qX)(1-\beta^2X)$, which is the stated polynomial; $a$ is determined up to sign, and $a^2$ is well defined. The value of $e_\ell$ is the computation of Lemma~\ref{lem:euler}. (iii) A finite abelian subgroup of $\SL_2$ in characteristic $0$ is cyclic; let $g$ generate $\rho(I_q)$, of order $m\in\{3,4,6\}$, with eigenvalues $\zeta^{\pm1}$, $\zeta\ne\zeta^{-1}$. The trace-zero endomorphisms commuting with $g$ form the line spanned by $g-g^{-1}$. Frobenius normalises $\rho(I_q)$, so it conjugates $g$ to some $g^k$ with the same eigenvalues, that is, to $g^{\pm1}$, and it acts on the line by $c=\pm1$, the same sign for the arithmetic and the geometric Frobenius. So $L_q(A_f,s)^{-1}=1-cq^{-s}$, $P_q(X)=1-cqX$ and $e_\ell=v_\ell(1-c/q)=v_\ell(q-c)$. (iv) A faithful two-dimensional representation of a non-abelian finite group in characteristic $0$ is absolutely irreducible, so the commutant of $\rho(I_q)$ consists of scalars and $(\ad V)^{I_q}=0$. For $q\ge5$ the image is tame, hence cyclic.
\end{proof}

\begin{remark}\label{rem:wild}
At $q=2,3$ the Kodaira type does not determine $P_q$, but PARI does, and Definition~\ref{def:terms} applies verbatim. For instance a type $\mathrm{IV}$ or $\mathrm{IV}^*$ at $2$ always has $P_2(X)=1+2X$ and $e_3=1$ (Lemma~\ref{lem:IV}); a type $\mathrm{II}$ or $\mathrm{II}^*$ at $2$ usually has $P_2=1$, but 336f1 (type $\mathrm{II}$ at $2$) has $P_2=1+2X$, $e_3(E,2)=1$ and $\degphi=24$, and 1232l1 (type $\mathrm{II}^*$ at $2$), the twist of 77a1 by $\chi_{-4}$ with $a_2(\text{77a1})=0$, has $P_2=(1-2X)(1+2X)^2=1+2X-4X^2-8X^3$, $e_3(E,2)=2$ and $\degphi=288$, so that the bound of Theorem~\ref{thm:main} at $\ell=3$ is attained by the prime $2$ alone. A type $\mathrm{I}_n^*$ at $2$ has potentially good reduction in a little over half of the cases (below conductor $10{,}000$: $58\%$, $2{,}340$ of $4{,}037$, among the optimal curves, and $61\%$, $5{,}935$ of $9{,}667$, among all curves), usually with $P_2=1$, and then $t_\ell=0$. Among all curves of conductor below $10{,}000$, every cubic factor at $2$ ($1{,}143$) and at $3$ ($2{,}283$) equals the formula of (ii) for a twist of good reduction found independently ($534$ and $1{,}308$ among the optimal non-CM curves); that a cubic factor has shape (ii) is automatic, (ii) being the only cubic shape. The prime $2$ can contribute at $\ell\ge5$: 592c1, the twist of 37a1 by $\chi_{-4}$ with $a_2(\text{37a1})=-2$, has type $\mathrm{II}^*$ at $2$, $P_2=(1-2X)(1+4X^2)$, $e_5(E,2)=1$ and $\degphi=80=2^4\cdot5$. At $q=\ell$ the Euler factor is not the relevant quantity, and the weights $w_3$ of Conjecture~\ref{conj:euler} are empirical.
\end{remark}

\begin{lemma}[the Tamagawa factor at a multiplicative prime]\label{lem:tamfactor}
Let $q$ be a prime of multiplicative reduction, $n=n_q$, $\ell$ an odd prime different from $q$, $T=T_\ell E$ and $M=\ad T$. Let $q_E\in\Q_q^\times$ be the Tate parameter, $u=q_Eq^{-n}\in\Z_q^\times$ its unit part, and $c(u)$ the class of the reduction of $u$ in $\F_q^\times/(\F_q^\times)^{\ell^{v}}$, $v=v_\ell(q-1)$, identified with an element of $\Z/\ell^{v}$ by the choice of a generator; $v_\ell(c(u))$ does not depend on that choice, $v_\ell(0)$ is read as $v$, and $v_\ell(c(u))\ge k$ if and only if the reduction of $u$ is an $\ell^k$-th power in $\F_q^\times$, for $k\le v$. Then $H^1(I_q,M)_{\mathrm{tor}}\simeq(\Z/\ell^{v_\ell(n)})^2$ and
\[
v_\ell\bigl(\Tam_q(\ad T)\bigr)=v_\ell(n)+\min\bigl(v_\ell(n),v_\ell(q-1),v_\ell(c(u))\bigr).
\]
In particular $v_\ell(\Tam_q(\ad T))\ge t_\ell(E,q)$, with equality unless $q\equiv1\pmod\ell$, $\ell\mid n$ and the reduction of $u$ is an $\ell$-th power in $\F_q^\times$. The same formula holds at a prime $q\ne\ell$ of additive potentially multiplicative reduction, $q=2$ included, with $n=-v_q(j_E)$ and $q_E$ the Tate parameter, determined by $j(q_E)=j_E$.
\end{lemma}

\begin{proof}
Over the unramified quadratic extension of $\Q_q$ the curve is a Tate curve, so $T$ has a basis $e_1,e_2$ in which $G=G_{\Q_q}$ acts through matrices $\begin{psmallmatrix}\chi_\ell\varepsilon&\varepsilon\kappa\\0&\varepsilon\end{psmallmatrix}$, with $\varepsilon$ unramified of order at most $2$ and $\kappa$ the Kummer cocycle of $q_E$ in $H^1(G,\Z_\ell(1))=\varprojlim\Q_q^\times/(\Q_q^\times)^{\ell^k}$; the nonsplit curve is the Tate curve twisted by $\varepsilon$, and $\varepsilon\kappa$, not $\kappa$, satisfies the cocycle relation for the twisted matrices. The inertia group $I=I_q$ acts through its maximal pro-$\ell$ tame quotient, generated by an element $t$ with $te_1=e_1$, $te_2=e_2+n'e_1$, $v_\ell(n')=v_\ell(n)$, since the restriction of $\kappa$ to $I$ is $n$ times the tame character and $\varepsilon$ is trivial on $I$; the kernel of $I\to\langle t\rangle$ has trivial pro-$\ell$ quotient, so by inflation-restriction $H^1(I,M')=H^1(\langle t\rangle,M')=M'/(t-1)M'$ for every $\Z_\ell$-module $M'$ of finite type with this action. In the basis $e=e_1\otimes e_2^\vee$, $h=e_1\otimes e_1^\vee-e_2\otimes e_2^\vee$, $f=e_2\otimes e_1^\vee$ of $M=\ad T$ one has $(t-1)e=0$, $(t-1)h=-2n'e$, $(t-1)f=n'h-n'^2e$, and as $2$ is a unit, $(t-1)M=n'(\Z_\ell e\oplus\Z_\ell h)$; so $H^1(I,M)=(\Z_\ell/n')\bar e\oplus(\Z_\ell/n')\bar h\oplus\Z_\ell\bar f$ and its torsion is $(\Z/\ell^{v_\ell(n)})^2$. The group $G/I$ acts on $H^1(I,M)$ by $(\sigma\cdot z)(i)=\sigma z(\sigma^{-1}i\sigma)$ for a cocycle $z$; for the arithmetic Frobenius $\Frob$, $\Frob^{-1}t\Frob=t^{q^{-1}}$ and $z(t^k)\equiv k\,z(t)$ modulo $(t-1)M$, so $\Frob$ acts on $M/(t-1)M$ as $q^{-1}\operatorname{Ad}(\rho(\Frob))$, that is, by conjugation by $\rho(\Frob)$ followed by multiplication by $q^{-1}$. Writing $\rho(\Frob)=\begin{psmallmatrix}\varepsilon q&c'\\0&\varepsilon\end{psmallmatrix}$ with $\varepsilon=\varepsilon(\Frob)$, the adjoint action sends $e$ to $qe$ and $h$ to $h-2c'\varepsilon e$, so on the torsion $\bar e\mapsto\bar e$ and $\bar h\mapsto q^{-1}(\bar h-2c'\varepsilon\bar e)$. The fixed points are the $x\bar e+y\bar h$ with $(q-1)y\equiv0$ and $c'y\equiv0\pmod{\ell^{v_\ell(n)}}$: $\ell^{v_\ell(n)}$ choices of $x$ and $\ell^{\min(v_\ell(n),v_\ell(q-1),v_\ell(c'))}$ of $y$. Finally, $q$ being a uniformiser of $\Q_q$, $H^1(\Q_q,\Z_\ell(1))$ splits as $\Z_\ell\kappa_q\oplus H^1_{\mathrm{ur}}$, and $\kappa=n\kappa_q+\kappa_u$, where $\kappa_u$ is unramified, inflated from the class of $u$ in $H^1(G/I,\Z_\ell(1))=\Z_\ell(1)/(q-1)\simeq\F_q^\times\otimes\Z_\ell$. The value $c'=\varepsilon\kappa(\Frob)$ is a cocycle value: it is $\varepsilon$ times the class of $u$ for the Frobenius lift fixing a compatible system of $\ell$-power roots of $q$, and in general it is well defined only modulo $\ell^{\min(v_\ell(n),v_\ell(q-1))}$ (changing the lift changes it by a multiple of $n'$, changing the cocycle by a multiple of $q-1$). Neither indeterminacy affects $\min(v_\ell(n),v_\ell(q-1),v_\ell(c'))$, and $v_\ell(c')$ may be replaced by $v_\ell(c(u))$ there. The identification of $\Tam_q$ with the Fitting ideal of the $G$-invariants of the torsion of $H^1(I,M)$ is \cite[I.4.2.2]{FPR94}, as recalled in \cite[p.~708]{DFG04}. As a check, the same computation for $M=T$ gives $(\Z/\ell^{v_\ell(n)})\bar e_1$ with $\Frob$ acting by $\varepsilon$, that is, the $\ell$-part of the Tamagawa number $c_q(E)$, as it should. At an additive potentially multiplicative prime, $E$ is over $\Q_q$ the twist of the Tate curve $E_{q_E}$ by a ramified quadratic character \cite[V.5]{SilvermanATAEC}; the $G_{\Q_q}$-module $\ad T$ is unchanged by the twist, and the computation above, which only uses that module, applies with $n=v_q(q_E)=-v_q(j_E)$. At $q=2$, $v_\ell(q-1)=0$ and the minimum vanishes.
\end{proof}

\begin{lemma}[the Tamagawa factor at a potentially good additive prime]\label{lem:tamadd}
Let $\ell$ be odd and let $q\ne\ell$ be a prime at which $E$ has additive, potentially good reduction. Then $H^1(I_q,\ad T)_{\mathrm{tor}}=0$; in particular $v_\ell(\Tam_q(\ad T))=0$.
\end{lemma}

\begin{proof}
Let $M=\ad T$ and let $K$ be the kernel of the action of $I=I_q$ on $M$; the quotient $\Delta=I/K$ is the image of the finite group $\rho(I)$ modulo $\rho(I)\cap\{\pm1\}$. Inflation-restriction gives the exact sequence $0\to H^1(\Delta,M)\to H^1(I,M)\to\operatorname{Hom}(K,M)^\Delta$. The group $K$ is the inertia group of a finite extension of $\Q_q^{\mathrm{ur}}$, and a continuous homomorphism from $K$ to the pro-$\ell$ group $M$ factors through its maximal pro-$\ell$ abelian quotient, which is $\Z_\ell(1)$ as $q\ne\ell$; so $\operatorname{Hom}(K,M)\simeq M(-1)$ is torsion-free, and the torsion of $H^1(I,M)$ lies in $H^1(\Delta,M)$, which is killed by $\#\Delta$. The elements of $\rho(I)$ have orders $1,2,3,4,6$ (proof of Proposition~\ref{prop:shapes}), so $\#\Delta$ is divisible only by $2$ and $3$, and $H^1(\Delta,M)=0$ for $\ell\ge5$. Let $\ell=3$. The $3$-primary part of $H^1(\Delta,M)$ injects into $H^1(P,M)$ for a Sylow $3$-subgroup $P$ of $\Delta$, by restriction and corestriction. A finite $3$-subgroup of $\SL_2$ in characteristic $0$ is cyclic, and it has exponent $3$, so $P$ is trivial or generated by the image of an element $\sigma\in\rho(I)$ of order $3$. Then $T$ is free of rank one over $\Z_3[X]/(X^2+X+1)$, with $X$ acting as $\sigma$, so in a suitable basis $\sigma=\begin{psmallmatrix}0&-1\\1&-1\end{psmallmatrix}$; with $e=\begin{psmallmatrix}0&1\\0&0\end{psmallmatrix}$, the elements $e$, $\sigma e\sigma^{-1}$, $\sigma^2e\sigma^{-2}$ have coordinates $(1,0,0)$, $(0,0,-1)$, $(1,-1,-1)$ in the basis $e,h,f$ of $\ad T$, a matrix of determinant $-1$. So $\ad T\simeq\Z_3[P]$ is free over $\Z_3[P]$ and $H^1(P,\ad T)=0$. Hence $H^1(I,M)_{\mathrm{tor}}=0$, and so is its group of $G_{\Q_q}$-invariants.
\end{proof}

By contrast $\Tam_q(T)$ can be divisible by $3$ at a prime of type $\mathrm{IV}$ or $\mathrm{IV}^*$: then $H^1(\langle\sigma\rangle,T)=T/(1-\zeta_3)T\simeq\F_3$.

\section{Proof of Theorem~\ref{thm:qual}}\label{sec:proof}

Throughout this section $\ell$ is odd and $\ell\nmid N$. Let $f=f_E\in S_2(\Gamma_0(N))$ be the newform of the isogeny class of $E$, let $E_0$ be the optimal curve in the class, and let $r=r_{E_0}$ be its congruence number: the largest integer $r$ such that $f\equiv g\pmod r$ coefficientwise for some $g\in S_2(\Gamma_0(N))$ with integer coefficients orthogonal to $f$ \cite[\S2]{ARS12}. By \cite[Theorem 2.2]{ARS12}, $\deg\varphi_{E_0}$ divides $r$ and $v_p(\deg\varphi_{E_0})=v_p(r)$ for every prime $p$ with $p^2\nmid N$; as $\ell\nmid N$ it suffices to prove $\ell\mid r$. Let $\T$ be the Hecke algebra of level $N$, the subring of $\operatorname{End}S_2(\Gamma_0(N))$ generated over $\Z$ by all $T_n$, $n\ge1$, including the operators $U_p$ for $p\mid N$, and write $\bar\rho=\bar\rho_{E,\ell}$, irreducible by hypothesis, hence absolutely irreducible, complex conjugation having the distinct eigenvalues $1$ and $-1$. The two lemmas below reduce the theorem to the production of a congruent newform of lower level.

\begin{lemma}[congruence criterion]\label{lem:crit}
Let $\phi_f\colon\T\to\Z$, $T_n\mapsto a_n(f)$, let $I=\ker\phi_f$ and let $J$ be the annihilator in $\T$ of the orthogonal complement $f^\perp$ of $f$. Then $\T/(I+J)$ is cyclic of order $r$; equivalently, $r$ generates the congruence ideal $\phi_f(J)=\phi_f(\operatorname{Ann}_\T I)$. Consequently, if there is a normalised eigenform $g\in S_2(\Gamma_0(N))\otimes\overline\Q$ for all of $\T$, orthogonal to $f$, with coefficients in the ring of integers $\mathcal{O}$ of a number field, and a prime $\lambda\mid\ell$ of $\mathcal{O}$ such that $a_n(g)\equiv a_n(f)\pmod\lambda$ for every $n\ge1$, then $\ell\mid r$. The same statements hold with $\Z$ replaced by $\Z_\ell$ and $r$ by its $\ell$-part.
\end{lemma}

\begin{proof}
Writing $\T\otimes\Q=\Q\times\T'$ with $\T'$ acting faithfully on $f^\perp$ (the $f$-line is $\T$-stable, and so is $f^\perp$, because $f$ is an eigenvector of the adjoints of the $T_n$ as well), one has $I=\T\cap(0\times\T')$ and $J=\T\cap(\Q\times0)$, so $I\cap J=0$ and $\T/(I+J)\simeq\Z/\phi_f(J)$ is cyclic, of order $r'$ say. The pairing $S_2(\Gamma_0(N);\Z)\times\T\to\Z$, $(g,t)\mapsto a_1(tg)$, is perfect \cite[Theorem 2.2]{Ribet83}, \cite[Lemma 2.1]{AU96} (see \cite[\S3]{ARS12}), and under it the forms killed by $I$ are $\Z f$ and the forms killed by $J$ are $L=f^\perp\cap S_2(\Gamma_0(N);\Z)$. Dualising the inclusion $\T\subset\T/I\times\T/J$, whose cokernel is $\T/(I+J)$, shows that $\Z f\oplus L$ has index $r'$ in $S_2(\Gamma_0(N);\Z)$. The quotient $S_2(\Gamma_0(N);\Z)/(\Z f\oplus L)$ embeds in $\Q f/\Z f$ by projection onto $\Q f$ along $f^\perp$, so it is cyclic, generated by the class of a form $h$ with projection $f/r'$; then $r'h=f+g_0$ with $g_0\in L$, so $f\equiv-g_0\pmod{r'}$ and $r'\le r$; conversely a congruence $f\equiv g_1\pmod r$ with $g_1\in L$ gives the integral form $(f-g_1)/r$, whose class has order $r$, so $r\mid r'$. Hence $r=r'=\#\T/(I+J)$. Now let $\phi_g\colon\T\to\mathcal{O}$ be the eigencharacter of $g$. It kills $J$ because $g\in f^\perp$, and $\phi_g\equiv\phi_f\pmod\lambda$ on all of $\T$ by hypothesis, so $\phi_g(I)\subset\lambda$. Thus $I+J$ lies in the kernel of $\T\to\mathcal{O}/\lambda$, a maximal ideal of residue characteristic $\ell$, and $\ell$ divides $\#\T/(I+J)=r$. Every step holds verbatim over $\Z_\ell$, the pairing being perfect after tensoring with $\Z_\ell$.
\end{proof}

\begin{lemma}[oldforms]\label{lem:old}
Let $h$ be a newform of weight $2$ on $\Gamma_0(M)$ with $M\mid N$ and $h\ne f$, with coefficients in $\mathcal{O}$, and let $\lambda\mid\ell$ be a prime of $\mathcal{O}$ such that $a_p(h)\equiv a_p(f)\pmod\lambda$ for all primes $p\nmid\ell N$. Assume that $p\nmid M$ for every prime $p\,\|\,N$ at which $\bar\rho$ is unramified. Then, after enlarging $\mathcal{O}$ and $\lambda$, the span of the forms $h(dz)$, $d\mid N/M$, contains an eigenform $g$ as in Lemma~\ref{lem:crit}. Consequently $\ell\mid r$.
\end{lemma}

\begin{proof}
The forms $h(dz)$, $d\mid N/M$, span a $\T$-stable subspace $V$ of $S_2(\Gamma_0(N))$ orthogonal to the newform $f$ (whether $M<N$, or $M=N$ and $h\ne f$). On $V$ the operators $T_n$ with $(n,N)=1$ act by the scalars $a_n(h)$, and $U_p$ for $p\mid M$ with $v_p(M)=v_p(N)$ acts by $a_p(h)$. For $p\mid N/M$ put $e_p=v_p(N/M)$; on the span of $h(p^jz)$, $0\le j\le e_p$, one has $U_ph(p^jz)=h(p^{j-1}z)$ for $j\ge1$, and $U_ph=a_p(h)h-p\,h(pz)$ if $p\nmid M$, $U_ph=a_p(h)h$ if $p\mid M$. The eigenvalues of $U_p$ there are therefore the roots $\alpha_p,\beta_p$ of $X^2-a_p(h)X+p$, together with $0$ if $e_p\ge2$, when $p\nmid M$; and $a_p(h)$, together with $0$ if $e_p\ge1$, when $p\mid M$. The operator $U_p$ is not semisimple on this span when $0$ is a repeated eigenvalue, but each listed eigenvalue has an eigenvector: the $p$-stabilisations $h-\beta_ph(pz)$, $h-\alpha_ph(pz)$ for $\alpha_p,\beta_p$, the form $h$ for $a_p(h)$, and the depleted forms $h-a_p(h)h(pz)+p\,h(p^2z)$ (if $p\nmid M$), $h-a_p(h)h(pz)$ (if $p\,\|\,M$) or $h$ itself (if $p^2\mid M$, where $a_p(h)=0$) for the eigenvalue $0$. Since $V$ is the tensor product of these local spaces and $\T$ is commutative, for every choice of one such eigenvalue $\mu_p$ at each $p\mid N/M$ there is a common eigenform $g\in V\otimes\overline\Q$ of all of $\T$ with $U_pg=\mu_pg$ for $p\mid N/M$, $U_pg=a_p(h)g$ for the other $p\mid N$, and $T_ng=a_n(h)g$ for $(n,N)=1$; normalise it by $a_1(g)=1$. We choose the $\mu_p$ and check $a_p(g)\equiv a_p(f)\pmod\lambda$ prime by prime; the congruence for all $n$ then follows from the multiplicativity of the coefficients, from $a_{p^k}=a_p^k$ for $p\mid N$, and from the recursion $a_{p^{k+1}}=a_pa_{p^k}-p\,a_{p^{k-1}}$ for $p\nmid N$, which $f$ and $g$ both satisfy.

\emph{$p\nmid\ell N$.} $a_p(g)=a_p(h)\equiv a_p(f)$ by hypothesis. Consequently $\bar\rho_h\simeq\bar\rho$ over a common finite extension of $\F_\ell$, by Chebotarev and Brauer-Nesbitt, $\bar\rho$ being irreducible, hence absolutely irreducible.

\emph{$p=\ell$.} Both $f$ and $h$ have weight $2$ and level prime to $\ell$. By the theorems of Deligne and Fontaine describing $\bar\rho_f|_{D_\ell}$ and $\bar\rho_h|_{D_\ell}$ \cite[\S2]{Edixhoven92}, the residue class of $a_\ell$ is determined by the restriction to $D_\ell$: it is $0$ when that restriction is irreducible, and otherwise the eigenvalue of Frobenius on the unramified quotient. That quotient is well defined: the two characters of the semisimplification restrict to $I_\ell$ as $\bar\chi_\ell$ and $1$, and $\bar\chi_\ell|_{I_\ell}\ne1$ since $\ell$ is odd, so the unramified quotient is the module of $I_\ell$-coinvariants. Hence $a_\ell(h)\equiv a_\ell(f)$.

\emph{$p^2\mid N$.} Here $E$ has additive reduction at $p$ and $a_p(f)=0$. If $v_p(M)=v_p(N)$ then $a_p(h)=0$, because a newform on $\Gamma_0(M)$ with $p^2\mid M$ has $a_p=0$ \cite[Theorem 3]{AtkinLehner70}. Otherwise $e_p\ge1$, and $e_p\ge2$ if $p\nmid M$, so $0$ is an available eigenvalue: take $\mu_p=0$.

\emph{$p\,\|\,N$, $\bar\rho$ unramified at $p$.} By hypothesis $p\nmid M$, so $e_p=1$ and $\mu_p\in\{\alpha_p,\beta_p\}$. Here $\rho_{E,\ell}|_{D_p}$ is an extension of an unramified character $\varepsilon$ by $\varepsilon\chi_\ell$ with $\varepsilon(\Frob_p)=a_p(f)=\pm1$ \cite[Theorem 3.1(e)]{DDT}, so the eigenvalues of $\bar\rho(\Frob_p)$ are $a_p(f)p$ and $a_p(f)$. As $\bar\rho_h\simeq\bar\rho$ is unramified at $p$ and $p\nmid M$, the characteristic polynomial of $\bar\rho_h(\Frob_p)$ is $X^2-a_p(h)X+p\equiv(X-a_p(f)p)(X-a_p(f))\pmod\lambda$, so one of $\alpha_p,\beta_p$ is congruent to $a_p(f)$ modulo a prime above $\lambda$ in $\mathcal{O}[\alpha_p]$; take it as $\mu_p$.

\emph{$p\,\|\,N$, $\bar\rho$ ramified at $p$.} Then $p\mid M$, since otherwise $\bar\rho_h$ would be unramified at $p$; so $v_p(M)=1=v_p(N)$ and $a_p(g)=a_p(h)=\pm1$. Both $\bar\rho|_{D_p}$ and $\bar\rho_h|_{D_p}$ are extensions of an unramified character by its twist by $\bar\chi_\ell$, with Frobenius eigenvalue $a_p(f)$, respectively $a_p(h)$, on the unramified quotient \cite[Theorem 3.1(e)]{DDT}, and both are non-split because $\bar\rho$ is ramified at $p$. A non-split extension of two characters has a unique stable line, so the isomorphism $\bar\rho|_{D_p}\simeq\bar\rho_h|_{D_p}$ matches the sublines and hence the quotients, and $a_p(h)\equiv a_p(f)$.

In all cases $a_p(g)\equiv a_p(f)\pmod\lambda$, and Lemma~\ref{lem:crit} applies.
\end{proof}

We use the optimal-level theorem of Ribet, Carayol and Diamond \cite{Ribet90,Carayol89,Diamond95} in the form of \cite[Theorem 3.15]{DDT}: \emph{let $\ell$ be odd and let $\bar\sigma\colon\Gal(\overline\Q/\Q)\to\GL_2(\overline{\F}_\ell)$ be irreducible and modular, absolutely irreducible on $\Gal(\overline\Q/\Q(\sqrt{-3}))$ if $\ell=3$; then $\bar\sigma$ arises from a newform of weight $2$ and level $N(\bar\sigma)\ell^{\delta(\bar\sigma)}$ with nebentypus of order prime to $\ell$, where $N(\bar\sigma)$ is the Serre conductor and $\delta(\bar\sigma)=0$ when $\bar\sigma|_{D_\ell}$ is good.} For $\bar\sigma=\bar\rho_{E',\ell}$ with $E'$ an elliptic curve of conductor prime to $\ell$, the restriction to $D_\ell$ comes from the finite flat group scheme $E'[\ell]$ and is good, and $\det\bar\sigma=\bar\chi_\ell$ forces the nebentypus, which is trivial modulo $\lambda$ and of order prime to $\ell$, to be trivial. So in that case $\bar\sigma$ arises from a newform of weight $2$ on $\Gamma_0(N(\bar\sigma))$. Recall also that $N(\bar\rho_{E',\ell})$ divides the conductor of $E'$ \cite{Carayol89}, and that $N(\bar\sigma)$ is prime to every prime at which $\bar\sigma$ is unramified.

\medskip\noindent\emph{Case (a).} Here $q^2\,\|\,N$ and $\ell=3$ (Lemma~\ref{lem:IV}, which gives conductor exponent $2$ at $q$, also for $q=2$). By Lemma~\ref{lem:IV} the Serre conductor $M=N(\bar\rho)$ has exponent exactly $1$ at $q$, so $M\mid N/q$. The optimal-level theorem gives a newform $g$ of weight $2$ on $\Gamma_0(M)$ with $\bar\rho_g\simeq\bar\rho$, that is, $a_p(g)\equiv a_p(f)\pmod\lambda$ for all $p\nmid 3N$. Take $h=g$: it has level $M<N$, and $p\nmid M$ for every prime $p$ at which $\bar\rho$ is unramified. Lemma~\ref{lem:old} gives $3\mid r$; at $p=q$ it uses the eigenvalue $0$ of $U_q$, available since $v_q(M)=1<2=v_q(N)$.

\medskip\noindent\emph{Case (b).} Here $q^2\,\|\,N$, $\ell=3$, and $E$ has type $\mathrm{II}$ or $\mathrm{II}^*$ at $q\ge5$. Let $\chi=\chi_{q^*}$ and $E'=E^{(q)}$. By Lemma~\ref{lem:twist}, $E'$ has type $\mathrm{IV}^*$ or $\mathrm{IV}$ at $q$, the same types as $E$ elsewhere, and conductor $N$; its newform is $f'=f\otimes\chi$, a newform of weight $2$ on $\Gamma_0(N)$ (both have $a_q=0$), and $\bar\rho_{E',3}=\bar\rho\otimes\chi$ satisfies the hypotheses of the optimal-level theorem since $\bar\rho$ does. By Lemma~\ref{lem:IV} applied to $E'$, the Serre conductor $M=N(\bar\rho\otimes\chi)$ has exponent exactly $1$ at $q$, and $M\mid N/q$. The optimal-level theorem gives a newform $g$ of weight $2$ on $\Gamma_0(M)$ with $a_p(g)\equiv a_p(f')=\chi(p)a_p(f)\pmod\lambda$ for all $p\nmid 3N$. As $q\,\|\,M$, Lemma~\ref{lem:twist}(iii) says that $h=g\otimes\chi$ is a newform on $\Gamma_0(Mq)$, and $Mq\mid N$. For $p\nmid 3N$,
\[
a_p(h)=\chi(p)a_p(g)\equiv\chi(p)^2a_p(f)=a_p(f)\pmod\lambda .
\]
Moreover $h\ne f$: otherwise $a_p(g)=\chi(p)a_p(h)=a_p(f')$ for all $p\nmid qN$, so $g=f'$ by strong multiplicity one, whereas $g$ has level $M<N$. (Note that $h\otimes\chi$ is the $q$-depletion of $g$, not $g$, as $a_q(g)=\pm1$.) Finally, if $p\,\|\,N$ and $\bar\rho$ is unramified at $p$, then $p\ne q$, so $\bar\rho\otimes\chi$ is unramified at $p$ as well and $p\nmid M$, hence $p\nmid Mq$. Lemma~\ref{lem:old} gives $3\mid r$.

\medskip\noindent\emph{Case (c).} Here $E$ has type $\mathrm{I}_n^*$ at an odd prime $q\ne\ell$ with $\ell\mid n$, so $q^2\,\|\,N$. Let $\chi=\chi_{q^*}$ and $E'=E^{(q)}$. By Lemma~\ref{lem:twist}, $E'$ has multiplicative reduction of type $\mathrm{I}_n$ at $q$ and conductor $N/q$; let $f'=f_{E'}$ be its newform, of weight $2$ on $\Gamma_0(N/q)$ with $q\,\|\,N/q$, $a_q(f')=\pm1$ and $a_p(f')=\chi(p)a_p(f)$ for $p\ne q$ (so $f\otimes\chi$ is the $q$-depletion of $f'$; for example 539c1 twisted by $\chi_{-7}$ lies in the class 77c, and $a_7(\text{77c1})=-1$). By Lemma~\ref{lem:tamagawa-exp} applied to $E'$, $\bar\rho\otimes\chi=\bar\rho_{E',\ell}$ is unramified at $q$, so the Serre conductor $M=N(\bar\rho\otimes\chi)$ is prime to $q$ and $M\mid N/q^2$. The optimal-level theorem (this is the case of Ribet's theorem \cite{Ribet90}, $q$ exactly dividing the level and $\bar\rho\otimes\chi$ unramified at $q$) gives a newform $g$ of weight $2$ on $\Gamma_0(M)$ with $a_p(g)\equiv\chi(p)a_p(f)\pmod\lambda$ for all $p\nmid\ell N$. As $q\nmid M$, $h=g\otimes\chi$ is a newform on $\Gamma_0(Mq^2)$ with $Mq^2\mid N$, congruent to $f$ at all $p\nmid\ell N$ by the computation of case (b), and $h\ne f$: otherwise $g$ and $f'$ would have the same $a_p$ for all $p\nmid qN$, so $g=f'$ by strong multiplicity one, whereas their levels $M$ and $N/q$ differ. The hypothesis of Lemma~\ref{lem:old} holds as in case (b), and $\ell\mid r$. \qed

\begin{example}\label{ex:qual}
The congruences produced by the proof can be seen directly on the tables (labels as in \cite{Cremona97}); in each case the two forms lie in the same space $S_2(\Gamma_0(N))$ and have integer coefficients, so a check up to Sturm's bound \cite{Sturm87}, $\frac16N\prod_{p\mid N}(1+\frac1p)$, proves the congruence for all $n$. (a) at $q=2$: 52a1 has type $\mathrm{IV}^*$ at $2$; its mod $3$ representation has Serre conductor $26$, and $f_{\text{52a}}\equiv f_{\text{26b}}(z)-f_{\text{26b}}(2z)\pmod 3$ for all $n$ (Sturm's bound $14$; $a_2(\text{26b1})=1$), while $\deg\varphi_{\text{52a1}}=3$: here $e_3(E,2)=1$ and the bound is attained. (b): the curve 121a1 has type $\mathrm{II}$ at $11$; its twist by $\chi_{-11}$ lies in the class 121c and has type $\mathrm{IV}^*$ at $11$; its mod $3$ representation has conductor $11$ and is congruent to the form of 11a1; twisting back, $f_{\text{11a}}\otimes\chi_{-11}$ is the form of 121d, and indeed $f_{\text{121a}}\equiv f_{\text{121d}}\pmod 3$ for all $n$ (Sturm's bound $22$), while $\deg\varphi_{\text{121a1}}=6$: here $e_3(E,11)=v_3(11+1)=1$ and the bound is attained. (c) at $\ell=3$: 539c1 has type $\mathrm{I}_3^*$ at $7$; its twist by $\chi_{-7}$ has conductor $77$ and type $\mathrm{I}_3$ at $7$, and $f_{\text{539c}}\equiv f_{\text{11a}}\otimes\chi_{-7}=f_{\text{539d}}\pmod3$ for all $n$ (Sturm's bound $112$); $\deg\varphi_{\text{539c1}}=2^5\cdot3^2$, and the two terms $e_3(E,7)=v_3(48)=1$ and $t_3(E,7)=v_3(3)=1$ are both attained. (c) at $\ell=5$: 980i1 has type $\mathrm{I}_5^*$ at $7$; its twist by $\chi_{-7}$ has conductor $140$ and type $\mathrm{I}_5$ at $7$, and $f_{\text{980i}}\equiv f_{\text{20a}}\otimes\chi_{-7}=f_{\text{980g}}\pmod5$ for all $n$ (Sturm's bound $336$); $\deg\varphi_{\text{980i1}}=2^6\cdot3^2\cdot5$. As $5\mid980$, this last example illustrates the congruence, not the theorem.
\end{example}

\begin{remark}
The hypothesis $\ell\nmid N$ is used in the optimal-level theorem, in Lemma~\ref{lem:old} at $p=\ell$, and to pass from the congruence number to the modular degree. The last step needs only $\ell^2\nmid N$, and it cannot be dropped there: Table 1 of \cite{ARS12} lists the curve 99A1, of conductor $9\cdot11$, with modular degree $4$ and congruence number $12$; likewise 726a1 has $\degphi=144$ and $r_E=1584=11\cdot144$, with $11^2\mid726$, and 539d1 has $\degphi=72$ and $r_E=504=7\cdot72$, with $7^2\mid539$. The data nevertheless show the conclusion of the theorem for every curve with $\ell\,\|\,N$ or $\ell^2\mid N$ and a prime of one of the three kinds; so do they at $q=3$ in cases (a) and (b), where $q=\ell$ and the theorem does not apply.
\end{remark}

\section{The congruence ideal and the adjoint Selmer group}\label{sec:quant}

\subsection{Congruence ideals}\label{sec:dfg}
Let $f$ be a newform of weight $2$ on $\Gamma_0(M)$ with rational coefficients, $\lambda=\ell$ a prime not dividing $2M$, and $\Sigma$ a finite set of primes. Put $M^\Sigma=M\prod_{p\in\Sigma}p^{\delta_p}$ and let $f^\Sigma=\sum_{(n,\Sigma)=1}a_n(f)q^n$ be the $\Sigma$-depleted eigenform, of level $M^\Sigma$. Diamond, Flach and Guo define \cite[\S\S1.5.3, 1.6.2, 1.7.3]{DFG04} a lattice $M_{f^\Sigma,\lambda}$ and an ideal $\eta^\Sigma_{f,\lambda}\subset\Z_\ell$ as follows. Their integral structure $M(M^\Sigma,1)_{!,\lambda}$ is the image of the $U_0(M^\Sigma)$-invariants of the compactly supported cohomology of a modular curve of full level in its cohomology, with its Hecke action \cite[\S\S1.2, 1.4]{DFG04}. The following comparison identifies it with the cohomology of $X_0(M^\Sigma)$.

\begin{proposition}[comparison of integral structures]\label{prop:integral}
Let $\ell$ be a prime. Pull-back along the covering of $X_0(M^\Sigma)$ by the modular curve of full level identifies $H^1(X_0(M^\Sigma),\Z)$ with the Betti lattice $M(M^\Sigma,1)_{!,B}$ of \cite[\S1.4.2]{DFG04}, compatibly with the Hecke operators, $U_p$ included; in particular $M(M^\Sigma,1)_{!,\lambda}=H^1(X_0(M^\Sigma),\Z_\ell)$, with no localisation.
\end{proposition}
\begin{proof}
Write $M=M^\Sigma$ and $U=U_0(M)$. For weight $2$ and trivial character the coefficient sheaf of \cite[\S1.2.2]{DFG04} is $\Z$. Choose $N\ge3$ with $M\mid N$ and every prime factor of $N$ in the set $S$ of \cite[\S1.4.1]{DFG04} (for instance $N=4M$). Then $U_N\subseteq U$, and by \cite[\S1.4.1]{DFG04} $M(M,1)_{?,B}=H^1_?(Y(N),\Z)^{G}$ for $?\in\{c,\mathrm{tf}\}$, where $Y(N)$ is the analytic modular curve of full level $N$, a disjoint union of copies of $\Gamma(N)\backslash\mathcal H$ indexed by $(\Z/N\Z)^\times$, and $G=U/U_N$ acts through its image in $\GL_2(\Z/N\Z)$ \cite[(5), Proposition 1.3]{DFG04}; $M(M,1)_{!,B}$ is the image of $M(M,1)_{c,B}\to M(M,1)_{\mathrm{tf},B}$, and none of this depends on $N$ \cite[Lemma 1.1]{DFG04}. Since $\det U=\hat\Z^\times$, $G$ permutes the components transitively, the stabiliser of a component acts through the image of $\Gamma_0(M)$, and the quotient map $\pi\colon Y(N)\to Y(N)/G=Y_0(M)$ is finite onto the coarse curve; the stabiliser in $G$ of a point of $Y(N)$ is finite.

\emph{Step 1: $\pi^*\colon H^1_c(Y_0(M),\Z)\to H^1_c(Y(N),\Z)^G$ is an isomorphism.} Let $H^\bullet_{G,c}$ denote $G$-equivariant cohomology with compact support. In the Hochschild-Serre spectral sequence $H^p(G,H^q_c(Y(N),\Z))\Rightarrow H^{p+q}_{G,c}(Y(N),\Z)$ the row $q=0$ vanishes, because every component of $Y(N)$ is a non-compact connected surface, so $H^0_c(Y(N),\Z)=0$; hence $H^1_{G,c}(Y(N),\Z)=H^1_c(Y(N),\Z)^G$. In the equivariant Leray spectral sequence of the finite map $\pi$, $H^p_c(Y_0(M),\mathcal R^q)\Rightarrow H^{p+q}_{G,c}(Y(N),\Z)$, the stalk of $\mathcal R^q$ at $y$ is $H^q(G_x,\Z)$ for a point $x$ above $y$ (Shapiro's lemma), so $\mathcal R^0=\Z$ and $\mathcal R^1=0$, since $H^1(C,\Z)=\operatorname{Hom}(C,\Z)=0$ for a finite group $C$. Hence $H^1_{G,c}(Y(N),\Z)=H^1_c(Y_0(M),\Z)$. The composite of the two edge isomorphisms is $\pi^*$.

\emph{Step 2: the image.} Under these identifications $M(M,1)_{c,B}\to M(M,1)_{\mathrm{tf},B}$ is the forgetful map $H^1_c(Y(N),\Z)^G\to H^1(Y(N),\Z)^G$, and as $\pi$ is proper its image is $\pi^*$ of the image of $H^1_c(Y_0(M),\Z)\to H^1(Y_0(M),\Z)$. Since $Y_0(M)$ is the compact surface $X_0(M)$ with finitely many points removed, $H^1_c(Y_0(M),\Z)\to H^1(X_0(M),\Z)$ is surjective and $H^1(X_0(M),\Z)\to H^1(Y_0(M),\Z)$ is injective, so that image is $H^1(X_0(M),\Z)$. Finally $\pi^*$ is injective on the torsion-free group $H^1(Y_0(M),\Z)$, since $\pi_*\pi^*$ is multiplication by $\deg\pi$. So $\pi^*$ maps $H^1(X_0(M),\Z)$ isomorphically onto $M(M,1)_{!,B}$.

\emph{Step 3: Hecke operators.} Both lattices are torsion-free and $\pi^*$ is injective, so equivariance may be checked after tensoring with $\C$. There the Hecke operators of \cite{DFG04} act on $\operatorname{Fil}^1M(M,1)_{!,\mathrm{dR}}\otimes\C=S_2(\Gamma_0(M))$ as the classical $T_p$, and as $U_p$ for $p\mid M$ \cite[Proposition 1.3]{DFG04}, and commute with complex conjugation $F_\infty$; $\pi^*$ is compatible with the comparison isomorphism (5) of \cite{DFG04} and with the Eichler-Shimura isomorphism, both being pull-back of $f(\tau)\,d\tau$; and the Hecke action on $H^1(X_0(M),\Z)$ is the one for which the Eichler-Shimura isomorphism is Hecke equivariant (Section~\ref{sec:quant}). Hence $\pi^*$ is Hecke equivariant, $U_p$ included.

\emph{Step 4: the $\lambda$-adic structure.} By \cite[\S1.2.3]{DFG04}, $\Z_\ell\otimes M_{?,B}\simeq M_{?,\ell}$ for every prime $\ell$, and taking invariants under a finite group and images commutes with the flat base change $\Z\to\Z_\ell$; so $M(M,1)_{!,\lambda}=\Z_\ell\otimes M(M,1)_{!,B}=H^1(X_0(M),\Z_\ell)$.
\end{proof}

\begin{remark}
Compact support is what makes the comparison integral: $H^0_c$ of the covering vanishes, so taking invariants under the covering group loses nothing, as in \cite[Lemma 1.1]{DFG04}. With ordinary cohomology it is not so: $H^1(Y(N),\Z_\ell)^G/\pi^*H^1(Y_0(M),\Z_\ell)$ is the group of characters $\Gamma_0(M)\to\Q_\ell/\Z_\ell$ of $\ell$-power order that factor through $G$, namely $\gamma\mapsto\psi(d\bmod p)$ for primes $p\mid M$ with $\ell\mid p-1$; for $M=11$ and $\ell=5$ it is the unramified covering $X_1(11)\to X_0(11)$.
\end{remark}

The lattice $M_{f^\Sigma,\lambda}$ is the kernel on $M(M^\Sigma,1)_{!,\lambda}$ of the kernel of the eigencharacter of $f^\Sigma$ on the full Hecke algebra, $U_p$ included \cite[\S\S1.1.2, 1.6.2]{DFG04}. The pairing $\hat\delta_!$ is the cup product composed with the operator $[UwU]_\omega$ attached to $w=\begin{psmallmatrix}0&-1\\M^\Sigma&0\end{psmallmatrix}$ \cite[\S1.5.3, (13), (14)]{DFG04}, which for weight $2$ and trivial character is the Atkin-Lehner involution $w_{M^\Sigma}$ with no scalar factor; on classes pulled back from $X_0(M^\Sigma)$ it is $\frac u2$ times the cup product on $X_0(M^\Sigma)$ composed with $w_{M^\Sigma}$, for an $\ell$-adic unit $u$, so for odd $\ell$ only a unit separates the two. Indeed, for a sufficiently small level $U'\subseteq U_0(M^\Sigma)$ with $\det U'=\hat\Z^\times$, such as $U_0(M^\Sigma)\cap U_1(4)$, the normalised pairing $\bar\delta$ of \cite[\S1.5.2]{DFG04} is Poincar\'e duality on $X_{U'}$ up to an $\ell$-adic unit when $\ell\nmid 2M^\Sigma$ \cite[\S1.5.2, before Lemma 1.2]{DFG04}; it is normalised by dividing by the index, so on $U_0(M^\Sigma)$-invariants it is $[U_0(M^\Sigma):U']^{-1}$ times its value on $U'$-invariants, and the covering $X_{U'}\to X_0(M^\Sigma)$ has degree $[U_0(M^\Sigma):U']/2$, because $-1$ acts trivially. It is perfect and alternating on $H^1(X_0(M^\Sigma),\Z_\ell)$, and the Hecke operators, $U_p$ included, are self-adjoint for it, because $w_{M^\Sigma}T_nw_{M^\Sigma}^{-1}$ is the adjoint of $T_n$ for the cup product (this is what the composition with $w$ is for; see \cite[\S\S1.6.1, 1.8.2]{DFG04} and \cite[\S4.4]{DDT}). The restriction of $\hat\delta_!$ to the rank-two lattice $M_{f^\Sigma,\lambda}$ takes its values in an ideal of $\Z_\ell$, the image of $\wedge^2M_{f^\Sigma,\lambda}$, which is $\eta^\Sigma_{f,\lambda}$, the \emph{naive $\Sigma$-finite congruence ideal} \cite[\S1.7.3]{DFG04}; equivalently the Gram determinant of the pairing on $M_{f^\Sigma,\lambda}$ is $(\eta^\Sigma_{f,\lambda})^2$ up to a unit.

\begin{theorem}[{\cite[Proposition 1.4(c)]{DFG04}}]\label{thm:dfg14}
Let $\ell\notin\Sigma$. If $\bar\rho_f$ is irreducible then
\[
\eta^\Sigma_{f,\lambda}=\eta^{\emptyset}_{f,\lambda}\prod_{p\in\Sigma}L_p^{\nv}(A_f,1)^{-1}.
\]
\end{theorem}
Only the primes with $\delta_p\ge1$ contribute: adding an additive prime to $\Sigma$ changes neither the level nor the ideal; adding a good prime $p$ multiplies the ideal by $(1-\frac\alpha\beta p^{-1})(1-p^{-1})(1-\frac\beta\alpha p^{-1})$, Wiles's level-raising factor; adding a multiplicative prime multiplies it by $1-p^{-2}$. The hypothesis $\ell\notin\Sigma$ is a standing assumption of \cite[\S1.7.3]{DFG04}, and it cannot be dropped: for $f=f_{\text{11a}}$, $\ell=3$ and $\Sigma=\{3\}$ one has $a_3=-1$ and $L^{\nv}_3(A_f,1)^{-1}=\frac23\cdot\frac{16-1}9=\frac{10}9$, of $3$-adic valuation $-2$, so the right side would not be contained in $\Z_3$. Every use below has $\Sigma$ a set of primes dividing the level and $\ell$ prime to the level.

For the optimal curve, the congruence ideal at level $N$ is the modular degree.

\begin{proposition}\label{prop:pullback}
Let $E$ be $X_0(N)$-optimal with newform $f$, $\varphi=\varphi_E$, $I_f$ the kernel of the eigencharacter $\phi_f\colon\T\to\Z$ of $f$ on the full Hecke algebra, and $\ell$ a prime. Then $\varphi^*\colon H^1(E,\Z_\ell)\to H^1(X_0(N),\Z_\ell)$ is injective, with image the kernel lattice $H^1(X_0(N),\Z_\ell)[I_f]$, and
\[
\langle\varphi^*a,\,w_N^*\varphi^*b\rangle_{X_0(N)}=\varepsilon_N(f)\,\degphi\,\langle a,b\rangle_E\qquad(a,b\in H^1(E,\Z_\ell)),
\]
where $\langle\ ,\ \rangle$ denotes the cup product and $\varepsilon_N(f)=\pm1$ the eigenvalue of $w_N$ on $f$. In particular, for $\ell$ odd, the ideal of values of $\frac12\langle x,w_N^*y\rangle$ on $H^1(X_0(N),\Z_\ell)[I_f]$ is $\degphi\,\Z_\ell$, and its Gram determinant is $(\degphi)^2$ up to a unit. If $\ell\nmid2N$, then $M_{f,\lambda}=\varphi^*H^1(E,\Z_\ell)$, and
\[
v_\lambda(\eta^\emptyset_{f,\lambda})=v_\ell(\degphi).
\]
\end{proposition}

\begin{proof}
We use Betti cohomology of the Riemann surfaces, with the Hecke action for which the Eichler-Shimura isomorphism $H^1(X_0(N),\C)=S_2(\Gamma_0(N))\oplus\overline{S_2(\Gamma_0(N))}$ is Hecke equivariant. By optimality the kernel $B$ of the map $J_0(N)\to E$ induced by $\varphi$ is connected; the exact sequence $H_1(J_0(N),\Z)\to H_1(E,\Z)\to\pi_0(B)=0$ shows that $\varphi_*\colon H_1(X_0(N),\Z)=H_1(J_0(N),\Z)\to H_1(E,\Z)$ is surjective, so its dual $\varphi^*\colon H^1(E,\Z)\to H^1(X_0(N),\Z)$ is injective with torsion-free cokernel, and the same holds after tensoring with $\Z_\ell$. The modular parametrisation satisfies $\varphi^*\omega_E=c_E\cdot2\pi i\,f(z)\,dz$ with the Manin constant $c_E>0$, so $\varphi^*H^1(E,\C)$, which is spanned by $\varphi^*\omega_E$ and its conjugate, is $\C f\,dz\oplus\C\overline{f\,dz}$. Since $f$ is a newform, the eigenspace of $\phi_f$ in $S_2(\Gamma_0(N))$ is the line of $f$, so this plane is $H^1(X_0(N),\C)[I_f]$. It is defined over $\Q$: $H^1(X_0(N),\Q)[I_f]$ is a two-dimensional $\Q$-form of it, and by flatness $H^1(X_0(N),\Z_\ell)[I_f]=H^1(X_0(N),\Z)[I_f]\otimes\Z_\ell$ is a saturated lattice of rank two in $H^1(X_0(N),\Q_\ell)[I_f]$. The image of $\varphi^*$ is a saturated lattice of rank two in the same $\Q_\ell$-plane; so the two coincide. Next, $\varphi^*(a\cup b)=\varphi^*a\cup\varphi^*b$ and $\varphi^*$ multiplies the fundamental class of $H^2$ by $\degphi$; so $\langle\varphi^*a,\varphi^*b\rangle_{X_0(N)}=\degphi\,\langle a,b\rangle_E$, the projection formula. By multiplicity one for newforms the plane is also the eigenspace of the $T_n$, $(n,N)=1$, for the eigenvalues $a_n(f)$, and $w_N$ commutes with these $T_n$ (not with the $U_p$, $p\mid N$); so $w_N$ preserves the plane, on which it acts by $\varepsilon_N(f)$ ($f|w_N=\varepsilon_N(f)f$ and $\varepsilon_N(f)$ is real); this gives the displayed formula. The cup product on $H^1(E,\Z)$ is unimodular, so the ideal of values is $\degphi\,\Z_\ell$, and the Gram determinant of a rank-two alternating form is the square of its ideal of values. Finally, if $\ell\nmid2N$, the integral structure of \cite{DFG04} is $H^1(X_0(N),\Z_\ell)$ (Proposition~\ref{prop:integral}), so $M_{f,\lambda}=H^1(X_0(N),\Z_\ell)[I_f]=\varphi^*H^1(E,\Z_\ell)$, and $\hat\delta_!$ is $\frac u2\langle x,w_N^*y\rangle$ with $u$ an $\ell$-adic unit (Section~\ref{sec:dfg}); so $\eta^\emptyset_{f,\lambda}=\degphi\,\Z_\ell$.
\end{proof}

No multiplicity-one hypothesis, and for the first part no irreducibility, is used. The lattice $M_{f,\lambda}$ of \cite{DFG04} is therefore $H^1(E,\Z_\ell)=\operatorname{Hom}(T_\ell E,\Z_\ell)\simeq T_\ell E(-1)$, by Galois-equivariance of $\varphi^*$ in \'etale cohomology, and its adjoint is $\ad T_\ell E$, the adjoint being insensitive to dualising and twisting.

We also need the following piece of commutative algebra, which is the inequality half of \cite[Lemma 4.17]{DDT}. Let $\mathcal{O}$ be a discrete valuation ring with finite residue field, $\T$ a commutative $\mathcal{O}$-algebra, finite and free as an $\mathcal{O}$-module, acting on an $\mathcal{O}$-module $L$ free of finite rank with a perfect $\mathcal{O}$-bilinear pairing for which $\T$ is self-adjoint; let $\pi\colon\T\to\mathcal{O}$ be a surjective $\mathcal{O}$-algebra homomorphism, $\wp=\ker\pi$, $I=\operatorname{Ann}_\T(\wp)$ and $\eta=\pi(I)$, the congruence ideal of $\pi$; and suppose that $L[\wp]$ is free of rank $d$ over $\mathcal{O}$ with Gram determinant $\Delta\ne0$.

\begin{lemma}\label{lem:ddt417}
$\eta^d\subseteq\Delta\mathcal{O}$, that is, $d\cdot v(\eta)\ge v(\Delta)$; and equality holds if $L$ is free over $\T$.
\end{lemma}
\begin{proof}
If $\eta=0$ there is nothing to prove; let $t\in I$ with $\pi(t)$ a generator of $\eta\ne0$. First $\wp\cap I=0$: for $x\in\wp\cap I$, $tx=0$ as $t$ kills $\wp$, and $(\pi(t)-t)x=0$ as $\pi(t)-t\in\wp$ and $x$ kills $\wp$; so $\pi(t)x=0$ and $x=0$, $\T$ being torsion-free. (This does not follow from reducedness: the full Hecke algebra can have nilpotents, for instance $U_2$ at level $88$ on the span of the $f_{\text{11a}}(dz)$, $d\mid8$.) Consequently $\pi|_I$ is injective, with image $\eta=\pi(t)\mathcal{O}$, so $I=\mathcal{O}t$. The sum $L[\wp]+L[I]$ is direct: $\pi(t)=t+(\pi(t)-t)$ lies in $I+\wp$ and is a nonzero scalar, so an element of $L$ killed by $\wp$ and by $I$ is killed by $\pi(t)$, hence is $0$, $L$ being torsion-free. For $x\in L$, one has $tx\in L[\wp]$ (as $\wp I=0$) and $(\pi(t)-t)x\in L[I]$ (as $s(\pi(t)-t)\in\wp\cap I=0$ for $s\in I$), so $Q=L/(L[\wp]\oplus L[I])$ is killed by $\eta$. The pairing identifies $L/L[\wp]^\perp$ with $\operatorname{Hom}(L[\wp],\mathcal{O})$, because $L[\wp]$ is saturated, and $L[\wp]^\perp=L[I]$: for $y\in L[I]$ and $x\in L[\wp]$, $\pi(t)\langle x,y\rangle=\langle tx,y\rangle=\langle x,ty\rangle=0$; conversely if $y\perp L[\wp]$ then $\langle ty,z\rangle=\langle y,tz\rangle=0$ for all $z$, since $tz\in L[\wp]$, so $ty=0$, and $y\in L[I]$ because $I=\mathcal{O}t$. Hence $Q$ is the cokernel of $L[\wp]\to\operatorname{Hom}(L[\wp],\mathcal{O})$, of order $\#\mathcal{O}/\Delta$, and it is generated by $d$ elements and killed by $\eta$, so $\#\mathcal{O}/\Delta\le\#(\mathcal{O}/\eta)^d$. If $L\simeq\T^d$ then $Q\simeq(\T/(\wp+I))^d\simeq(\mathcal{O}/\eta)^d$.
\end{proof}

\begin{lemma}\label{lem:hecke}
Let $E$ be $X_0(N)$-optimal and $\ell$ an odd prime with $\ell\nmid 2N$ and $\bar\rho$ irreducible. Then $v_\ell(\degphi)=v_\ell(r_E)=v_\lambda(\eta^\emptyset_{f,\lambda})$.
\end{lemma}
\begin{proof}
$v_\ell(\degphi)=v_\ell(r_E)$ is \cite[Theorem 2.2]{ARS12}, as $\ell^2\nmid N$, and $v_\ell(\degphi)=v_\lambda(\eta^\emptyset_{f,\lambda})$ is Proposition~\ref{prop:pullback}.
\end{proof}

Let $\mathfrak m$ be the maximal ideal of $\T\otimes\Z_\ell$ containing $\ell$ and $\ker\phi_f$. Lemma~\ref{lem:ddt417} applied to $L=H^1(X_0(N),\Z_\ell)_{\mathfrak m}$ gives, without the projection formula, only $v_\ell(r_E)\ge v_\lambda(\eta^\emptyset_{f,\lambda})$, with equality when $L$ is free over $\T_{\mathfrak m}$. That argument is the one used in Step 2 of the proof of Theorem~\ref{thm:quant}, for a depleted form that is not attached to an optimal quotient. In it the congruence ideal is computed in $\T_{\mathfrak m}$: the other local factors of $\T\otimes\Z_\ell$ lie in the kernel $\wp$, $\phi_f$ factors through $\T_{\mathfrak m}$, and $\operatorname{Ann}(\wp)$ decomposes along the idempotents; and the idempotent $e_{\mathfrak m}$ is self-adjoint for $\hat\delta_!$, so $H^1(X_0(N),\Z_\ell)$ splits orthogonally into its local components and the pairing is perfect on each. Without the hypotheses the inequality of Lemma~\ref{lem:ddt417} can be strict: at 99a1 and $\ell=3$, where $3\mid N$, $\degphi=4$ and $r_E=12$, so $H^1(X_0(99),\Z_3)_{\mathfrak m}$ is not free over $\T_{\mathfrak m}$.

\begin{remark}\label{rem:free}
Freeness of $H^1(X_0(N),\Z_\ell)_{\mathfrak m}$ over $\T_{\mathfrak m}$ holds when $J_0(N)[\mathfrak m]$ is two-dimensional, since the localised Hecke algebra is then Gorenstein and the Tate module is free of rank two \cite{Tilouine97}; and $J_0(N)[\mathfrak m]$ is two-dimensional for $\ell\nmid2N$ and $\bar\rho$ irreducible: \cite[Theorem 3.5]{RibetStein} gives it for $J_1(N)$, the only possible exception there having Serre weight $\ell$, which does not occur in weight $2$ with $\ell\nmid N$, and $J_0(N)[\mathfrak m]$ injects into $J_1(N)[\mathfrak m']$ under the pull-back map, whose kernel, the Shimura subgroup, is Eisenstein. (The identification of $T_\ell J_0(N)$ with $H^1(X_0(N),\Z_\ell)$ as Hecke modules uses Poincar\'e duality composed with $w_N$, which is not Galois equivariant; freeness transfers nevertheless.) After Proposition~\ref{prop:pullback} this is not needed for Theorem~\ref{thm:main} or for Theorem~\ref{thm:quant}; with Lemma~\ref{lem:ddt417} it gives a second proof of $v_\ell(r_E)=v_\ell(\degphi)$ for $\ell\nmid2N$ and $\bar\rho$ irreducible, and it compares the Hecke and the cohomological congruence ideals of forms not attached to optimal quotients.
\end{remark}

\subsection{The theorem of Diamond, Flach and Guo}
Let now $f=f_E$, $\Sigma$ the set of primes dividing $N$, $\lambda=\ell\nmid2N$, $A=A_{f,\lambda}=\ad T$ (by Proposition~\ref{prop:pullback}, $\varphi_E^*$ identifies the lattice $M_{f,\lambda}$ of \cite{DFG04} with $H^1(E,\Z_\ell)\simeq T_\ell E(-1)$, so its adjoint is $\ad T$), $W=A\otimes\Q_\ell/\Z_\ell$, $T^D=\operatorname{Hom}(A,\Z_\ell(1))\simeq A(1)$ (via the trace form), $V^D=T^D\otimes\Q_\ell$ and $W^D=V^D/T^D$. Let $H^1_f(\Q,W)$ be the Bloch-Kato Selmer group and $H^1_\Sigma(\Q,W)\supseteq H^1_f(\Q,W)$ the group defined in \cite[\S2.1]{DFG04}, with no condition at the primes of $\Sigma$, the Bloch-Kato condition at $\ell$ and the unramified condition elsewhere. Let $S_f$ be the set of primes $\lambda$ such that $\lambda\mid 2N$ or $\bar\rho|_{\Gal(\overline\Q/\Q(\sqrt{\ell^*}))}$ is not absolutely irreducible \cite[(31)]{DFG04}; the hypothesis of Theorem~\ref{thm:main} is $\ell\notin S_f$.

\begin{proposition}\label{prop:Sf}
Let $\ell$ be an odd prime with $\ell\nmid N$ and $E[\ell]$ irreducible. If $\ell\ge5$, then $\ell\notin S_f$. If $\ell=3$, then $3\in S_f$ if and only if the image of $\bar\rho_{E,3}$ is conjugate to $N_s(\F_3)$, and then $E$ is ordinary at $3$. In all cases $\ell\notin S_f$ if and only if $h_\ell(E)=0$.
\end{proposition}

\begin{proof}
Let $K=\Q(\sqrt{\ell^*})$. As $\bar\rho$ is absolutely irreducible (complex conjugation has distinct eigenvalues), $\bar\rho|_{G_K}$ fails to be absolutely irreducible exactly when $\bar\rho\simeq\bar\rho\otimes\chi_{\ell^*}$ (Clifford theory, as in the proof of Lemma~\ref{lem:h0}). Suppose so, and restrict to the inertia group $I_\ell$, where $\chi_{\ell^*}=\omega^{(\ell-1)/2}$, with $\omega$ the fundamental character of level one, of order $\ell-1$. As $E$ has good reduction at $\ell$, by \cite[\S1]{Serre72} either $E$ is ordinary at $\ell$ and the semisimplification of $\bar\rho|_{I_\ell}$ is $\omega\oplus1$, or $E$ is supersingular at $\ell$ and $\bar\rho|_{I_\ell}\otimes\F_{\ell^2}\simeq\omega_2\oplus\omega_2^\ell$, with $\omega_2$ of order $\ell^2-1$ and $\omega_2^{\ell+1}=\omega$. In the ordinary case $\{\omega,1\}=\{\omega^{(\ell+1)/2},\omega^{(\ell-1)/2}\}$; as $\omega^{(\ell-1)/2}\ne1$, this forces $\omega^{(\ell+1)/2}=1$, so $(\ell-1)\mid(\ell+1)/2$ and $\ell=3$. In the supersingular case $\omega^{(\ell-1)/2}=\omega_2^{(\ell^2-1)/2}=-1$, and $\{\omega_2,\omega_2^\ell\}=\{-\omega_2,-\omega_2^\ell\}$ would need $\omega_2^{\ell-1}=-1=\omega_2^{(\ell^2-1)/2}$, so $(\ell+1)\mid(\ell-1)/2$, which is impossible. So $\ell\ge5$ never occurs, and for $\ell=3$ the curve is ordinary at $3$; the image is then $N_s(\F_3)$ by Proposition~\ref{prop:ns}. The last statement follows from Lemma~\ref{lem:h0}.
\end{proof}

So for $\ell\nmid2N$ the hypothesis of Theorem~\ref{thm:main} reads: $E[\ell]$ is irreducible, and, if $\ell=3$, the mod $3$ image is not $N_s(\F_3)$.

\begin{theorem}[{\cite[Theorem 3.7]{DFG04}}]\label{thm:dfg37}
If $\ell\notin S_f$, then $H^1_\Sigma(\Q,W)$ is finite of length $v_\lambda(\eta^\Sigma_{f,\lambda})$.
\end{theorem}
The theorem is stated in \cite{DFG04} for a finite set $\Sigma$ not containing $\ell$ such that $M_{f,\lambda}$ is minimally ramified outside $\Sigma$, which holds for the set of primes dividing $N$, outside of which $\rho$ is unramified; this is the form given as Theorem 0.2 there. It covers $\ell=3$ in weight $2$: $S_f$ excludes only the primes dividing $2N$ and the induced case.

\begin{theorem}[{\cite[Theorem 2.7]{DFG04}}]\label{thm:dfg27}
If $\ell\notin S_f$, then $H^0(\Q,V^D)=H^1_f(\Q,V^D)=0$ (and the same for $A\otimes\Q_\ell$).
\end{theorem}

\begin{lemma}[{\cite[Lemma 2.1]{DFG04}}]\label{lem:dfg21}
Suppose $H^0(\Q,V^D)=H^1_f(\Q,V^D)=0$, and let $\Sigma$ be a nonempty finite set of primes, with $\ell\notin\Sigma$, containing every prime $p\ne\ell$ at which $W$ is ramified. Put $H^1_f(\Q_p,T^D)=\iota^{-1}H^1_f(\Q_p,V^D)$, $\iota\colon H^1(\Q_p,T^D)\to H^1(\Q_p,V^D)$. Then there is an exact sequence
\[
0\to H^1_f(\Q,W)\to H^1_\Sigma(\Q,W)\to\bigoplus_{p\in\Sigma\cup\{\infty\}}H^1_f(\Q_p,T^D)^*\to H^0(\Q,W^D)^*\to0,
\]
where $^*$ is the Pontryagin dual.
\end{lemma}
The condition on $\Sigma$ implies the one of \cite{DFG04}, which asks that $\Sigma$ contain every prime at which $H^1_f(\Q_p,W)\ne H^1_{\mathrm{ur}}(\Q_p,W)$. For $p\in\Sigma$ the group $H^1_f(\Q_p,V^D)=H^1_{\mathrm{ur}}(\Q_p,V^D)=(V^D)^{I_p}/(\Frob_p-1)$ vanishes, because $\det(1-\Frob_p^{-1}\mid(V^D)^{I_p})=L_p(A_f,1)^{-1}\ne0$; for the same reason $H^0(\Q_p,V^D)=0$, so the hypothesis $H^0(\Q,V^D)=0$ is automatic here, and only $H^1_f(\Q,V^D)=0$ needs Theorem~\ref{thm:dfg27}. So $H^1_f(\Q_p,T^D)$ is the torsion of $H^1(\Q_p,T^D)$, a finite group. This computation needs $p\ne\ell$, which is why $\ell\notin\Sigma$ is part of the lemma: at $p=\ell$ the group $H^1_f(\Q_\ell,V^D)$ has dimension $2$. At $\infty$, $H^1(\R,T^D)$ is killed by $2$, so $H^1_f(\R,T^D)=0$ for $\ell$ odd.

\begin{lemma}[{\cite[I.4.2.2]{FPR94}, as recalled in the proof of \cite[Theorem 2.15]{DFG04}}]\label{lem:fpr}
For $p\in\Sigma$,
\[
\Fitt_{\Z_\ell}H^1_f(\Q_p,T^D)=L_p(A_f,1)^{-1}\cdot\Tam_p(T^D),\qquad\text{and}\qquad\Tam_p(T^D)=\Tam_p(A)
\]
\cite[(57)]{DFG04}; here $L_p(A_f,1)^{-1}=L_p(A_f(1),0)^{-1}$ is a nonzero rational number which is an $\ell$-adic integer, of valuation $e_\ell(E,p)$ at an additive prime and $v_\ell(1-p^{-2})$ at a multiplicative one.
\end{lemma}

\begin{lemma}\label{lem:h0}
Let $\ell$ be odd with $\bar\rho$ irreducible. Then $H^0(\Q,W^D)=H^0(\Q,\ad T(1)\otimes\Q_\ell/\Z_\ell)$ is finite and cyclic, of length $h_\ell(E)$, the largest $k$ such that $\rho\bmod\ell^k\simeq(\rho\otimes\chi_\ell)\bmod\ell^k$. Moreover $h_\ell(E)=0$ for every $\ell\ge5$, and $h_3(E)\le1$, with equality if and only if $\bar\rho_{E,3}\simeq\bar\rho_{E,3}\otimes\chi_{-3}$, if and only if $\bar\rho_{E,3}$ is not absolutely irreducible on $\Gal(\overline\Q/\Q(\sqrt{-3}))$, if and only if the image of $\bar\rho_{E,3}$ is conjugate to the normaliser $N_s(\F_3)$ of a split Cartan subgroup. This holds for every $E$, with or without complex multiplication; in particular $H^0(\Q,W^D)=0$ if $\ell\notin S_f$.
\end{lemma}
\begin{proof}
The representation $\bar\rho$ is absolutely irreducible, complex conjugation having distinct eigenvalues $\pm1$. The $\ell^k$-torsion of $H^0(\Q,W^D)$ is $H^0(\Q,\ad(T/\ell^k)(1))$, the group of trace-zero $G_\Q$-maps $T/\ell^k\to(T/\ell^k)(1)$; every such map has trace $0$ automatically, since $\operatorname{tr}\phi=\chi_\ell(g)\operatorname{tr}\phi$ for every $g$ and $1-\chi_\ell(c)=2$ is a unit for complex conjugation $c$. For $k=1$ these are the $G_\Q$-maps $\bar\rho\to\bar\rho\otimes\bar\chi_\ell$, which form a space of dimension at most one by Schur's lemma; so $H^0(\Q,W^D)$, a submodule of $(\Q_\ell/\Z_\ell)^3$ with $\ell$-torsion of rank at most one, is finite cyclic or isomorphic to $\Q_\ell/\Z_\ell$; the determinant argument below shows that it is finite. An element of exact order $\ell^k$ is a map $\phi\colon T/\ell^k\to(T/\ell^k)(1)$ whose reduction modulo $\ell$ is nonzero, hence an isomorphism, so $\phi$ is an isomorphism by Nakayama's lemma, and $\rho\bmod\ell^k\simeq(\rho\otimes\chi_\ell)\bmod\ell^k$; conversely such an isomorphism gives such an element. Taking determinants, $\chi_\ell(g)^2\equiv1\pmod{\ell^k}$ for all $g\in G_\Q$. Since $\chi_\ell(G_\Q)=\Z_\ell^\times$ contains $2$, and $4\not\equiv1$ modulo $\ell$ for $\ell\ge5$ and modulo $9$, this forces $k=0$ for $\ell\ge5$ and $k\le1$ for $\ell=3$. So $H^0(\Q,W^D)$ is finite, $h_\ell(E)=0$ for $\ell\ge5$, and $h_3(E)\le1$, with equality if and only if $\bar\rho\simeq\bar\rho\otimes\bar\chi_3^{\cyc}$, and $\bar\chi_3^{\cyc}=\chi_{-3}$. Let $K=\Q(\sqrt{-3})$. If $\phi\colon\bar\rho\to\bar\rho\otimes\chi_{-3}$ is an isomorphism, it commutes with $\bar\rho(G_K)$ and is not scalar (otherwise $\chi_{-3}$ would be trivial), so $\bar\rho|_{G_K}$ is not absolutely irreducible. Conversely, if $\bar\rho|_{G_K}\otimes\overline\F_3$ has a stable line with character $\psi$, Frobenius reciprocity gives a nonzero map $\operatorname{Ind}_{G_K}^{G_\Q}\psi\to\bar\rho\otimes\overline\F_3$, an isomorphism since $\bar\rho$ is absolutely irreducible, and $\operatorname{Ind}\psi\otimes\chi_{-3}\simeq\operatorname{Ind}\psi$; the isomorphism $\bar\rho\simeq\bar\rho\otimes\chi_{-3}$ then exists over $\F_3$, since spaces of homomorphisms commute with extension of scalars. The last equivalence is the first statement of Proposition~\ref{prop:ns} ($\bar\rho$ irreducible means $E[3]$ irreducible), whose proof uses this lemma only for the value $h_3(E)=1$.
\end{proof}

\subsection{Proof of Theorem~\ref{thm:main}}
Let $\ell\notin S_f$, $\Sigma$ the set of primes dividing $N$. By Lemma~\ref{lem:hecke}, $v_\ell(\degphi)=v_\lambda(\eta^\emptyset_{f,\lambda})$. By Theorem~\ref{thm:dfg14},
\[
v_\lambda(\eta^\emptyset_{f,\lambda})=v_\lambda(\eta^\Sigma_{f,\lambda})-\sum_{p\in\Sigma}v_\ell\bigl(L_p^{\nv}(A_f,1)^{-1}\bigr).
\]
By Theorem~\ref{thm:dfg37}, $v_\lambda(\eta^\Sigma_{f,\lambda})=\len H^1_\Sigma(\Q,W)$. By Theorem~\ref{thm:dfg27} the hypotheses of Lemma~\ref{lem:dfg21} hold, and $H^0(\Q,W^D)=0$ by Lemma~\ref{lem:h0}, so
\[
\len H^1_\Sigma(\Q,W)=\len H^1_f(\Q,W)+\sum_{p\in\Sigma}\len H^1_f(\Q_p,T^D),
\]
and by Lemma~\ref{lem:fpr} each local term is $v_\ell(L_p(A_f,1)^{-1})+v_\ell(\Tam_p(A))$. Hence
\[
v_\ell(\degphi)=\len H^1_f(\Q,W)+\sum_{p\in\Sigma}v_\ell\bigl(\Tam_p(\ad T)\bigr)+\sum_{p\in\Sigma}v_\ell\Bigl(\frac{L_p(A_f,1)^{-1}}{L^{\nv}_p(A_f,1)^{-1}}\Bigr),
\]
and the last sum is $\sum_{p\mid N}e_\ell(E,p)$ by Definition~\ref{def:terms}: at a multiplicative prime the true and the naive factors coincide, at an additive prime the naive factor is $1$. This is \eqref{eq:identity}. For \eqref{eq:thm}, the Selmer length is nonnegative, $v_\ell(\Tam_p(\ad T))\ge t_\ell(E,p)$ at every potentially multiplicative prime $p$, multiplicative or additive, $p=2$ included, by Lemma~\ref{lem:tamfactor}, and $t_\ell(E,p)=0$ at the other primes. \qed

\begin{corollary}\label{cor:explicit}
Under the hypothesis of Theorem~\ref{thm:main},
\begin{multline*}
v_\ell(\degphi)=\len H^1_f(\Q,\ad T\otimes\Q_\ell/\Z_\ell)+\sum_{q\mid N}e_\ell(E,q)\\
+\sum_{q:\,v_q(j_E)<0}\Bigl(t_\ell(E,q)+\min\bigl(t_\ell(E,q),v_\ell(q-1),v_\ell(c(u_q))\bigr)\Bigr),
\end{multline*}
where $u_q$ is the unit part of the Tate parameter at $q$ and $c(u_q)$ the class of its reduction, as in Lemma~\ref{lem:tamfactor}.
\end{corollary}
\begin{proof}
Theorem~\ref{thm:main}, with $\Tam_q(\ad T)$ trivial at the potentially good additive primes (Lemma~\ref{lem:tamadd}) and given by Lemma~\ref{lem:tamfactor} at the potentially multiplicative ones.
\end{proof}

So the excess of $v_\ell(\degphi)$ over the explicit local bound is exactly the length of the adjoint Bloch-Kato Selmer group. Table~\ref{tab:excess} gives its distribution on the optimal curves without complex multiplication of conductor below $10{,}000$ with $\ell\nmid N$, $E[\ell]$ irreducible and, for $\ell=3$, image not $\Ns$, with $e_\ell$ from PARI's \texttt{lfunsympow}. The excess is never negative, as the corollary requires; for instance 1184f1 (explicit bound $2$, $\degphi=1728$) and 1748e1 (explicit bound $3$, $\degphi=3240$) have Selmer length $1$ at $\ell=3$.

\begin{table}[ht]\centering\small
\begin{tabular}{lrrrrrr}
\toprule
$\ell$ & curves & excess $0$ & $1$ & $2$ & $3$ & $\ge4$\\
\midrule
$3$ & $13{,}203$ & $8{,}688$ & $3{,}288$ & $942$ & $248$ & $37$\\
$5$ & $23{,}361$ & $19{,}019$ & $3{,}993$ & $328$ & $19$ & $2$\\
$7$ & $28{,}134$ & $25{,}133$ & $2{,}888$ & $111$ & $2$ & $0$\\
\bottomrule
\end{tabular}
\caption{The excess $v_\ell(\degphi)-\sum_qe_\ell(E,q)-\sum_{v_q(j_E)<0}\bigl(t_\ell(E,q)+\min(t_\ell(E,q),v_\ell(q-1),v_\ell(c(u_q)))\bigr)$, which by Corollary~\ref{cor:explicit} is $\len H^1_f(\Q,\ad T\otimes\Q_\ell/\Z_\ell)$, on the optimal non-CM curves of conductor below $10{,}000$ with $\ell\nmid N$, $E[\ell]$ irreducible and, at $\ell=3$, image not $\Ns$.}\label{tab:excess}
\end{table}

\begin{remark}\label{rem:degree}
The Euler factors also appear in an exact formula. By the formula of Zagier and Watkins \cite{Zagier85,Watkins02}, $\degphi\cdot\operatorname{vol}(\Lambda_E)=4\pi^2c_E^2\langle f,f\rangle$, with the unnormalised Petersson product and the Manin constant $c_E>0$ in the numerator; and $\langle f,f\rangle=N\,L^{\nv}(\Sym f,2)/(8\pi^3)$ at every level \cite{Rankin39,Shimura76}, \cite[Lemma 2.12]{DFG04}. Hence
\[
\degphi=\frac{c_E^2\,N\,L^{\nv}(\ad f,1)}{2\pi\operatorname{vol}(\Lambda_E)}=B_E\prod_{q^2\mid N}L_q(A_f,1)^{-1},\qquad B_E=\frac{c_E^2\,N\,L(\ad f,1)}{2\pi\operatorname{vol}(\Lambda_E)}\in\Q^\times,
\]
so that $v_\ell(\degphi)=v_\ell(B_E)+\sum_{q\mid N}e_\ell(E,q)$ for every $\ell\nmid N$, with no hypothesis, and \eqref{eq:thm} is equivalent to $v_\ell(B_E)\ge\sum_qt_\ell(E,q)$. With $c_E=1$ the first formula reproduces $\degphi$ to $30$ digits for 11a1, 121a1, 539d1, 1232l1 and 726a1 (degrees $1$, $6$, $72$, $288$, $144$).
\end{remark}

\begin{remark}\label{rem:ko}
(1) At a potentially multiplicative prime $q\equiv1\pmod\ell$ with $\ell\mid n$ whose Tate parameter has an $\ell$-th power as unit part, Lemma~\ref{lem:tamfactor} gives $v_\ell(\Tam_q)\ge t_\ell(E,q)+1$, and \eqref{eq:identity} raises the bound accordingly; Section~\ref{sec:tam} tests this on the tables. (2) The formula of Kim and Ota \cite[Theorem 1.1]{KimOta} concerns the same congruence ideal: for $N=N^+N^-$ with $N^-$ square-free, $\ell\ge5$, $\ell\nmid N$, $\ell\notin S_f$, and $\bar\rho$ ramified at every $q\mid N^-$ with $q\equiv\pm1\pmod\ell$, it says $\ord_\lambda\eta_f(N)=\ord_\lambda\eta_f(N^+,N^-)+\sum_{q\mid N^-}t_f(q)$, where $\eta_f(N^+,N^-)$ is the congruence ideal in the $N^-$-new quotient of the Hecke algebra and $t_f(q)$ the Tamagawa exponent \cite[Definition 2.12]{KimOta}, which is $t_\ell(E,q)$ by Lemma~\ref{lem:tamagawa-exp}. Since $\eta_f(N^+,N^-)$ is integral, this gives $v_\ell(r_E)\ge\sum_{q\mid N^-}t_\ell(E,q)$, which \eqref{eq:thm} contains, without the conditions $\ell\ge5$ and $q\not\equiv\pm1$; the equality is theirs.
\end{remark}

\subsection{Twist monotonicity and the Euler terms without the hypothesis on \texorpdfstring{$\Q(\sqrt{\ell^*})$}{Q(sqrt l*)}}\label{sec:twisting}
\begin{theorem}\label{thm:twist}
Let $E/\Q$ have conductor $N$, $S$ a set of odd primes, $D=\prod_{q\in S}q$, $\chi_S=\prod_{q\in S}\chi_{q^*}$, and $E_S=E\otimes\chi_S$ of conductor $N_S$. If $q^2\mid N$ for every $q\in S$ and $\operatorname{lcm}(N_S,D^2)\mid N$, then $r_{E_S}\mid r_E$; if moreover $\operatorname{lcm}(N,D^2)\mid N_S$, then $r_{E_S}=r_E$.
\end{theorem}
\begin{proof}
Let $f=f_E$, $g=f_{E_S}$ and $\iota\colon h\mapsto h\otimes\chi_S=\sum\chi_S(n)a_n(h)q^n$, which maps $S_2(\Gamma_0(N_S))$ into $S_2(\Gamma_0(\operatorname{lcm}(N_S,D^2)))\subset S_2(\Gamma_0(N))$, preserves integrality, and sends an eigenform of all $T_n$, $(n,N_S)=1$, to an eigenform with eigenvalues twisted by $\chi_S(n)$. Then $\iota(g)=f$: both are multiplicative, $a_p(f)=\chi_S(p)a_p(g)$ for $p\nmid D$ since the twist is unramified outside $S$, and $a_{q^k}(f)=0$ for $q\in S$ since $q^2\mid N$. If $h$ is a generalised eigenform of level dividing $N_S$ for a system of $T_p$-eigenvalues, $p\nmid N_S$, differing from that of $g$ at infinitely many $p$, then $\iota(h)$ is a generalised eigenform for a system differing from that of $f$ at infinitely many $p$, so $\iota(h)\in f^\perp$; hence $\iota$ maps $g^\perp$, the $\T$-stable complement of the $g$-line (the sum of the generalised eigenspaces of the other systems of eigenvalues), into $f^\perp$. A congruence $g\equiv h\pmod r$ with $h\in g^\perp$ integral therefore gives $f\equiv\iota(h)\pmod r$ with $\iota(h)\in f^\perp$ integral, and $r\mid r_E$ by Lemma~\ref{lem:crit}. With $r=r_{E_S}$ this is the first claim; the second hypothesis makes the situation symmetric.
\end{proof}
The hypotheses hold when (a) $E$ has type $\mathrm{I}_n^*$ at each $q\in S$ ($N_S=N/D$), (b) type $\mathrm{I}_0^*$ at each $q\in S$ ($N_S=N/D^2$, $E_S$ good at $S$), (c) type $\mathrm{II}$, $\mathrm{III}$, $\mathrm{IV}$ or a dual at each $q\in S$, $q\ge5$ ($N_S=N$, equality). On the working set, for the $157{,}938$ pairs (curve, prime of type $\mathrm{I}_n^*$) and the $292{,}580$ pairs (curve, prime of potentially good additive type), $v_\ell(\degphi)\ge v_\ell(\deg\varphi_{E^{(q)}_0})$, with $E^{(q)}_0$ the optimal curve of the isogeny class of $E^{(q)}$, holds for $\ell\in\{3,5,7,11,13\}$, $\ell^2\nmid N$, with no exception, with equality in every pair of case (c). In cases (a) and (b) Theorem~\ref{thm:quant} predicts a difference equal to the Euler term, which is positive in case (a) at $\ell=3$ for every $q\ne3$, as $3\mid q^2-1$: for instance $v_3$ of the degrees is $2$ against $1$ for 539c1 and 77c1, $2$ against $1$ for 980i1 and 140b1, $1$ against $0$ for 121d1 and 11a1, and $2$ against $0$ for 539d1 (type $\mathrm{I}_0^*$, $e_3(E,7)=2$) and 11a1. Below conductor $10{,}000$ the difference equals the Euler term in all $25{,}477$ pairs (curve, $\ell$) with $S\ne\emptyset$, $\ell\in\{3,5,7,11,13\}$, $\ell\nmid N$ and $E[\ell]$ irreducible; when $\ell\notin S_f$ this is a theorem (Theorem~\ref{thm:quant}).

\emph{The set-up.} Let $E/\Q$ be any elliptic curve of conductor $N$, $S$ the set of odd primes at which $E$ has type $\mathrm{I}_0^*$ or $\mathrm{I}_n^*$, $D=\prod_{q\in S}q$, $E_S=E\otimes\chi_S$ of conductor $N_S$, $f=f_E$, $g=f_{E_S}$, and $\iota$ the twist map of Theorem~\ref{thm:twist}. By Lemma~\ref{lem:twist}, valid at every odd $q$ for these exchanges, $E_S$ has good reduction at the $\mathrm{I}_0^*$ primes of $S$ and multiplicative reduction at the $\mathrm{I}_n^*$ primes, so $\delta_q=2$, resp.\ $1$, for $g$ at $q\in S$, $D^2\mid N$, and $N_S\prod_{q\in S}q^{\delta_q}=N$: the $S$-depleted form $g^S$ lives at level $N$, and $\iota(g^S)=f$, since $\iota$ kills the coefficients at multiples of $D$ and so does not distinguish $g$ from $g^S$; likewise $\iota(f)=g^S$, the coefficients of $f\otimes\chi_S$ being $a_n(g)$ for $n$ prime to $D$ and $0$ otherwise. Let $\T$ be the full Hecke algebra of level $N$, $\pi\colon\T\to\Z$ the eigencharacter of $g^S$ and $\eta^{\mathrm{Hecke}}(g^S)=\pi(\operatorname{Ann}_\T\ker\pi)$ its congruence ideal. The generalised eigenspace of $\pi$ in $S_2(\Gamma_0(N))\otimes\overline\Q$ is the line of $g^S$: strong multiplicity one identifies the forms with the $T_p$-eigenvalues of $g$, $p\nmid N$, with the span of the $g(dz)$, $d\mid N/N_S$; on that span the operators $U_q$, $q\in S$, have the eigenvalue $0$ with multiplicity one, and the $U_p$ with $p\mid N_S$, $p\notin S$, act by the scalars $a_p(g)$. Hence the proof of Lemma~\ref{lem:crit}, with $f^\perp$ replaced by the $\T$-stable complement $C$ of the line of $g^S$ (the sum of the generalised eigenspaces of the other eigencharacters; not the Petersson complement, which is not $\T$-stable when the $U_p$ are involved), and carried out over $\Z_\ell$ with the pairing $S_2(\Gamma_0(N);\Z_\ell)\times\T_\ell\to\Z_\ell$, $\T_\ell=\T\otimes\Z_\ell$, shows for every prime $\ell$ that $\ell^{v_\ell(\eta^{\mathrm{Hecke}}(g^S))}$ is the largest power $\ell^t$ for which $g^S\equiv h\pmod{\ell^t}$ coefficientwise with $h\in C\cap S_2(\Gamma_0(N);\Z_\ell)$.

\begin{lemma}[the twist transport]\label{lem:step1}
In this situation, $v_\ell(r_E)\le v_\ell(\eta^{\mathrm{Hecke}}(g^S))$ for every prime $\ell$, and
\[
v_\ell(r_E)=v_\ell\bigl(\eta^{\mathrm{Hecke}}(g^S)\bigr)\qquad\text{for every prime }\ell\notin S.
\]
No other hypothesis on $\ell$ is needed: $\ell=2$, $\ell\mid N$, $E[\ell]$ reducible and $\ell\in S_f$ are allowed, and $E$ need not be optimal.
\end{lemma}

\begin{proof}
\emph{The inequality $\le$.} By Lemma~\ref{lem:crit} there is an integral $h\in f^\perp$ with $f\equiv h\pmod{r_E}$. The form $f$ is new, so $f^\perp$ is the sum of the generalised eigenspaces of the $T_p$, $p\nmid N$, for systems of eigenvalues that differ from that of $f$ at infinitely many $p$. The map $\iota$ preserves integrality, maps $S_2(\Gamma_0(N))$ into $S_2(\Gamma_0(\operatorname{lcm}(N,D^2)))=S_2(\Gamma_0(N))$, and intertwines $T_p$ with $\chi_S(p)T_p$ for $p\nmid N$; so $\iota(h)$ is integral and is a sum of generalised eigenvectors for systems that differ from that of $g$ at infinitely many $p$, whence $\iota(h)\in C$. Then $g^S=\iota(f)\equiv\iota(h)\pmod{r_E}$, and $v_\ell(r_E)\le v_\ell(\eta^{\mathrm{Hecke}}(g^S))$ for every $\ell$ by the characterisation above.

\emph{The inequality $\ge$ for $\ell\notin S$.} Let $t=v_\ell(\eta^{\mathrm{Hecke}}(g^S))$ and $\mathfrak{m}$ the maximal ideal of $\T_\ell$ containing $\ker\pi$ and $\ell$; it contains $U_q$ for $q\in S$. Take $h\in C\cap S_2(\Gamma_0(N);\Z_\ell)$ with $g^S\equiv h\pmod{\ell^t}$. The idempotent $e_{\mathfrak m}$ of $\T_\ell$ fixes $g^S$ (as $\pi$ factors through $\T_{\mathfrak m}$) and preserves $C$ and the integral forms, so $h'=e_{\mathfrak m}h$ lies in $C\cap S_2(\Gamma_0(N);\Z_\ell)$ and $g^S\equiv h'\pmod{\ell^t}$. The form $h'$ lies in the sum of the generalised eigenspaces of the eigencharacters $\phi\ne\pi$ of $\T$ congruent to $\pi$ modulo the maximal ideal above $\ell$, and such a $\phi$ has $\phi(U_q)\equiv0$ for $q\in S$. Every eigenvalue of $U_q$ on $S_2(\Gamma_0(N))$ is $0$, $\pm1$ or a root of $X^2-a_q(h_0)X+q$ for a newform $h_0$, hence $0$ or a unit at every prime not above $q$; as $\ell\ne q$, $\phi(U_q)=0$. (This is the only place where $\ell\notin S$ is used.) Each such $\phi$ comes from a newform $h_0$ of level dividing $N$; if $h_0=g$, then $\phi$ agrees with $\pi$ on the $T_p$, $p\nmid N$, on the $U_q$, $q\in S$, and on the $U_p$, $p\mid N_S$, $p\notin S$, which act by scalars on the span of the $g(dz)$; so $\phi=\pi$, which is excluded. Hence $h_0\ne g$, and $\iota$ maps the generalised eigenspace of $\phi$ into $S_2(\Gamma_0(N))$, into the generalised eigenspace of the system $\chi_S(p)a_p(h_0)$ of the $T_p$, $p\nmid N$, which differs from the system $\chi_S(p)a_p(g)$ of $f$ at infinitely many $p$ by strong multiplicity one; so $\iota(h')\in f^\perp$, integral. Then $f=\iota(g^S)\equiv\iota(h')\pmod{\ell^t}$, and the $\Z_\ell$-version of Lemma~\ref{lem:crit} gives $t\le v_\ell(r_E)$.
\end{proof}

At level $539$ the four congruence modules of $f_{\text{539c}}$, $f_{\text{539d}}$ and of the depleted forms $(f_{\text{77c}})^{\{7\}}$, $(f_{\text{11a}})^{\{7\}}$ have orders $288$, $504$, $288$ and $504$; their $7$-parts, $7\in S$, are not covered by Lemma~\ref{lem:step1}.

\begin{proof}[Proof of Theorem~\ref{thm:quant}]
Use the set-up above, with $\Sigma=S$ and $\mathcal{O}=\Z_\ell$. Note that $\ell\notin S$, as $\ell\nmid N$; in particular $\ell\ne3$ when $3\in S$.

\emph{Step 1: the twist map and the full Hecke algebra.} By Lemma~\ref{lem:step1}, $v_\ell(r_E)=v_\ell(\eta^{\mathrm{Hecke}}(g^S))$; only the inequality
\[
v_\ell(r_E)\ \ge\ v_\ell(\eta^{\mathrm{Hecke}}(g^S))
\]
is needed below. Let $\mathfrak m$ be the maximal ideal of $\T_\ell$ containing $\ker\pi$ and $\ell$. The congruence ideal $\eta^{\mathrm{Hecke}}(g^S)$ is computed in $\T_{\mathfrak m}$: the other local factors of $\T_\ell$ lie in $\ker\pi$, and the annihilator decomposes along the idempotents.

\emph{Step 2: Hecke versus cohomological congruence ideal.} Apply Lemma~\ref{lem:ddt417} to the module $L=H^1(X_0(N),\Z_\ell)_{\mathfrak m}$, with the pairing of Section~\ref{sec:dfg}, perfect on the $\mathfrak m$-component because $e_{\mathfrak m}$ is self-adjoint, and to $\pi$: $L[\wp]$ is the rank-two lattice $M_{g^S,\lambda}$ (Proposition~\ref{prop:integral}; it lies in the $\mathfrak m$-component), whose Gram determinant is $(\eta^S_{g,\lambda})^2$ up to a unit. Hence
\[
v_\ell(\eta^{\mathrm{Hecke}}(g^S))\ \ge\ v_\lambda(\eta^S_{g,\lambda}).
\]
Freeness of $L$ over $\T_{\mathfrak m}$ is not needed for this inequality.

\emph{Step 3: the formula of Diamond, Flach and Guo.} By Theorem~\ref{thm:dfg14} applied to $g$ (its residual representation $\bar\rho\otimes\chi_S$ is irreducible, $\ell\nmid 2N_S$ and $\ell\notin\Sigma=S$),
\[
v_\lambda(\eta^S_{g,\lambda})=v_\lambda(\eta^\emptyset_{g,\lambda})+\sum_{q\in S}v_\ell\bigl(L_q^{\nv}(A_g,1)^{-1}\bigr)=v_\lambda(\eta^\emptyset_{g,\lambda})+\sum_{q\in S}e_\ell(E,q),
\]
the last equality because $A_g=A_f$ and, at $q\in S$, $g$ has $\delta_q\ge1$, so that its naive factor is the true one, which is the factor that $f$ (with $a_q(f)=0$) omits; the values are those of Lemma~\ref{lem:euler}, whose rows $\mathrm{I}_n^*$ and $\mathrm{I}_0^*$ hold at every odd $q$, $q=3$ included.

\emph{Step 4: the level $N_S$.} By Lemma~\ref{lem:hecke} applied to the optimal curve $E_S'$ of the isogeny class of $E_S$ ($\ell\nmid2N_S$, and $\bar\rho_g$ is irreducible), $v_\lambda(\eta^\emptyset_{g,\lambda})=v_\ell(\deg\varphi_{E_S'})=v_\ell(r_{E_S})$, with no freeness hypothesis. Putting the steps together and using $v_\ell(\degphi)=v_\ell(r_E)$ (as $\ell\nmid N$) proves the first inequality; the second holds because $r_{E_S}$ is an integer.

\emph{The equality.} Let $\ell\notin S_f$. Then $\ell\notin S_{f_{E_S}}$, since $\bar\rho_{E_S,\ell}=\bar\rho\otimes\chi_S$ and $N_S\mid N$. Apply Corollary~\ref{cor:explicit} to $E$ and to $E_S'$. The $G_\Q$-modules $\ad T_\ell E$ and $\ad T_\ell E_S'$ are isomorphic, the twist and the isogenies (of degree prime to $\ell$) not changing the adjoint, so the Selmer groups coincide, and so do the Tamagawa factors, which are local invariants of $\ad T$. At a prime $q\mid N$ not in $S$, $\chi_S$ is unramified, so $E$ and $E_S$ have the same type (Lemma~\ref{lem:twist}(i)), the same naive and true factors, and the same $e_\ell$. At $q\in S$, $e_\ell(E_S,q)=0$, $E_S$ being multiplicative or good there, while the Tamagawa factors agree: both are trivial at an $\mathrm{I}_0^*$ prime (Lemma~\ref{lem:tamadd}, and $E_S$ is good), and at an $\mathrm{I}_n^*$ prime they coincide because $\ad T$ does. Hence $v_\ell(\degphi)-v_\ell(\deg\varphi_{E_S'})=\sum_{q\in S}e_\ell(E,q)$, and $v_\ell(\deg\varphi_{E_S'})=v_\ell(r_{E_S})$ as $\ell^2\nmid N_S$.
\end{proof}

\begin{remark}\label{rem:afterquant}
(1) Theorem~\ref{thm:quant} is the quantitative form of the twist-and-lower argument of Section~\ref{sec:proof}, read in the other direction: the level is raised, from $N_S$ to $N_S\prod q^{\delta_q}=N$, and the price is the level-raising Euler factor. For $\ell=3$ it says that a single prime of type $\mathrm{I}_n^*$ (necessarily $q\ge5$), with $3\mid n$ or not, forces $3\mid\degphi$, and that a prime of type $\mathrm{I}_0^*$ with $q\equiv1\pmod3$ does too; neither is in Theorem~\ref{thm:qual}, and neither needs the hypothesis on $\Q(\sqrt{-3})$. (2) When $\ell\notin S_f$ every inequality of the proof is an equality, in particular that of Step 2; no multiplicity one at level $N$ or $N_S$ is used. Step 1 is an equality at every prime $\ell\notin S$, with no hypothesis on $\ell$ at all (Lemma~\ref{lem:step1}): the twist map sends $f$ to $g^S$ and $f^\perp$ into the $\T$-stable complement of $g^S$. (3) Step 1 is where the argument would fail for the potentially good types $\mathrm{II}$, $\mathrm{IV}$, $\mathrm{III}$ and their duals: there the quadratic twist does not lower the level, $f$ is its own depletion, and Theorem~\ref{thm:dfg14} gives no factor. For these types the omitted factor is reached only through Theorem~\ref{thm:dfg37}, that is, with the hypothesis on $\Q(\sqrt{\ell^*})$, and Section~\ref{sec:anomaly} shows that the hypothesis cannot be dropped there.
\end{remark}

\section{The boundary of the hypothesis}\label{sec:boundary}

\begin{proof}[Proof of Proposition~\ref{prop:ns}]
Since $\det\bar\rho_{E,3}=\chi_{-3}$ cuts out $\Q(\sqrt{-3})$, the image of $\Gal(\overline\Q/\Q(\sqrt{-3}))$ is the subgroup of determinant $1$ of the image $H$ of $\bar\rho_{E,3}$. When $E[3]$ is irreducible, $H$ is conjugate to one of $\GL_2(\F_3)$, the normaliser $N_{ns}(\F_3)$ of a nonsplit Cartan subgroup (order $16$) and the normaliser $N_s(\F_3)$ of a split Cartan subgroup (order $8$); the list is in \cite{Zywina15}, and we derive it directly: these are the subgroups with surjective determinant acting irreducibly on $\F_3^2$ that contain an element of determinant $-1$ and trace $0$ (complex conjugation), the nonsplit Cartan subgroup itself being excluded by the last condition, and the list follows from these three properties by inspection of the subgroups of $\GL_2(\F_3)$, with or without complex multiplication. Their determinant-one subgroups are $\SL_2(\F_3)$, the quaternion group $Q_8$ and the cyclic group $C_4$ of order $4$ generated by $J=\begin{psmallmatrix}0&1\\-1&0\end{psmallmatrix}$. The first two act absolutely irreducibly on $\F_3^2$: a subgroup of $\GL_2(\overline\F_3)$ fixing a line lies in a Borel subgroup, so its commutator subgroup is abelian of exponent $3$, whereas the commutator subgroups $\{\pm1\}$ of $Q_8$ and $Q_8$ of $\SL_2(\F_3)$ are not. The third is abelian, hence absolutely reducible. In the third case $J$ has the characteristic polynomial $X^2+1$, irreducible over $\F_3$, so the restriction of $\bar\rho$ to $\Gal(\overline\Q/\Q(\sqrt{-3}))$ is irreducible over $\F_3$ and splits as $\psi\oplus\psi^{-1}$ over $\F_9$, for a character $\psi$ of order $4$ with values in $\F_9^\times$; by Frobenius reciprocity $\bar\rho\otimes\F_9\simeq\operatorname{Ind}\psi$, so $\bar\rho\simeq\bar\rho\otimes\chi_{-3}$, the isomorphism descending to $\F_3$, and $H^0(\Q,\ad\bar\rho(1))\ne0$ and $h_3(E)=1$ by Lemma~\ref{lem:h0}. Finally the four roots of the $3$-division polynomial correspond to the four lines of $\F_3^2$, on which $N_s(\F_3)$ has two orbits of size two (the two lines of the Cartan subgroup and the two others), whereas $N_{ns}(\F_3)$ and $\GL_2(\F_3)$ are transitive; so the division polynomial, which has no rational root when $E[3]$ is irreducible, factors as two quadratics in the first case and is irreducible otherwise.
\end{proof}

The curves with image $N_s(\F_3)$ have a rigid local structure, which accounts for the observations of Section~\ref{sec:anomaly}.

\begin{proposition}\label{prop:ns-structure}
Let $E/\Q$, with or without complex multiplication, have mod $3$ image conjugate to $N_s(\F_3)=\{\pm1,\pm J,\pm h,\pm w\}$, with $J=\begin{psmallmatrix}0&1\\-1&0\end{psmallmatrix}$, $h=\operatorname{diag}(1,-1)$ and $w=hJ$; let $C_s=\{\pm1,\pm h\}$ be the split Cartan subgroup, $C_4=\langle J\rangle=N_s(\F_3)\cap\SL_2(\F_3)$, and let $K_1$, $K_2$ be the quadratic fields cut out by $\bar\rho_{E,3}^{-1}(C_s)$ and $\bar\rho_{E,3}^{-1}(\{\pm1,\pm w\})$, the fields of definition of the roots of the two quadratic factors of the $3$-division polynomial $\psi_3$ (Proposition~\ref{prop:ns}). Let $R$ be the set of primes $p\ne3$ with $\bar\rho_{E,3}(I_p)=C_4$.
\begin{enumerate}
\item[(i)] For every prime $p\ne3$, $\bar\rho_{E,3}(I_p)$ lies in $C_4$. Consequently every prime $p\ne3$ with $v_p(j_E)<0$ has $3\mid v_p(j_E)$, so $t_3(E,p)\ge1$; no prime $p\ne3$ has type $\mathrm{IV}$ or $\mathrm{IV}^*$; and no prime $p\ge5$ has type $\mathrm{II}$ or $\mathrm{II}^*$.
\item[(ii)] $R$ is the set of primes $p\ne3$ ramified in $K_1$, and also the set of primes $p\ne3$ ramified in $K_2$; it is nonempty. At every $p\in R$ the reduction is potentially good, the image of $I_p$ in $\operatorname{Aut}(T_\ell E)$ is cyclic of order $4$ for every prime $\ell\ne p$, $(\ad V_\ell E)^{I_p}$ is a line on which either Frobenius acts by $\chi_{-3}(p)$, $P_p(X)=1-\chi_{-3}(p)pX$ and $e_3(E,p)=v_3(p-\chi_{-3}(p))\ge1$; in particular $P_2(X)=1+2X$ and $e_3(E,2)=1$ if $2\in R$.
\item[(iii)] If $p\in R$ and $p\ge5$, then $E$ has type $\mathrm{III}$ or $\mathrm{III}^*$ at $p$ and $p\equiv\pm1\pmod{12}$. Conversely every prime $p\ge5$ of type $\mathrm{III}$ or $\mathrm{III}^*$ lies in $R$.
\item[(iv)] The splitting field of $\psi_3$ is biquadratic, with quadratic subfields $K_1$, $K_2$ and $\Q(\sqrt{-3})$.
\item[(v)] $a_p(E)\equiv0\pmod3$ for every prime $p\nmid3N$ with $p\equiv2\pmod3$.
\item[(vi)] If $3\nmid N$, then $E$ is ordinary at $3$.
\item[(vii)] $h_3(E)=1$.
\end{enumerate}
\end{proposition}

\begin{proof}
Write $\bar\rho=\bar\rho_{E,3}$. (i) For $p\ne3$ the determinant $\chi_{-3}$ of $\bar\rho$ is unramified at $p$, so $\bar\rho(I_p)\subseteq N_s(\F_3)\cap\SL_2(\F_3)=C_4$, which has no element of order $3$ or $6$. If $v_p(j_E)<0$, then on $I_p$ the representation $\rho_{E,3}$ is $\chi(\sigma)\begin{psmallmatrix}1&n\theta(\sigma)\\0&1\end{psmallmatrix}$ in a suitable basis, with $\chi$ quadratic, $n=-v_p(j_E)$ and $\theta$ the $3$-adic tame character (proof of Lemma~\ref{lem:tamfactor}); on the kernel of $\chi$, where $\theta$ takes a unit value, this is a unipotent of order $3$ modulo $3$ unless $3\mid n$. So $3\mid n$. At a prime $p\ne3$ of type $\mathrm{IV}$ or $\mathrm{IV}^*$, $\bar\rho(I_p)$ has order $3$ (Lemma~\ref{lem:IV}), and at a prime $p\ge5$ of type $\mathrm{II}$ or $\mathrm{II}^*$, $\rho(I_p)$ is cyclic of order $6$ \cite{Kraus90} and injects into $\GL_2(\F_3)$; neither is possible.

We prove (iv) next, because (ii) uses the fields $K_1$ and $K_2$ it describes.

(iv) The splitting field of $\psi_3$ is the field cut out by the kernel $\bar\rho^{-1}(\{\pm1\})$ of the action on the four lines of $E[3]$, with group $N_s(\F_3)/\{\pm1\}\simeq(\Z/2)^2$. Its quadratic subfields are the fixed fields of the three subgroups of index two containing $-1$: $C_4$, whose fixed field is $\Q(\sqrt{-3})$, the field of the determinant; $C_s$, the stabiliser of each of the two lines of the Cartan subgroup; and $\{\pm1,\pm w\}$, the stabiliser of each of the two other lines. The last two give $K_1$ and $K_2$.

(ii) Let $p\ne3$. By (i), $\bar\rho(I_p)\subseteq C_4$, and $C_4\cap C_s=C_4\cap\{\pm1,\pm w\}=\{\pm1\}$. So $\bar\rho(I_p)\not\subseteq C_s$ if and only if $\bar\rho(I_p)=C_4$ if and only if $\bar\rho(I_p)\not\subseteq\{\pm1,\pm w\}$; that is, $p$ is ramified in $K_1$ if and only if $p\in R$ if and only if $p$ is ramified in $K_2$. As $C_s$ contains $h$, of determinant $-1$, $C_s\ne C_4$ and $K_1\ne\Q(\sqrt{-3})$; every quadratic field other than $\Q(\sqrt{-3})$ is ramified at some prime $p\ne3$, so $R\ne\emptyset$. (This does not use $3\nmid N$.) Let $p\in R$. The reduction at $p$ is not potentially multiplicative, since then $\bar\rho(I_p)$ would consist of plus or minus unipotent elements, and $C_4$ has elements of order $4$; so $\rho_{E,3}(I_p)$ is finite, injects into $\GL_2(\F_3)$, and is cyclic of order $4$, generated by an element $\tau$ with $\tau^2=-1$ reducing to $J^{\pm1}$; the kernel of $\rho_{E,\ell}|_{I_p}$ does not depend on $\ell\ne p$ \cite[\S2]{SerreTate68}, so the image is cyclic of order $4$ on every $T_\ell E$. The trace-zero endomorphisms commuting with $\tau$ form the line spanned by $\tau$. Frobenius normalises $\langle\tau\rangle$, so it conjugates $\tau$ to $\tau^{\pm1}$, and the sign can be read modulo $3$. The elements of $N_s(\F_3)$ not in $C_4$ form the coset $hC_4$, of determinant $-1$, and they conjugate $J$ to $J^{-1}$, while those of $C_4$ commute with $J$; so $\bar\rho(\Frob_p)$ commutes with $J$ if and only if its determinant $\chi_{-3}(p)$ is $1$. Hence either Frobenius acts on the line by $\chi_{-3}(p)$, $L_p(A_f,s)^{-1}=1-\chi_{-3}(p)p^{-s}$, $P_p(X)=1-\chi_{-3}(p)pX$, and, the reduction being additive, $e_3(E,p)=v_3(1-\chi_{-3}(p)p^{-1})=v_3(p-\chi_{-3}(p))\ge1$, as $p\equiv\chi_{-3}(p)\pmod3$. At $p=2$, $\chi_{-3}(2)=-1$.

(iii) For $p\ge5$ the image is tame, so $\rho(\Frob_p)\tau\rho(\Frob_p)^{-1}=\tau^p$; by (ii) this is $\tau^{\chi_{-3}(p)}$, so $p\equiv\chi_{-3}(p)\pmod4$, that is, $\chi_{-4}(p)=\chi_{-3}(p)$ and $p\equiv\pm1\pmod{12}$. A cyclic image of order $4$ means semistability defect $4$, that is, $v_p(\Delta_E)\in\{3,9\}$, the types $\mathrm{III}$, $\mathrm{III}^*$ \cite{Kraus90}. Conversely, at a prime $p\ge5$ of type $\mathrm{III}$ or $\mathrm{III}^*$ the image $\rho_{E,3}(I_p)$ is cyclic of order $4$ \cite{Kraus90} and injects into $\GL_2(\F_3)$, so $\bar\rho(I_p)$ is a subgroup of order $4$ of $C_4$, and $p\in R$.

(v) For $p\nmid3N$ with $p\equiv2\pmod3$, $\det\bar\rho(\Frob_p)=\chi_{-3}(p)=-1$, so $\bar\rho(\Frob_p)$ lies in $hC_4=\{\pm h,\pm w\}$, whose elements have trace $0$; and $a_p(E)\equiv\operatorname{tr}\bar\rho(\Frob_p)\pmod3$.

(vi) The image $N_s(\F_3)$ acts irreducibly on $\F_3^2$, so $E[3]$ is irreducible, and $3\in S_f$ by Proposition~\ref{prop:ns}; if $3\nmid N$, the proof of Proposition~\ref{prop:Sf} shows that $E$ is ordinary at $3$. (vii) is Proposition~\ref{prop:ns}.
\end{proof}

The argument of (ii) does not use $3\nmid N$. Among the $90$ optimal curves without complex multiplication with image $\Ns$ of conductor below $10{,}000$ ($44$ of them with $3\mid N$), $74$ have a prime $p\ge5$ in $R$, of type $\mathrm{III}$ or $\mathrm{III}^*$ and $p\equiv1$ or $11\pmod{12}$, and $16$ have $R=\{2\}$, where the type is $\mathrm{III}$ or $\mathrm{III}^*$ and $P_2=1+2X$. The proposition also covers the $13$ optimal curves with complex multiplication and image $N_s(\F_3)$ in that range ($j=-32768$ or $8000$: 121b1, 256a1, 256d1, 1089e1, 1936f1, 2304c1, 2304k1, 3025a1, 5929a1, 6400f1, 6400r1, 7744a1, 7744s1); $7$ of them have $R=\{11\}$ and $6$ have $R=\{2\}$, of type $\mathrm{III}$ at $2$. That the type at $2$ is always $\mathrm{III}$ or $\mathrm{III}^*$ when $R=\{2\}$ is not proved.

\begin{remark}\label{rem:ns}
When $\ell=3$ and the image is $N_s(\F_3)$, Theorem~\ref{thm:dfg37} does not apply, nor does Theorem~\ref{thm:dfg27}. The exact sequence of Lemma~\ref{lem:dfg21} still applies if $H^1_f(\Q,V^D)=0$ ($H^0(\Q,V^D)=0$ being automatic), and it then reads
\[
\len H^1_\Sigma(\Q,W)=\len H^1_f(\Q,W)+\sum_{p\in\Sigma}\bigl(v_\ell(L_p(A_f,1)^{-1})+v_\ell(\Tam_p(A))\bigr)-h_\ell(E).
\]
The vanishing of $H^1_f(\Q,V^D)$, equivalently of $H^1_f(\Q,\ad V_3E)$, is not known for these curves: we assume it here. If in addition $v_\lambda(\eta^\Sigma_{f,\lambda})$ were still the length of $H^1_\Sigma(\Q,W)$, as in Theorem~\ref{thm:dfg37}, then \eqref{eq:identity} would hold with $h_\ell(E)$ subtracted on the right: this is the shape of Conjecture~\ref{conj:euler}. Section~\ref{sec:anomaly} shows that on the tables the correction is exactly what is observed: $h_3(E)=1$ for every curve with image $N_s(\F_3)$, the inequality \eqref{eq:thm} fails for $31\%$, $23\%$ and $18\%$ of them on the three ranges, always by exactly one, and never for any other curve.
\end{remark}

\section{The computations}\label{sec:data}

\subsection{Data and protocol}
We used Cremona's tables \cite{Cremona97} in the form of the \texttt{ecdata} repository at commit \texttt{25cec5e}, which contain every elliptic curve over $\Q$ of conductor below $500000$ with its Kodaira symbols, Tamagawa numbers, modular degree, isogeny class, torsion and the images of the mod $\ell$ Galois representations in the notation of \cite{Sutherland16}. We use these labels as they appear in the \texttt{galrep} file of the tables, and with each label we give its intrinsic description (for $\Ns$, a $3$-division polynomial that is a product of two quadratics; for \texttt{3Cs.1.1} and \texttt{3B.1.1}, a rational point of order $3$, Section~\ref{sec:eis}); we have not compared them with the tables of \cite{Sutherland16} themselves. Optimality is determined in the tables for every class of conductor at most $400000$ and for $359{,}009$ further classes up to $500000$. The ranges are the working set ($N\le300000$), the hold-out ($300000<N\le400000$), the extension ($400000<N\le500000$) and the out-of-table set of Section~\ref{sec:oot}. The component-group inequality \eqref{eq:main} and Conjecture~\ref{conj:eis} were found on the working set, written down, and then tested once on the hold-out, which had been held back, and afterwards on the extension; nothing was adjusted after the hold-out test. The twist-symmetrised form \eqref{eq:twist} replaced an earlier, weaker formulation of the same phenomenon on the working set and was then tested once on the hold-out and the extension. The Euler-factor rule came next: it was formulated on the working set from the minima by Kodaira type (Section~\ref{sec:iso}), the terms at $2$ and $3$ were fixed afterwards by PARI's Euler factors, and the rule was tested on all four ranges; the correction term $h_\ell(E)$ was identified last, from the exceptions. Since the hold-out and the extension had already been used for \eqref{eq:main} and \eqref{eq:twist}, they are a confirmation and not a blind test for Conjecture~\ref{conj:euler}; the out-of-table curves, generated before the rule was formulated, are. They contain no curve with image $\Ns$, so the correction term $h_3$ was fitted on, and only confirmed by, data already seen; and we do not count the out-of-table curves as a blind test of \eqref{eq:twist}. Nor do we count them as a blind test of the weights $w_3$, which they could test only through curves with $3\mid N$ and potentially good reduction at $3$. On the extension range optimality is determined only for the $359{,}009$ classes named above; for a class whose optimal curve is not determined, the right side of \eqref{eq:rule} at $\ell=3$ is nevertheless the same for every curve of the class. Indeed, when $E[3]$ is irreducible the class has no $3$-isogeny, so every isogeny in it is a composite of isogenies of prime degree $p\ne3$, and an isogeny of prime degree $p$ preserves the Kodaira type at a prime $q\ne p$ of potentially good reduction \cite[Theorem 5.4(1)]{DokDok15}; so $w_3$ is constant on the class. In the presence of a $3$-isogeny, where \eqref{eq:rule} does not apply at $\ell=3$, $w_3$ can vary: below conductor $10{,}000$, among the classes with $3\mid N$, at least two curves and potentially good reduction at $3$, $w_3$ is constant on all $1{,}295$ classes without a $3$-isogeny, as it must be, and varies in $828$ of the $1{,}293$ classes with one (for instance 54a, whose curves have types $\mathrm{IV}^*$, $\mathrm{II}^*$, $\mathrm{II}$ at $3$). The other terms of \eqref{eq:rule} are invariant under isogenies of degree prime to $\ell$. The Kodaira types and discriminant exponents at the bad primes were recomputed with PARI/GP 2.15 \cite{PARI} from the $a$-invariants, the Euler factors of $\Sym f$ at every bad prime were computed with \texttt{lfunsympow} and \texttt{lfuneuler}, the Tate parameters with \texttt{ellinit} over $\Q_q$ (at $q=2$ the Tate parameter enters no count, since the extra Tamagawa term vanishes there, $v_\ell(2-1)=0$, and $t_\ell$ is read from $v_2(j_E)$; we note that PARI/GP 2.17.2 returns $q_E^5$ in place of $q_E$ for \texttt{ellinit(E, O(2\^{}k)).tate} when $v_2(j_E)=-1$, while the $q$-expansion of $j$ gives $q_E$ directly), and fifteen modular degrees spanning the configurations were recomputed with \texttt{ellmoddegree}, all agreeing with the tables. The mod $3$ images of the tables were checked against the factorisation of the $3$-division polynomial (Proposition~\ref{prop:ns}) on every curve for which it mattered. The checks on conductor below $10{,}000$ (or $30{,}000$ where stated) reported in this note were repeated with PARI/GP 2.17.2 \cite{PARI217}; on that range they reproduce the counts obtained with PARI/GP 2.15 on all $38{,}042$ optimal curves, and \texttt{ellmoddegree} agrees with the tables on all of them.

\subsection{The rule on whole ranges}
Table~\ref{tab:uniform} checks \eqref{eq:rule}, with PARI's Euler factors at every bad prime $q\ne\ell$, the terms $t_\ell$, $w_3$ as defined, and $h_3(E)=1$ for the curves with image $\Ns$ (and $h_\ell=0$ otherwise, Lemma~\ref{lem:h0}), on the four ranges at $\ell=3,5,7,11,13$: no violation at any $\ell$. It also reports the number of violations when the correction term is dropped (all are curves with image $\Ns$, Section~\ref{sec:anomaly}) and the behaviour of the rule with the terms of the types $\mathrm{III}$, $\mathrm{III}^*$ left out, which is the form that holds for every curve of the tables without any correction; the weights $w_3$ at $q=\ell=3$ only sharpen the bound, and at $\ell\ge5$ nothing is needed at $q=\ell$.

\begin{table}[ht]\centering\footnotesize
{\footnotesize
\begin{tabular}{llrrrrrr}
\toprule
set & $\ell$ & curves & \eqref{eq:rule} $>0$ & viol. & equality & viol.\ without $h$ & III-free: $>0$ / viol. / equality\\
\midrule
working set & 3 & 1,233,729 & 867,578 & 0 & 57.5\% & 364 & 861,074 / 0 / 56.8\% \\
 & 5 & 1,303,224 & 332,771 & 0 & 75.6\% & 0 & 330,477 / 0 / 75.5\% \\
 & 7 & 1,307,753 & 174,266 & 0 & 84.1\% & 0 & 173,488 / 0 / 84.1\% \\
 & 11 & 1,309,103 & 59,849 & 0 & 91.6\% & 0 & 59,495 / 0 / 91.6\% \\
 & 13 & 1,308,945 & 35,817 & 0 & 92.4\% & 0 & 35,635 / 0 / 92.4\% \\
\midrule
hold-out & 3 & 409,577 & 298,435 & 0 & 54.6\% & 68 & 295,973 / 0 / 53.9\% \\
 & 5 & 428,251 & 115,702 & 0 & 74.2\% & 0 & 114,860 / 0 / 74.0\% \\
 & 7 & 429,280 & 60,868 & 0 & 83.2\% & 0 & 60,580 / 0 / 83.2\% \\
 & 11 & 429,502 & 21,271 & 0 & 90.6\% & 0 & 21,097 / 0 / 90.6\% \\
 & 13 & 429,498 & 12,751 & 0 & 92.4\% & 0 & 12,681 / 0 / 92.3\% \\
\midrule
extension & 3 & 348,227 & 257,165 & 0 & 54.0\% & 43 & 255,507 / 0 / 53.2\% \\
 & 5 & 358,244 & 100,034 & 0 & 73.3\% & 0 & 99,401 / 0 / 73.2\% \\
 & 7 & 358,717 & 53,034 & 0 & 82.7\% & 0 & 52,766 / 0 / 82.7\% \\
 & 11 & 358,830 & 19,235 & 0 & 90.6\% & 0 & 19,136 / 0 / 90.6\% \\
 & 13 & 358,816 & 11,561 & 0 & 91.9\% & 0 & 11,484 / 0 / 91.9\% \\
\midrule
out of table & 3 & 893 & 828 & 0 & 52.5\% & 0 & 827 / 0 / 51.6\% \\
 & 5 & 893 & 112 & 0 & 82.1\% & 0 & 111 / 0 / 82.0\% \\
 & 7 & 893 & 52 & 0 & 90.4\% & 0 & 52 / 0 / 90.4\% \\
 & 11 & 893 & 19 & 0 & 94.7\% & 0 & 19 / 0 / 94.7\% \\
 & 13 & 893 & 12 & 0 & 100.0\% & 0 & 12 / 0 / 100.0\% \\
\bottomrule
\end{tabular}
}
\caption{Conjecture~\ref{conj:euler} with PARI's factors at all primes: curves with $E[\ell]$ irreducible, not CM, with a positive bound in \eqref{eq:rule}; violations; equality rate among them; violations when the correction term $h_\ell$ is dropped; and the rule without the type $\mathrm{III}$ terms (``III-free'': positive bounds, violations, equality rate).}\label{tab:uniform}
\end{table}

\subsection{Calibration and sharpness}\label{sec:sharp}
The zero counts are not tautological. Among the $366{,}151$ working-set curves with $E[3]$ irreducible, not CM, and right-hand side $0$ in \eqref{eq:rule}, the prime $3$ does not divide $\degphi$ for $53.2\%$. If the valuation of the degree were independent of the local data, the $867{,}578$ working-set curves with a positive bound at $\ell=3$ would produce roughly $605{,}000$ violations (summing over them the proportion of bound-$0$ curves whose valuation is smaller than the bound); there are none. This shows that the statement is strong, far from vacuous, not that it is valid; the figure depends on the reference population. The bounds are attained: among working-set curves with a positive bound, $v_\ell(\degphi)$ equals the right-hand side of \eqref{eq:rule} in $57.5\%$ of the cases at $\ell=3$, $75.6\%$ at $\ell=5$, $84.1\%$ at $\ell=7$, $91.6\%$ at $\ell=11$ and $92.4\%$ at $\ell=13$, and the mean excess at $\ell=3$ is $0.63$. Under the hypothesis of Theorem~\ref{thm:main} the identity \eqref{eq:identity} says what the excess is: the length of the Bloch-Kato Selmer group of the adjoint plus the excess of the Tamagawa factors over the terms $t_\ell$ (Corollary~\ref{cor:explicit}). The whole odd part of $\degphi$ is another matter: for the $1{,}226{,}216$ working-set optimal non-CM curves with no isogeny of odd degree, the odd part of $\degphi$ equals $\prod_\ell\ell^{\,b_\ell}$, with $b_\ell$ the right-hand side of \eqref{eq:twist}, in only $1.85\%$ of cases (the proportion is range dependent: it is $17.1\%$ below conductor $10{,}000$); the quotient is most often $9$, $3$, $15$, $45$, $27$ or $21$, and its smallest prime factor is $3$ in $63\%$ of the cases ($58.7\%$ below conductor $10{,}000$). Below conductor $10{,}000$ the smallest prime factor of the quotient divides $N$ in $47.5\%$ of the cases with a nontrivial quotient, and some prime factor of the quotient divides $N$ in $51.8\%$. So the local factors account for a definite, sharp part of the modular degree and for little of the rest, which is governed by the adjoint Selmer group and by the size of the Petersson norm.

\subsection{The Tamagawa factors}\label{sec:tam}
Lemma~\ref{lem:tamfactor} and the identity \eqref{eq:identity} predict that the bound can be raised, at a prime $q$ of type $\mathrm{I}_n$ with $\ell\mid n$ and $q\equiv1\pmod\ell$, by $\min(v_\ell(n),v_\ell(q-1),v_\ell(c(u)))$, where $v_\ell(c(u))\ge k$ means that the unit part $u$ of the Tate parameter at $q$ reduces to an $\ell^k$-th power in $\F_q^\times$. We computed the Tate parameters with PARI (at $q=2$ they do not enter, Section~\ref{sec:data}) for the working-set curves that satisfy the hypothesis of Theorem~\ref{thm:main} ($\ell\nmid N$, and at $\ell=3$ image not $\Ns$) and have such a prime. At $\ell=3$ there are $46{,}222$ such curves, $12{,}095$ of them with a positive extra term ($12{,}125$ in total); the strengthened inequality has no exception, and the number of curves attaining the bound rises from $20{,}822$ ($45.0\%$) to $28{,}007$ ($60.6\%$). At $\ell=5$ there are $10{,}940$ such curves, $1{,}646$ with a positive extra term; again no exception, and equality rises from $7{,}333$ ($67.0\%$) to $8{,}610$ ($78.7\%$). The same refinement applies at the primes of type $\mathrm{I}_n^*$ (Lemma~\ref{lem:tamfactor}); including them, below conductor $10{,}000$ there are $1{,}050$ such curves at $\ell=3$, with equality rising from $56.8\%$ to $68.0\%$, and $252$ at $\ell=5$, with equality rising from $203$ to $220$ curves, and no violation once the CM curves are excluded. This tests the identity, not only the inequality: the extra divisibility of the modular degree predicted by the Tamagawa factor of the adjoint is always present, the Selmer length accounting for the rest.

\subsection{Each term in isolation}\label{sec:iso}
For working-set curves with exactly one additive prime $q\ge5$, $q\ne\ell$, and no multiplicative prime with $\ell\mid n$, Table~\ref{tab:iso} gives the minimum of $v_\ell(\degphi)$ by Kodaira type and by $e_\ell(E,q)$. In every cell, for $\ell=3,5,7$, the minimum equals $e_\ell$ and every curve has $v_\ell(\degphi)\ge e_\ell$; for $\mathrm{I}_n^*$ the excess over $e_\ell$ has minimum exactly $v_\ell(n)$ in every cell. The cells with $e=0$ ($\mathrm{III}$: $5{,}534$ curves, $\mathrm{I}_0^*$: $5{,}394$) have minimum $0$, and for $\mathrm{I}_0^*$ the term runs up to $e=6$ (minimum $6$). Multiplicative primes show no Euler term: for curves with a single multiplicative prime $q$ with $\ell\mid n$ and no additive prime, $\min(v_\ell(\degphi)-v_\ell(n))=0$ for every value of $v_\ell(q^2-1)$ up to $6$ (at $\ell=3$ the value $0$ occurs only for $q=\ell=3$, since $3\mid q^2-1$ for every other $q$). For two primes of potentially good type the terms add: for the $2{,}502$ working-set curves with exactly two odd primes of type $\mathrm{II}$ or $\mathrm{II}^*$ and no other contribution, $9\mid\degphi$ always, and the mechanism is not a congruence between $f$ and its own quadratic twist by the character ramified at both primes, which never occurs ($0$ of $2{,}502$).

\begin{table}[ht]\centering\footnotesize
{\footnotesize
\begin{tabular}{l rr rr rr rr}
\toprule
 & \multicolumn{2}{c}{$e=1$} & \multicolumn{2}{c}{$e=2$} & \multicolumn{2}{c}{$e=3$} & \multicolumn{2}{c}{$e=4$}\\
type ($\ell=3$) & min & $n$ & min & $n$ & min & $n$ & min & $n$\\
\midrule
$\mathrm{II}$ & 1 & 8,327 & 2 & 821 & 3 & 59 & 4 & 2\\
$\mathrm{II}^*$ & 1 & 4,184 & 2 & 296 & 3 & 12 & 4 & 1\\
$\mathrm{IV}$ & 1 & 4,290 & 2 & 296 & 3 & 12 & 4 & 1\\
$\mathrm{IV}^*$ & 1 & 8,221 & 2 & 821 & 3 & 59 & 4 & 2\\
$\mathrm{III}$ & 1 & 1,271 & 2 & 76 & 3 & 9 & -- & --\\
$\mathrm{III}^*$ & 1 & 1,271 & 2 & 76 & 3 & 9 & -- & --\\
$\mathrm{I}_0^*$ & 1 & 1,370 & 2 & 4,432 & 3 & 2,737 & 4 & 247\\
$\mathrm{I}_n^*$ & 1 & 30,497 & 2 & 2,589 & 3 & 145 & 4 & 13\\
\bottomrule
\end{tabular}
}
\caption{Working set, $\ell=3$: minimum of $v_3(\degphi)$ (number of curves) over curves with $E[3]$ irreducible whose only additive prime is a single $q\ge5$ of the given type with $e_3(E,q)=e$, and no multiplicative $3$-contribution. The same holds at $\ell=5$ and $7$.}\label{tab:iso}
\end{table}

\subsection{The two corollary forms}\label{sec:forms}
Table~\ref{tab:counts} gives, for each range and each small prime $\ell$, the number of optimal non-CM curves with $E[\ell]$ irreducible and a positive right-hand side in \eqref{eq:main} and in \eqref{eq:twist}, the number of violations of each inequality, and the number of Eisenstein curves with a positive bound together with the violations of Conjecture~\ref{conj:eis}. There are none. Among working-set curves with a positive bound, equality holds in \eqref{eq:twist} for $39.4\%$ at $\ell=3$, $71.1\%$ at $\ell=5$, $81.2\%$ at $\ell=7$ and $91.5\%$ at $\ell=13$, and in \eqref{eq:main} for $36.9\%$, $69.7\%$, $80.4\%$ and $91.2\%$; the rule raises these to the rates of Section~\ref{sec:sharp} and bounds $867{,}578$ curves instead of $705{,}563$ at $\ell=3$. The hypothesis of Theorem~\ref{thm:main} fails on the working set for the $1{,}170$ non-CM curves with image $\Ns$ at $\ell=3$ and for the curves with $\ell\mid N$ (which include $455{,}707$ of the $913{,}244$ working-set pairs $(E,\ell)$ with $E[\ell]$ irreducible and a positive multiplicative bound at $\ell=3$, and $175{,}798$ at $\ell\ge5$); by Proposition~\ref{prop:Sf} there are no other failures among the non-CM curves with $E[\ell]$ irreducible, the curves whose mod $\ell$ representation is induced from $\Q(\sqrt{\ell^*})$ with $\ell\ge5$ all having $\ell\mid N$. (The curves with complex multiplication and $j\in\{-32768,8000\}$ also have image $N_s(\F_3)$ and fail at $\ell=3$ when $3\nmid N$, $20$ of them below conductor $10{,}000$; the \texttt{galrep} file gives them no label.) None of them violates \eqref{eq:rule}.

\begin{table}[ht]
\centering\footnotesize
{\footnotesize
\begin{tabular}{llrrrrrr}
\hline
set & $\ell$ & curves & \eqref{eq:main} $>0$ & \eqref{eq:twist} $>0$ & viol.\ \eqref{eq:main} & viol.\ \eqref{eq:twist} & Eis.\ $>0$ / viol. \\
\hline
working set, $N \le 300000$ & 3 & 1,311,066 & 598,833 & 705,563 & 0 & 0 & 67,845 / 0 \\
 & 5 & 1,311,066 & 235,510 & 262,531 & 0 & 0 & 4,454 / 0 \\
 & 7 & 1,311,066 & 121,947 & 136,707 & 0 & 0 & 759 / 0 \\
 & 11 & 1,311,066 & 42,334 & 47,040 & 0 & 0 & 0 / 0 \\
 & 13 & 1,311,066 & 27,215 & 29,943 & 0 & 0 & 26 / 0 \\
\hline
hold-out, $300000 < N \le 400000$ & 3 & 429,936 & 206,749 & 242,691 & 0 & 0 & 18,225 / 0 \\
 & 5 & 429,936 & 79,550 & 89,164 & 0 & 0 & 989 / 0 \\
 & 7 & 429,936 & 41,387 & 46,767 & 0 & 0 & 137 / 0 \\
 & 11 & 429,936 & 14,244 & 15,970 & 0 & 0 & 0 / 0 \\
 & 13 & 429,936 & 9,212 & 10,280 & 0 & 0 & 2 / 0 \\
\hline
extension, $400000 < N \le 500000$ & 3 & 359,009 & 181,075 & 214,033 & 0 & 0 & 9,894 / 0 \\
 & 5 & 359,009 & 70,514 & 78,650 & 0 & 0 & 474 / 0 \\
 & 7 & 359,009 & 36,734 & 41,317 & 0 & 0 & 70 / 0 \\
 & 11 & 359,009 & 13,347 & 14,975 & 0 & 0 & 0 / 0 \\
 & 13 & 359,009 & 8,467 & 9,465 & 0 & 0 & 5 / 0 \\
\hline
out of table, $500000 < N \le 4\cdot 10^6$ & 3 & 893 & 561 & 718 & 0 & 0 & 0 / 0 \\
 & 5 & 893 & 6 & 7 & 0 & 0 & 0 / 0 \\
 & 7 & 893 & 0 & 1 & 0 & 0 & 0 / 0 \\
\hline
\end{tabular}
}
\caption{Verification counts for the two corollary forms. ``\eqref{eq:main} $>0$'' is the number of curves with $E[\ell]$ irreducible, not CM, for which the right-hand side of \eqref{eq:main} is positive, and ``viol.'' the number of violations; likewise for \eqref{eq:twist}. The last column counts the Eisenstein curves with a positive bound in \eqref{eq:main} and the violations of Conjecture~\ref{conj:eis}.}\label{tab:counts}
\end{table}

\subsection{The boundary: the normaliser of a split Cartan subgroup at \texorpdfstring{$\ell=3$}{l=3}}\label{sec:anomaly}
Table~\ref{tab:ns} records what happens to the inequality \eqref{eq:thm}, with all types included and no correction, on the curves for which the hypothesis of Theorem~\ref{thm:main} fails at $\ell=3$. There are $364$, $68$, $43$ and $0$ violations on the four ranges, all at $\ell=3$; every one of them is a curve with image $\Ns$ in the tables, and the $3$-division polynomial of every one of them factors as two quadratics, as Proposition~\ref{prop:ns} requires; no curve with image $\GL_2(\F_3)$ or $N_{ns}(\F_3)$ violates \eqref{eq:thm}. (For the curves with $3\mid N$ the count depends on whether the terms at $q=3$ are included, that is, the term $t_3(E,3)$; the weights $w_3$ change nothing. Below conductor $10{,}000$ there are $36$ violations with them, as \eqref{eq:thm} is written, and $24$ without; the twelve curves that differ, 726a1, 726f1, 2178b1, 2178h1, 3840h1, 3840k1, 3840o1, 3840r1, 5070h1, 5070p1, 5808y1 and 5808z1, all have $t_3(E,3)=1$ from a type $\mathrm{I}_n$ or $\mathrm{I}_n^*$ at $3$.) Conversely every $\Ns$ curve of the tables has a prime of type $\mathrm{III}$ or $\mathrm{III}^*$ with a positive term $e_3$, and the inequality fails for $31\%$, $23\%$ and $18\%$ of them, always by exactly one; the violation involves a prime $p\equiv-1\pmod3$ with $t_3(E,p)\ge1$ in every case (an observation, not proved). Proposition~\ref{prop:ns-structure} proves most of the rest: every $\Ns$ curve has a prime $p\ne3$ with $e_3(E,p)\ge1$, of type $\mathrm{III}$ or $\mathrm{III}^*$ and $p\equiv\pm1\pmod{12}$ when $p\ge5$, every prime $p\ge5$ of type $\mathrm{III}$ or $\mathrm{III}^*$ is such a prime, and $t_3(E,p)\ge1$ at every potentially multiplicative prime $p\ne3$. For every $\Ns$ curve $h_3(E)=1$, and $h_3(E)=0$ for the other curves, by Lemma~\ref{lem:h0} and Proposition~\ref{prop:ns}. So the correction of Conjecture~\ref{conj:euler} is exactly one on exactly these curves, and with it there is no exception. The residue classes of the type $\mathrm{III}$ primes of the violating curves are $q\equiv1$ and $q\equiv11\pmod{12}$ in comparable numbers, with some $q=2$; what the violations have in common is the image, not the prime.

\begin{table}[ht]\centering\scriptsize\setlength{\tabcolsep}{3pt}
{\scriptsize\setlength{\tabcolsep}{3pt}
\begin{tabular}{lrrrrrrr}
\toprule
set & $\Ns$ & with III term & failing & others failing & $h_3=1$ & $q\equiv1$ / $11\ (12)$ / $q=2$ & shortfall $1$\\
\midrule
working set & 1,170 & 1,170 & 364 (31\%) & 0 & 1,170 of 1,170 & 140 / 160 / 72 & 364 of 364\\
hold-out & 290 & 290 & 68 (23\%) & 0 & 290 of 290 & 34 / 30 / 4 & 68 of 68\\
extension & 237 & 237 & 43 (18\%) & 0 & 237 of 237 & 16 / 25 / 2 & 43 of 43\\
out of table & 0 & 0 & 0 (0\%) & 0 & 0 of 0 & 0 / 0 / 0 & 0 of 0\\
\bottomrule
\end{tabular}
}
\caption{The curves with image $\Ns$ (mod $3$ image the normaliser of a split Cartan subgroup), with $E[3]$ irreducible, not CM, optimal: how many have a positive type $\mathrm{III}$ term, how many violate \eqref{eq:thm} (all types, no correction), how many other curves do, how many have $h_3(E)=1$, the residue classes mod $12$ of the type $\mathrm{III}$ primes of the violating curves (or $q=2$), and the shortfall.}\label{tab:ns}
\end{table}

\subsection{The wild primes and the prime \texorpdfstring{$\ell$}{l}}\label{sec:wild}
Table~\ref{tab:isolated} isolates the additive types at $\ell=3$ with the weights of \eqref{eq:twist}: for the untwisted types it lists, for working-set curves with no multiplicative prime $q$ with $3\mid n_q$ and all additive primes of a single Kodaira type, the minimum of $v_3(\degphi)$ by type and number of primes; for the twisted types the minimum over the curves whose only contribution to \eqref{eq:twist} comes from primes of the given type at odd $q$, or at $q=2$, with the given total weight. The rows show that the prime $2$ contributes nothing through the types $\mathrm{II}$, $\mathrm{II}^*$ ($v_3(\degphi)=0$ for $14{,}182$ of the $44{,}672$ curves whose only candidate is such a prime) and nothing through a type $\mathrm{I}_n^*$ with $3\,\|\,n$, although with $9\mid n$ the minimum is $1$. Definition~\ref{def:terms} with PARI's factor explains every one of these minima, the factor having one of the four shapes of Proposition~\ref{prop:shapes}: at $2$ the Kodaira type does not determine it; a type $\mathrm{II}$ or $\mathrm{II}^*$ at $2$ usually has $P_2=1$ and $e_3=0$, but not always (336f1 and 1232l1, Remark~\ref{rem:wild}); a type $\mathrm{IV}$ or $\mathrm{IV}^*$ at $2$ has $e_3=1$ (Lemma~\ref{lem:IV}); a type $\mathrm{I}_n^*$ at $2$ has $t_3=v_3(-v_2(j))$ when it is potentially multiplicative and $0$ otherwise, and the index $n$ of the symbol is not the valuation of the Tate parameter at $2$, which is why $n$ was not the right invariant. With the weights of \eqref{eq:twist} at $2$ counted as if $2$ were odd, \eqref{eq:twist} fails thousands of times in each range; with the terms of Definition~\ref{def:terms} it never does. At $q=3$ with potentially good reduction and $\ell=3$, the minima of $v_3(\degphi)$ over curves with a single wild additive prime are $1$ for $\mathrm{II},\mathrm{IV},\mathrm{III}^*$, $2$ for $\mathrm{II}^*,\mathrm{IV}^*$ and $0$ for $\mathrm{III},\mathrm{I}_0^*$: these are the weights $w_3$ of Conjecture~\ref{conj:euler}, which lie outside the range of Theorem~\ref{thm:main} ($\ell\mid N$) and which we cannot derive from an Euler factor. For $\ell=5,7$ no term is needed at $q=\ell$ or at $q=2,3$ beyond Definition~\ref{def:terms}; the term at $2$ given there can be positive (592c1, Remark~\ref{rem:wild}). That no term is needed at $q=\ell\ge5$ is true but not sharp, by the following observation.

\begin{remark}[a twist relation at $q=\ell$; data only]\label{rem:oq1}
For each optimal curve $E$ without complex multiplication of conductor below $10{,}000$ with $\ell\mid N$, $\ell\in\{3,5,7,11,13\}$, compare $v_\ell(\degphi)$ with $v_\ell$ of the degree of the optimal curve of the isogeny class of the twist $E^{(\ell)}$ of $E$ by $\chi_{\ell^*}$, counting the curves whose twist has conductor below $10{,}000$. (a) If $E$ has type $\mathrm{II}^*$, $\mathrm{III}^*$ or $\mathrm{IV}^*$ at $\ell$, the difference is exactly $1$ in all $4{,}412$ cases ($2{,}141$, $1{,}210$, $659$, $239$, $163$ at $\ell=3,5,7,11,13$; $4{,}251$ of them with $E[\ell]$ irreducible). (b) If $E$ has type $\mathrm{I}_n^*$ at $\ell$, the difference is $0$ in all cases ($3{,}639$, $1{,}360$, $677$, $240$, $173$). (c) If $E$ has type $\mathrm{I}_0^*$ at $\ell$ and $E[\ell]$ is irreducible, the difference is $v_\ell\bigl((\ell+1)^2-a_\ell(E^{(\ell)})^2\bigr)$ in all $1{,}915$ cases; it fails in $13$ cases with $E[\ell]$ reducible, for instance 126a1 at $\ell=3$ and 275b1 at $\ell=5$. At $\ell=3$, (a) is the pattern of the weights $w_3$, each starred type weighing one more than its partner under the exchanges of Lemma~\ref{lem:twist}(ii). At $\ell\ge5$ it suggests a weight $1$ for the types $\mathrm{II}^*$, $\mathrm{III}^*$, $\mathrm{IV}^*$ at $q=\ell$; added to \eqref{eq:rule}, it gives no violation below conductor $10{,}000$ at $\ell=5,7,11,13$ and is attained by most of these curves. In (b) and (c) the difference is the value at $q=\ell$ of the numerators $v_\ell(q^2-1)$ and $v_\ell(q-1)+v_\ell((q+1)^2-a'^2)$ of the table of Section~\ref{sec:localterms}. These are observations on the whole range at once, with no hold-out, and we offer no explanation; they bear on the first of the open questions of Section~\ref{sec:open}.
\end{remark}

\begin{table}[ht]
\centering\small
{\small
\begin{tabular}{lccc}
\hline
type & weight 1 & weight 2 & weight 3 \\
\hline
IV (any $q$) & 1 (25,414) & 2 (969) & 3 (11) \\
IV$^*$ (any $q$) & 1 (30,168) & 2 (2,268) & 3 (45) \\
III (any $q$) & 0 (34,904) & 0 (2,208) & 0 (34) \\
III$^*$ (any $q$) & 0 (25,574) & 0 (1,307) & 0 (29) \\
$\mathrm{I}_0^*$ (any $q$) & 0 (42,924) & 0 (3,093) & 0 (53) \\
II, II$^*$ at odd $q$ & 1 (56,509) & 2 (2,502) & 3 (12) \\
$\mathrm{I}_n^*$ at odd $q$, weight $v_3(n)$ & 1 (28,175) & 2 (3,453) & 3 (89) \\
II, II$^*$ at $q=2$ & 0 (44,672) & -- & -- \\
$\mathrm{I}_n^*$ at $q=2$, weight $v_3(n)$ & 0 (21,176) & 1 (5,329) & 1 (342) \\
\hline
\end{tabular}
}
\caption{Minimum of $v_3(\degphi)$ (number of curves) over working-set curves with $E[3]$ irreducible whose only possible contributions at $3$ come from the primes named in the row, with the given number of primes (untwisted rows) or total weight of \eqref{eq:twist} (twisted rows).}\label{tab:isolated}
\end{table}

\subsection{Out of table}\label{sec:oot}
To test the statements beyond the tables, we generated curves with prescribed local types and conductor between $500000$ and $4\cdot10^6$ (families with a type $\mathrm{IV}$ or $\mathrm{IV}^*$ prime, with two or three type $\mathrm{II}$ primes, with mixed types, and random curves), kept those with trivial isogeny class (so that the curve is optimal and $E[\ell]$ is irreducible for every $\ell$) and not CM, and computed their modular degrees with \texttt{ellmoddegree}. Table~\ref{tab:oot} summarises them; none has image $\Ns$, and Conjecture~\ref{conj:euler} and its two corollary forms hold for all of them (Tables~\ref{tab:uniform} and~\ref{tab:counts}, last rows), with equality in \eqref{eq:rule} for $52.5\%$ of the $828$ curves with a positive bound at $\ell=3$.

\begin{table}[ht]
\centering\small
{\small
\begin{tabular}{lrrr}
\hline
family & curves & largest conductor & median seconds per degree \\
\hline
II,II & 98 & 3,968,800 & 0.4 \\
II,II,II & 1 & 2,832,200 & 0.2 \\
IV & 334 & 3,988,012 & 1.1 \\
IV* & 104 & 3,994,600 & 4.9 \\
IVmixed & 266 & 3,983,504 & 1.3 \\
random & 90 & 2,710,716 & 32.1 \\
all & 893 & 3,994,600 & 1.5 \\
\hline
\end{tabular}
}
\caption{Curves outside the tables, with modular degrees computed independently.}\label{tab:oot}
\end{table}

\subsection{The Eisenstein case}\label{sec:eis}
Among the $75{,}430$ working-set optimal non-CM curves with a rational $3$-isogeny, $67{,}845$ have a positive bound in \eqref{eq:main} and $2{,}777$ a positive deficit $D$; $D=2$ occurs $22$ times. No curve without a rational $3$-torsion point has a deficit, whatever its $3$-isogeny structure. With a rational point of order $3$, mod $3$ image of Borel type and no point of order $9$, $D\le1$; $D=2$ occurs only when $E[3]\simeq\Z/3\oplus\mu_3$ (the image labelled \texttt{3Cs.1.1} in the tables, in the notation of \cite{Sutherland16}: the subgroup of order $2$ of the split Cartan subgroup that fixes a point; seventeen curves, all of rank $0$, with torsion $\Z/3$ or $\Z/6$, for instance 14a1, with torsion $\Z/6$, types $\mathrm{I}_6$ at $2$ and $\mathrm{I}_3$ at $7$, and $\degphi=1$) or when $E$ has a point of order $9$ (five curves, with image \texttt{3B.1.1}, a Borel image fixing a point of order $3$, for instance 714i1, with $\degphi=1080=2^3\cdot3^3\cdot5$, and 6942n1, with $\degphi=65664=2^7\cdot3^3\cdot19$). The two sources of a deficit of $2$ never combine: none of the $2{,}065$ curves of conductor at most $400000$ with mod $3$ image \texttt{3Cs.1.1} has a point of order $9$ (their torsion is $\Z/3$ or $\Z/6$), and the $18$ optimal curves with a point of order $9$ all have image \texttt{3B.1.1}. This is a theorem.

\begin{proposition}\label{prop:no9}
Let $E/\Q$ be without complex multiplication, with $E[3]\simeq\Z/3\oplus\mu_3$. Then $E(\Q)$ has no point of order $9$. Consequently, for every non-CM $E/\Q$, $v_3(\#E(\Q)_{\mathrm{tors}})+[E[3]\simeq\Z/3\oplus\mu_3]\le2$, and for $\ell\ge5$, $v_\ell(\#E(\Q)_{\mathrm{tors}})\le1$.
\end{proposition}
\begin{proof}
Suppose $P\in E(\Q)$ has order $9$. Then $3P$ is a rational point of order $3$, so $\langle3P\rangle$ is the line $\Z/3$ of $E[3]$ ($\mu_3$ has no nonzero rational point), and $\langle P\rangle\cap E[3]=\langle3P\rangle$ meets $\mu_3$ trivially. Let $\phi\colon E\to E'=E/\mu_3$, rational, and $\hat\phi\colon E'\to E$ its dual. From $\phi\circ\hat\phi=[3]$, $\hat\phi(E'[3])\subseteq\ker\phi=\mu_3$, with equality as $\#\hat\phi(E'[3])=9/\#\ker\hat\phi=3$. Let $C=\hat\phi^{-1}(\langle P\rangle)\subset E'$, a $G_\Q$-stable subgroup of order $27$. If $C$ contained $E'[3]$, then $\mu_3=\hat\phi(E'[3])\subseteq\langle P\rangle$, a contradiction; so $C$ is cyclic, and $E'\to E'/C$ is a rational cyclic isogeny of degree $27$. A rational cyclic $27$-isogeny exists only for curves with complex multiplication (by the order of discriminant $-27$, $j=-2^{15}\cdot3\cdot5^3$) \cite{Ligozat75}; the level $27$ is in the list of the $N$ with $Y_0(N)(\Q)\ne\emptyset$ of \cite[Theorem 1]{Kenku}; so $E'$, hence $E$, has complex multiplication, a contradiction. The consequences follow from Mazur's list of torsion groups \cite{Mazur77}: $v_3(\#E(\Q)_{\mathrm{tors}})\le2$, with equality only for $\Z/9$, and $v_\ell(\#E(\Q)_{\mathrm{tors}})\le1$ for $\ell\ge5$.
\end{proof}

So the bound $D\le2$ of Conjecture~\ref{conj:eis} at $\ell=3$ follows from its first bound, and is not an additional conjecture. The bounds of Conjecture~\ref{conj:eis} hold with no exception on all three ranges, and the torsion bound is attained with $D>0$ in $2{,}773$ cases. For $\ell=5$ and $7$ the deficits are at most $1$ and occur only with a rational $\ell$-torsion point ($192$ and $18$ cases in the working set). Measured against the stronger bound \eqref{eq:twist} the deficit is not controlled by the torsion: it exceeds the bound of Conjecture~\ref{conj:eis} for $1{,}148$ working-set curves at $\ell=3$ ($189$ and $39$ on the hold-out and the extension) and for $141$ at $\ell=5$. In the Eisenstein case the additive terms do not contribute reliably, and Conjecture~\ref{conj:eis} is stated for \eqref{eq:main} only.

\section{The prime 2}\label{sec:two}
At $\ell=2$ the modular degree is the subject of Watkins's conjecture $2^{\operatorname{rank}E(\Q)}\mid\degphi$ \cite{Watkins02} and of a substantial literature \cite{Dummigan06,CalegariEmerton09,Yazdani11,CaroPasten22,HatleyKundu,BKP}, and the mod $2$ representation has a rational $2$-isogeny for a large fraction of curves. We record only the following. For working-set optimal curves without a rational $2$-isogeny and not CM, the inequality $v_2(\degphi)\ge\sum_{q\,\|\,N}v_2(n_q)+\#\{q:\ \text{type }\mathrm{III},\mathrm{III}^*,\mathrm{I}_0^*\text{ or }\mathrm{I}_n^*\}$ holds with no exception on all three ranges ($964{,}432$, $325{,}510$ and $313{,}741$ curves), while counting $\mathrm{I}_0^*$ with weight $2=v_2(\#\Phig_q)$ fails ($123$, $26$ and $20$ exceptions). The isolated-configuration minima suggest that types $\mathrm{III}$, $\mathrm{III}^*$ and $\mathrm{I}_n^*$ contribute $3$ each and types $\mathrm{II}$, $\mathrm{IV}$ and $\mathrm{I}_0^*$ contribute $1$ each, in line with the theorem of Calegari and Emerton \cite{CalegariEmerton09} (we use its statement and numbering in the arXiv version, where it is Theorem 1.1): a curve of odd modular degree has at most two odd bad primes and even analytic rank, and either it has a rational $2$-torsion point, or its conductor is prime, it is supersingular at $2$ and $E(\R)$ is connected, or it has complex multiplication and $N\in\{27,32,49,243\}$ (for instance 27a1 and 243b1, of degrees $1$ and $9$, with torsion $\Z/3$). We have not attempted to reconcile the exponents with the results of \cite{Dummigan06,CaroPasten22,HatleyKundu,BKP}, nor with the Euler factors at $\ell=2$, which lie outside the range of \cite{DFG04}.

\section{Open questions}\label{sec:open}
\begin{enumerate}
\item Extend Theorem~\ref{thm:main} to the primes $\ell$ dividing $N$, where the data show the inequality with the weights $w_3$ at $q=\ell=3$: identify these weights with the local term at $\ell$ of the congruence-ideal formula, where the Euler factor is replaced by a condition at $\ell$. The twist relation of Remark~\ref{rem:oq1}, observed at $\ell=3,5,7,11,13$, is a lead: it reproduces the pattern of the $w_3$ and suggests a weight at $q=\ell\ge5$ as well.
\item Prove the corrected identity of Remark~\ref{rem:ns} for the curves with image $N_s(\F_3)$, that is, establish $H^1_f(\Q,\ad V_3E)=0$ and $v_\lambda(\eta^\Sigma_{f,\lambda})=\len H^1_\Sigma(\Q,W)$ without the hypothesis on $\Q(\sqrt{-3})$ in this dihedral case; the data say that the conclusion holds with the correction $h_3(E)=1$, and the equality of Theorem~\ref{thm:quant} holds in the $8$ pairs with image $\Ns$ below conductor $10{,}000$, which lie outside its equality clause.
\item The Tamagawa factors of $\ad T$ at the additive primes are known (Lemmas~\ref{lem:tamfactor} and~\ref{lem:tamadd}), so that by Corollary~\ref{cor:explicit} the excess of $v_\ell(\degphi)$ over the explicit local bound is exactly the length of the adjoint Selmer group (Table~\ref{tab:excess}). Compute that Selmer group independently for the curves exceeding the bound, such as 1184f1 and 1748e1 at $\ell=3$, so as to account for the excess curve by curve.
\item Identify the deficit of Conjecture~\ref{conj:eis} with a natural invariant of the Eisenstein ideal at $\ell$, presumably through the cuspidal subgroup of $J_0(N)$ \cite{Mazur77}.
\item Reconcile Section~\ref{sec:two} with the literature on Watkins's conjecture and with the Euler factors at $\ell=2$.
\item A smaller point: whether, in Proposition~\ref{prop:ns-structure}, the type at $2$ is always $\mathrm{III}$ or $\mathrm{III}^*$ when $R=\{2\}$. It is so for all $16$ optimal curves without and all $6$ with complex multiplication below conductor $10{,}000$ where this happens.
\end{enumerate}

\appendix
\section{Reproducibility}\label{app:repro}
All scripts, the pinned data commit and the result tables are in the repository \url{https://github.com/dsr001001/mcptrial} (directory \texttt{ecmine}). The symmetric-square Euler factors are computed by \texttt{scripts/step10\_sym2\_euler.py}, the rule is checked on the whole ranges by \texttt{scripts/step10\_rule.py} and \texttt{scripts/step11\_tw\_boundary.py} (correction term, image statistics), the Tamagawa factors by \texttt{scripts/step12\_tamagawa.py}, the two corollary forms by \texttt{scripts/conj\_check.py}; each runs on the full tables in minutes to hours. A single curve can be checked against \eqref{eq:rule} in PARI/GP (version 2.15 or later) with the following functions, which return $[v_\ell(\degphi),\text{bound}]$; they are the file \texttt{scripts/euler\_rule.gp}. The curve must be the optimal curve of its class and $E[\ell]$ must be irreducible: for instance 990h1 returns \texttt{[4, 3]} while the optimal curve 990h3 of its class returns \texttt{[3, 3]}, and 11a1 at the Eisenstein prime $\ell=5$ returns \texttt{[0, 1]}.
{\footnotesize\begin{verbatim}
rule(E, l, {iii = 0}) =
{ my(R = ellglobalred(E), P = R[4][,1], K = R[5], S = lfunsympow(E, 2), b = 0, j = E.j);
  for(i = 1, #P,
    my(q = P[i], k = K[i][2], e = 0, t = 0, vj = valuation(j, q));
    if(q != l,
      my(F = 1/lfuneuler(S, q), r = subst(F, x, 1/q^2)); \\ P_q(q^{-2}): omitted factor
      if(k >= 5, r /= (1 - 1/q^2));                     \\ multiplicative: naive one out
      if(abs(k) != 3 || iii, e = valuation(r, l)),      \\ III, III* (codes 3,-3) if iii
      if(l == 3, e = if(k == 4 || k == 2, 1,
                     if(k == -4 || k == -2, 2, if(k == -3, 1, 0)))));   \\ w_3 at q = 3
    if(vj < 0, t = valuation(-vj, l));                   \\ Tamagawa exponent
    b += e + t);
  if(iii && l == 3, b -= ns3(E));                        \\ the correction term h_3
  [valuation(ellmoddegree(E), l), b]; }
ns3(E) = my(F = factor(elldivpol(E, 3))); (#F[,1] == 2 && poldegree(F[1,1]) == 2);
\\ psi_3 has degree 4 and is squarefree, so two factors, one quadratic, means two quadratics
\end{verbatim}}
Here \texttt{K[i][2]} is PARI's Kodaira code: $n+4$ for $\mathrm{I}_n$ and $-(n+4)$ for $\mathrm{I}_n^*$ when $n\ge1$, $1$ for $\mathrm{I}_0$ and $-1$ for $\mathrm{I}_0^*$, $\pm2$ for $\mathrm{II}/\mathrm{II}^*$, $\pm3$ for $\mathrm{III}/\mathrm{III}^*$, $\pm4$ for $\mathrm{IV}/\mathrm{IV}^*$. The function \texttt{rule(E,l)} leaves the type $\mathrm{III}$ terms out and \texttt{rule(E,l,1)} is \eqref{eq:rule}; at $q=\ell$ the term $t_\ell$ is that of Definition~\ref{def:terms}, which is defined there. For example \texttt{rule(ellinit([1,1,1,-30,-76]), 3)} (the curve 121a1) returns \texttt{[1, 1]}; for 726a1, a curve with image $\Ns$ and type $\mathrm{III}$ at $11$, \texttt{rule(E,3)} and \texttt{rule(E,3,1)} both return \texttt{[2, 2]}, the type $\mathrm{III}$ term and the correction cancelling. The function \texttt{ellmoddegree} needs about $32$ MB of stack with PARI 2.17 for conductors in the thousands (with PARI 2.15 and larger conductors, \texttt{parisizemax} of a few hundred megabytes). The pipeline was rerun from a clean clone and reproduced every count of Table~\ref{tab:counts}. Lemma~\ref{lem:twist}(ii) was also checked numerically, with the script \texttt{twist\_lemma\_check.gp} of the repository, on $4{,}577$ random pairs (curve, odd bad prime), with no exception, and its exchanges at $q=3$ on all $5{,}122$ pairs (curve, twist by $\Q(\sqrt{-3})$) below conductor $10{,}000$, and Lemma~\ref{lem:euler} against \texttt{lfunsympow} on every curve of the tables.

\section*{Statement on the use of an AI system}
This work was carried out with an AI system (Claude, Anthropic) at every stage, and it would not have been done without it. The author set the goal, a search for new relations in Cremona's tables, chose the data and the protocol (a frozen snapshot of the tables, a working set, a range held back for one-time tests, and a list of tests fixed before any was run), and decided at each step what to pursue and what to keep. Under that direction, and in discussion with the author, the AI system did the rest: it built the feature table and the list of known relations to be excluded; ran the pre-registered search, which produced the candidate relation between the modular degree and the local invariants; searched the literature and identified the results of Ribet-Takahashi, Pollack-Weston, Kim-Ota, Agashe-Ribet-Stein and Diamond-Flach-Guo behind the statements; replaced the Tamagawa number by the geometric component group and formulated the conjectures and their successive refinements, up to the Euler-factor rule and its correction term; wrote and ran the verification code in Python and PARI/GP on the four ranges, generated the out-of-table curves and computed their modular degrees and Tate parameters; read the published text of \cite{DFG04} and assembled the proof of Theorem~\ref{thm:main} from its statements, computed the Tamagawa factor of Lemma~\ref{lem:tamfactor}, and proposed and wrote the proofs of Theorems~\ref{thm:qual} and~\ref{thm:quant} and of Proposition~\ref{prop:ns}; drafted the text and the tables; and prepared the responses to reviews of drafts, two of them written by AI systems, from which the identification of the local factors and of the boundary of the hypothesis arose. Several of the results and proofs (among them Propositions~\ref{prop:pullback}, \ref{prop:Sf}, \ref{prop:shapes}, \ref{prop:ns-structure} and~\ref{prop:no9}, Lemmas~\ref{lem:tamadd} and~\ref{lem:step1}, Corollary~\ref{cor:explicit}, the equality in Theorem~\ref{thm:quant} and the exchanges at $q=3$ in Lemma~\ref{lem:twist}), and the checks of the counts below conductor $10{,}000$, were worked out with the same AI system while writing a companion textbook on this note. The note was prepared with an AI system under the author's direction and has not yet been refereed; the author has to review it before it is circulated. The scripts, the result files and the pinned data commit are in the repository named in the appendix, so that every count can be rerun; the proofs cite the published statements they use, with their numbering. The author takes full responsibility for all statements, proofs, computations and errors.

\bibliographystyle{amsplain}
\bibliography{refs}
\end{document}
